\documentclass[11pt,a4paper,oneside]{book}

% =========================================================
% EUFISICA books — Física e Química A — 10.º ano
% Autor: José Gonçalves
% Edição revista, alinhada com as Aprendizagens Essenciais de março de 2026
% =========================================================

% ---------- Codificação, língua e tipografia ----------
\usepackage[T1]{fontenc}
\usepackage[utf8]{inputenc}
\usepackage[portuguese]{babel}
\usepackage{lmodern}
\usepackage{microtype}

% ---------- Página e cabeçalhos ----------
\usepackage[a4paper,margin=2.2cm,headheight=15pt]{geometry}
\usepackage{fancyhdr}
\pagestyle{fancy}
\fancyhf{}
\fancyhead[L]{\small\bfseries\color{eublue}EUFISICA books}
\fancyhead[R]{\small\color{eudark}\nouppercase{\leftmark}}
\fancyfoot[L]{\scriptsize\color{gray}José Gonçalves}
\fancyfoot[R]{\bfseries\color{eublue}\thepage}
\renewcommand{\headrulewidth}{0.6pt}
\renewcommand{\footrulewidth}{0.3pt}
\renewcommand{\chaptermark}[1]{\markboth{\thechapter.\ #1}{}}

% ---------- Matemática, unidades, química ----------
\usepackage{amsmath,amssymb,mathtools}
\usepackage{siunitx}
\sisetup{locale=FR,per-mode=symbol,detect-all=true}
\usepackage[version=4]{mhchem}

% ---------- Tabelas e listas ----------
\usepackage{booktabs,tabularx,longtable,array,multirow}
\usepackage{float}
\usepackage{caption}
\usepackage{colortbl}
\usepackage{enumitem}
\usepackage{hanging}
\setlist[itemize]{topsep=4pt,itemsep=2pt}
\setlist[enumerate]{topsep=4pt,itemsep=3pt}

% ---------- Cor e caixas ----------
\usepackage{xcolor}
\definecolor{eublue}{HTML}{0D47A1}
\definecolor{eucyan}{HTML}{0097A7}
\definecolor{eugreen}{HTML}{2E7D32}
\definecolor{euorange}{HTML}{EF6C00}
\definecolor{eured}{HTML}{C62828}
\definecolor{eugrey}{HTML}{ECEFF1}
\definecolor{eudark}{HTML}{263238}
\definecolor{euyellow}{HTML}{FFF8E1}
\definecolor{eulightblue}{HTML}{EAF2FB}
\definecolor{eulightcyan}{HTML}{E8F7F8}
\definecolor{eulightgreen}{HTML}{EDF7EE}
\definecolor{eulightorange}{HTML}{FFF1E6}
\definecolor{eupurple}{HTML}{6A1B9A}
\usepackage[most]{tcolorbox}
\tcbset{
  enhanced,
  boxrule=0.8pt,
  arc=2.2mm,
  outer arc=2.2mm,
  left=3mm,right=3mm,top=2.2mm,bottom=2.2mm,
  before skip=9pt,after skip=9pt,
  fonttitle=\bfseries,
  attach boxed title to top left={xshift=4mm,yshift*=-2mm},
  boxed title style={arc=1.2mm,boxrule=0pt,left=2mm,right=2mm,top=1mm,bottom=1mm}
}

\newtcolorbox{objetivos}{
  colback=eulightblue,colframe=eublue,colbacktitle=eublue,coltitle=white,
  title=Objetivos de aprendizagem,breakable,borderline west={2.2mm}{0pt}{eublue}}
\newtcolorbox{mundo}{
  colback=euyellow,colframe=euorange,colbacktitle=euorange,coltitle=white,
  title=Ciência no mundo real,breakable,borderline west={2.2mm}{0pt}{euorange}}
\newtcolorbox{stem}{
  colback=eulightgreen,colframe=eugreen,colbacktitle=eugreen,coltitle=white,
  title=Desafio STEM,breakable,borderline west={2.2mm}{0pt}{eugreen}}
\newtcolorbox{nota}{
  colback=eulightcyan,colframe=eucyan,colbacktitle=eucyan!85!black,coltitle=white,
  title=Nota importante,breakable,borderline west={2.2mm}{0pt}{eucyan}}
\newtcolorbox{seguranca}{
  colback=red!3,colframe=eured,colbacktitle=eured,coltitle=white,
  title=Segurança,breakable,borderline west={2.2mm}{0pt}{eured}}
\newtcolorbox{exemplo}[1][]{
  colback=white,colframe=eublue,colbacktitle=eublue,coltitle=white,
  title=Exemplo resolvido,#1,breakable,borderline west={1.4mm}{0pt}{eublue}}
\newtcolorbox{atividade}[1][]{
  colback=white,colframe=eucyan,colbacktitle=eucyan!85!black,coltitle=white,
  title=Atividade experimental / investigação,#1,breakable,borderline west={1.4mm}{0pt}{eucyan}}
\newtcolorbox{exercicios}[1][]{
  colback=eugrey!55,colframe=eudark,colbacktitle=eudark,coltitle=white,
  title=Exercícios,#1,breakable,borderline west={1.4mm}{0pt}{eudark}}
\newtcolorbox{solucao}[1][]{
  colback=eulightgreen,colframe=eugreen,colbacktitle=eugreen,coltitle=white,
  title=Solução,#1,breakable,borderline west={1.4mm}{0pt}{eugreen}}
\newtcolorbox{curriculo}[1][]{
  colback=eulightorange,colframe=euorange,colbacktitle=euorange,coltitle=white,
  title=Ligação às Aprendizagens Essenciais,#1,breakable,borderline west={1.4mm}{0pt}{euorange}}
\newtcolorbox{metodo}[1][]{
  colback=eulightcyan,colframe=eucyan,colbacktitle=eucyan!85!black,coltitle=white,
  title=Estratégia de resolução,#1,breakable,borderline west={1.4mm}{0pt}{eucyan}}
\newtcolorbox{sabias}[1][]{
  colback=purple!3,colframe=eupurple,colbacktitle=eupurple,coltitle=white,
  title=Sabias que...?,#1,breakable,borderline west={1.4mm}{0pt}{eupurple}}
\newtcolorbox{conceito}[1][]{
  colback=white,colframe=eublue!75!black,colbacktitle=eublue!75!black,coltitle=white,
  title=Ideia-chave,#1,breakable,borderline west={1.4mm}{0pt}{eublue!75!black}}
\newtcolorbox{leituragrafico}[1][]{
  colback=white,colframe=euorange,colbacktitle=euorange,coltitle=white,
  title=Ler o gráfico,#1,breakable,borderline west={1.4mm}{0pt}{euorange}}

% ---------- TikZ e circuitos ----------
\usepackage{tikz}
\usetikzlibrary{arrows.meta,positioning,calc,patterns,shapes.geometric,decorations.pathmorphing,decorations.pathreplacing,angles,quotes,intersections}
\usepackage[europeanresistors]{circuitikz}

% ---------- Gráficos ----------
\usepackage{pgfplots}
\usepackage{pgf-spectra}
\pgfplotsset{compat=1.18}
\pgfplotsset{
  every axis/.append style={
    tick label style={font=\small},
    label style={font=\small},
    title style={font=\bfseries\small},
    grid style={draw=gray!22},
    axis line style={draw=eudark!70},
    legend style={font=\small,draw=none,fill=white,fill opacity=.85,text opacity=1}
  }
}
\captionsetup{font=small,labelfont={bf,color=eublue},labelsep=period,justification=centering}
\captionsetup[table]{position=above,skip=6pt}
\captionsetup[figure]{position=below,skip=5pt}

% ---------- Hiperligações ----------
\usepackage{xurl}
\usepackage{hyperref}
\hypersetup{
  colorlinks=true,
  linkcolor=eublue,
  urlcolor=eucyan,
  citecolor=eugreen,
  pdftitle={Física e Química A — 10.º ano},
  pdfauthor={José Gonçalves}
}

% ---------- Títulos ----------
\usepackage{titlesec}
\titleformat{\chapter}[display]
  {\normalfont\bfseries\color{eudark}}
  {\filleft\colorbox{eublue}{\parbox{2.7cm}{\centering\color{white}\Large\chaptername\\[-1mm]\fontsize{30}{32}\selectfont\thechapter}}}
  {1.2ex}
  {\titlerule[1.1pt]\vspace{1ex}\fontsize{25}{29}\selectfont\raggedright\color{eublue}}
\titleformat{\section}
  {\Large\bfseries\color{eublue}}{\thesection}{0.65em}{}[\vspace{-2pt}\color{eucyan!65}\titlerule]
\titleformat{\subsection}
  {\large\bfseries\color{eudark}}{\thesubsection}{0.65em}{}
\titleformat{\part}[display]
  {\centering\normalfont\bfseries\color{white}}
  {\colorbox{euorange}{\parbox{0.94\textwidth}{\centering\Large PARTE \thepart}}}
  {1.2ex}
  {\color{eublue}\Huge}

% ---------- Comandos úteis ----------
\newcommand{\EUFISICA}{\textsc{EUFISICA}}
\newcommand{\vect}[1]{\boldsymbol{#1}}
\newcommand{\dd}{\,\mathrm{d}}
\newcommand{\Ecin}{E_{\mathrm c}}
\newcommand{\Epg}{E_{\mathrm{p,g}}}
\newcommand{\Eel}{E_{\mathrm{el}}}
\newcommand{\DeltaH}{\Delta H}
\newcommand{\NA}{N_{\mathrm A}}

% =========================================================
% DOCUMENTO
% =========================================================
\begin{document}

% =========================================================
% CAPA
% =========================================================
\begin{titlepage}
\thispagestyle{empty}
\begin{tikzpicture}[remember picture,overlay]
  \fill[eudark] (current page.south west) rectangle (current page.north east);
  \fill[eublue] (current page.south west) rectangle ([xshift=4.4cm]current page.north west);
  \fill[eucyan,opacity=.30] ([xshift=4.4cm]current page.south west) rectangle ([xshift=4.9cm]current page.north west);

  % átomo estilizado
  \coordinate (A) at ([xshift=-3.2cm,yshift=2.0cm]current page.center);
  \fill[euorange] (A) circle (4pt);
  \draw[white,opacity=.85,rotate around={25:(A)}] (A) ellipse (2.0cm and .72cm);
  \draw[white,opacity=.85,rotate around={85:(A)}] (A) ellipse (2.0cm and .72cm);
  \draw[white,opacity=.85,rotate around={145:(A)}] (A) ellipse (2.0cm and .72cm);
  \fill[eucyan] ($(A)+(1.72,0.35)$) circle (3pt);
  \fill[eugreen] ($(A)+(-.8,1.22)$) circle (3pt);
  \fill[euorange] ($(A)+(-.7,-1.20)$) circle (3pt);

  % trajetória e energia
  \draw[euorange,very thick,-{Stealth[length=4mm]}]
    ([xshift=1.2cm,yshift=-4.3cm]current page.center)
    .. controls +(2.1,1.1) and +(-1.2,-.6) ..
    ([xshift=6.3cm,yshift=-1.2cm]current page.center);
  \foreach \x/\y in {1.7/-4.0,2.8/-3.6,4.0/-2.9,5.1/-2.1}
    \fill[white] ([xshift=\x cm,yshift=\y cm]current page.center) circle (1.8pt);

  % onda eletromagnética
  \begin{scope}[shift={(current page.center)}]
    \draw[eucyan,line width=1.25pt,samples=120,smooth,domain=0:7]
      plot ({0.8+\x},{4.5+0.42*sin(220*\x)});
  \end{scope}
  % molécula estilizada
  \coordinate (M) at ([xshift=5.2cm,yshift=4.3cm]current page.center);
  \draw[white!70,line width=1.2pt] ($(M)+(-.85,0)$)--(M)--($(M)+(.85,.45)$);
  \fill[eured] (M) circle (7pt);
  \fill[white] ($(M)+(-.85,0)$) circle (5pt);
  \fill[white] ($(M)+(.85,.45)$) circle (5pt);
  % mini circuito gráfico
  \draw[white!55,line width=1pt] ([xshift=2cm,yshift=-5.6cm]current page.center) rectangle ++(5.2,1.1);
  \draw[euorange,line width=1.5pt] ([xshift=2.25cm,yshift=-5.05cm]current page.center)--++(.7,0)--++(.25,.35)--++(.45,-.7)--++(.45,.7)--++(.45,-.7)--++(.25,.35)--++(.75,0);
\end{tikzpicture}

\vspace*{2.1cm}
\hspace*{5.1cm}
\begin{minipage}{0.62\textwidth}
{\color{white}\Large\textbf{\EUFISICA\ books}}\\[1.2cm]
{\color{white}\fontsize{31}{35}\selectfont\bfseries Física e Química A}\\[0.35cm]
{\color{eucyan}\fontsize{25}{29}\selectfont\bfseries 10.º ano}\\[0.9cm]
{\color{white}\Large Conceitos, gráficos, laboratório e desafios STEM}\\[1.6cm]
{\color{white}\large Autor: \textbf{José Gonçalves}}\\[0.25cm]
{\color{white}\normalsize Ensino Secundário Português}
\end{minipage}

\vfill
\hspace*{5.1cm}{\color{white!75}\small Edição EUFISICA books}
\end{titlepage}

\frontmatter

\chapter*{Nota ao leitor}
Este livro foi concebido para apoiar a aprendizagem de \textbf{Física e Química A no 10.º ano do ensino secundário português}. Esta edição integra também esquemas e problemas selecionados das fichas de trabalho de Física do autor e foi revista por confronto com as \textbf{Aprendizagens Essenciais publicadas pela Direção-Geral da Educação em março de 2026}. A organização segue os três grandes temas do referencial atual: \textit{Elementos Químicos e sua organização}, \textit{Propriedades e Transformações da Matéria} e \textit{Energia e sua conservação}. O texto valoriza a experimentação, a resolução de problemas, o tratamento de dados, a comunicação científica, a ligação ao quotidiano e abordagens STEM/STEAM — Ciência, Tecnologia, Engenharia, Artes e Matemática.

O objetivo não é apenas memorizar fórmulas. Ao longo do livro, cada conceito é associado a perguntas do mundo real: Como sabemos de que são feitas as estrelas? Porque é que alguns materiais conduzem eletricidade? Como funciona uma bateria? Porque aquece um travão? Como se mede a eficiência de um painel fotovoltaico? Como se pode reduzir o consumo energético de uma casa?

\begin{nota}
As atividades laboratoriais devem ser realizadas com supervisão adequada, respeitando as regras de segurança do laboratório e as fichas de dados de segurança das substâncias utilizadas.
\end{nota}


\chapter*{Como usar este livro}
\addcontentsline{toc}{chapter}{Como usar este livro}
Cada capítulo foi construído em quatro camadas complementares:
\begin{enumerate}
  \item \textbf{conceitos e modelos}, com explicações progressivas e relações matemáticas;
  \item \textbf{exemplos resolvidos}, onde se explicita o raciocínio e não apenas o resultado;
  \item \textbf{trabalho laboratorial e STEM}, com planeamento, recolha e tratamento de dados;
  \item \textbf{exercícios de consolidação e integração}, cujas soluções se encontram no final.
\end{enumerate}

Ao resolver um problema, começa por identificar o sistema, escreve os dados com unidades, escolhe o modelo físico ou químico adequado, efetua os cálculos e termina verificando se o resultado é plausível. Em questões experimentais, distingue sempre \textbf{o que foi medido diretamente} do que foi \textbf{calculado a partir das medições}.

\begin{figure}[H]
\centering
\begin{tikzpicture}[node distance=1.2cm and 1.1cm,>=Stealth]
\tikzset{etapa/.style={rounded corners=2mm,minimum width=3.3cm,minimum height=1.05cm,align=center,font=\small\bfseries,draw=eublue,fill=eulightblue}}
\node[etapa] (q) {1. Perguntar\\e prever};
\node[etapa,right=of q] (m) {2. Modelar\\e planear};
\node[etapa,right=of m] (m2) {3. Medir\\ou calcular};
\node[etapa,below=of m2] (a) {4. Analisar\\e representar};
\node[etapa,left=of a] (c) {5. Concluir\\com evidência};
\node[etapa,left=of c] (r) {6. Rever\\e comunicar};
\draw[->,thick,eucyan] (q)--(m); \draw[->,thick,eucyan] (m)--(m2);
\draw[->,thick,eucyan] (m2)--(a); \draw[->,thick,eucyan] (a)--(c);
\draw[->,thick,eucyan] (c)--(r); \draw[->,thick,eucyan] (r)--(q);
\end{tikzpicture}
\caption{Ciclo de trabalho científico usado ao longo do livro.}
\end{figure}

\chapter*{Mapa curricular e verificação de cobertura}
\addcontentsline{toc}{chapter}{Mapa curricular e verificação de cobertura}

\begin{curriculo}
A edição foi revista por confronto com as \textbf{Aprendizagens Essenciais de Física e Química A do 10.º ano — março de 2026}, atualmente publicadas pela Direção-Geral da Educação. O referencial oficial organiza o 10.º ano em três grandes temas: \textit{Elementos Químicos e sua organização}, \textit{Propriedades e Transformações da Matéria} e \textit{Energia e sua conservação}. O quadro seguinte mostra onde são trabalhados os conhecimentos e capacidades nucleares, incluindo a vertente experimental.
\end{curriculo}

\small
\begin{longtable}{p{.19\textwidth}p{.28\textwidth}p{.43\textwidth}}
\toprule
Domínio & Núcleo curricular & Cobertura no livro \\
\midrule
\endfirsthead
\toprule
Domínio & Núcleo curricular & Cobertura no livro \\
\midrule
\endhead
Elementos químicos & Massa e tamanho dos átomos & Escalas e ordens de grandeza; estrutura atómica; isótopos; unidade de massa atómica; massa atómica relativa; mole; constante de Avogadro; massa molar; medições de massa e volume; AL 1.1 -- gota de água. \\
Elementos químicos & Energia dos eletrões & Espetro eletromagnético; fotões; emissão e absorção; hidrogénio; quantização; átomos polieletrónicos; energia de remoção; orbitais $s$, $p$ e $d$; configurações eletrónicas até $Z=23$; AL 1.2 -- teste de chama. \\
Elementos químicos & Tabela Periódica & Evolução histórica; grupos, períodos e blocos; configurações eletrónicas; raio atómico; energia de ionização; formação de iões; gases nobres; AL 1.3 -- densidade relativa de metais por picnometria. \\
Propriedades da matéria & Ligação química & Energia e distância internuclear; ligações covalente, iónica e metálica; Lewis; octeto; geometria molecular; polaridade; parâmetros de ligação; hidrocarbonetos; grupos funcionais; forças de Van der Waals e pontes de hidrogénio; miscibilidade. \\
Propriedades da matéria & Gases e dispersões & Lei de Avogadro; volume molar; massa volúmica; fração molar; composição da troposfera; poluentes e efeito de estufa; composição quantitativa de soluções; preparação por dissolução e por diluição. \\
Propriedades da matéria & Transformações químicas & Quebra e formação de ligações; exotermicidade/endotermicidade; entalpia; energias de ligação; fotoquímica; reação do cloreto de prata; ozono estratosférico e troposférico; radicais livres. \\
Energia e conservação & Energia e movimentos & Centro de massa; energia cinética e potencial; trabalho de uma força constante em movimentos retilíneos; teorema da energia cinética; forças conservativas e não conservativas; energia mecânica; trabalho dissipativo e rendimento; plano inclinado; queda e ressalto. \\
Energia e conservação & Fenómenos elétricos & Corrente; diferença de potencial; resistência; circuitos em série e paralelo; potência elétrica; efeito Joule; conservação da energia em circuitos; gerador real; determinação experimental das características de uma pilha. \\
Energia e conservação & Fenómenos térmicos e radiação & Temperatura e equilíbrio térmico; energia cinética média das partículas; calor; condução, convecção e radiação; infravermelho; Primeira Lei; variação da energia interna; capacidade térmica mássica; entalpias mássicas de transição de fase; painel fotovoltaico e irradiância; determinação experimental de $c$ e da fusão do gelo; rendimento, Segunda Lei e degradação da energia. \\
\bottomrule
\end{longtable}
\normalsize

\begin{nota}
As Aprendizagens Essenciais definem o que o aluno deve saber e saber fazer; não constituem uma lista fechada de exemplos. Por isso, o livro acrescenta contextos, deduções simples, exercícios e projetos STEM/STEAM sempre que ajudam a consolidar as aprendizagens sem ultrapassar o nível do 10.º ano. A designação e numeração ``AL'' usada em algumas propostas é uma \textbf{organização editorial de apoio}; o documento de março de 2026 explicita as competências experimentais a desenvolver, mas deixa os temas concretos do trabalho prático abertos ao contexto da turma e da escola.
\end{nota}

\chapter*{Competências transversais nas AE de 2026}
\addcontentsline{toc}{chapter}{Competências transversais nas AE de 2026}
O referencial curricular atual não se limita a uma lista de conteúdos. Dá relevo à forma como o aluno \textbf{investiga, resolve problemas, trabalha experimentalmente e comunica ciência}. Por isso, estas competências aparecem de modo recorrente ao longo do livro.

\begin{table}[H]
\centering
\caption{Quatro dimensões que orientam o trabalho ao longo do livro.}
\renewcommand{\arraystretch}{1.18}\begin{tabularx}{\textwidth}{p{.23\textwidth}X}
\toprule
Dimensão & Como é trabalhada neste livro \\
\midrule
Conhecimento científico & Explicações, modelos, relações matemáticas, exemplos resolvidos e aplicação a situações novas. \\
Trabalho prático & Planeamento, segurança, escolha de instrumentos, medição, incerteza, tratamento de dados, avaliação de procedimentos e comunicação de conclusões. \\
Resolução de problemas & Problemas graduados, contextos do mundo real, tarefas integradoras de Física e Química e desafios STEM/STEAM. \\
Comunicação científica & Justificação de raciocínios, leitura e construção de tabelas/gráficos, relatórios curtos, argumentação baseada em evidências e avaliação crítica de fontes. \\
\bottomrule
\end{tabularx}\renewcommand{\arraystretch}{1.0}
\end{table}

\begin{metodo}[title=Investigar com fontes digitais e Inteligência Artificial]
Ferramentas digitais, simulações e sistemas de Inteligência Artificial generativa podem apoiar uma investigação, mas não substituem a validação científica. Sempre que uses uma ferramenta deste tipo: (1) formula uma pergunta verificável; (2) confirma números, unidades e afirmações em fontes credíveis; (3) distingue dados de interpretações; (4) regista as fontes utilizadas; e (5) reformula a conclusão com base nas evidências. Uma resposta convincente não é necessariamente uma resposta correta.
\end{metodo}

\begin{stem}
\textbf{Ciência cidadã e sustentabilidade.} Escolhe um problema observável na escola ou na comunidade — consumo elétrico, conforto térmico, qualidade do ar, iluminação ou produção fotovoltaica — e define uma questão investigável. Recolhe dados, identifica limitações, representa resultados e propõe uma intervenção mensurável. Esta estrutura aproxima o trabalho escolar da lógica de projetos STEAM e de ciência cidadã.
\end{stem}

\tableofcontents

\mainmatter

% =========================================================
\part{Ferramentas para pensar e medir em Ciência}
% =========================================================

\chapter{Medição, incerteza, dados e trabalho científico}

\begin{objetivos}
No final deste capítulo deverás ser capaz de:
\begin{itemize}
  \item distinguir grandeza, unidade e valor medido;
  \item utilizar corretamente unidades SI e notação científica;
  \item avaliar resolução, precisão e exatidão de uma medição;
  \item calcular médias, desvios, incertezas simples e erros relativos percentuais;
  \item representar dados em tabelas e gráficos e interpretar tendências;
  \item planear uma investigação simples, identificando variáveis e controlos.
\end{itemize}
\end{objetivos}

\section{Medir é comparar}
Uma medição compara uma grandeza com uma unidade. Escrever apenas ``12'' não é suficiente: \SI{12}{cm}, \SI{12}{s} e \SI{12}{g} representam realidades completamente diferentes.

\begin{table}[H]
\centering
\caption{Algumas grandezas e unidades SI relevantes.}
\begin{tabular}{lll}
\toprule
Grandeza & Símbolo habitual & Unidade SI \\
\midrule
Comprimento & $l$, $x$, $d$ & metro (m) \\
Massa & $m$ & quilograma (kg) \\
Tempo & $t$ & segundo (s) \\
Temperatura & $T$ & kelvin (K) \\
Quantidade de matéria & $n$ & mole (mol) \\
Corrente elétrica & $I$ & ampere (A) \\
Energia & $E$ & joule (J) \\
Potência & $P$ & watt (W) \\
\bottomrule
\end{tabular}
\end{table}

\section{Ordens de grandeza, notação científica e algarismos significativos}
Quando se comparam escalas muito diferentes, é útil escrever números em notação científica,
\[
x=a\times10^n,\qquad 1\leq |a|<10.
\]
A \textbf{ordem de grandeza} é a potência de dez que melhor representa a escala do valor. Este tipo de raciocínio é particularmente importante quando passamos de dimensões humanas para dimensões moleculares e atómicas.

\begin{exemplo}[title=Exemplo resolvido -- ordem de grandeza]
O diâmetro de um átomo típico é da ordem de \SI{1e-10}{m}, enquanto um cabelo humano tem uma espessura da ordem de \SI{1e-4}{m}. A razão entre as escalas é aproximadamente
\[
\frac{10^{-4}}{10^{-10}}=10^6.
\]
Assim, a espessura de um cabelo é cerca de \textbf{um milhão de vezes} superior ao diâmetro típico de um átomo.
\end{exemplo}

Os \textbf{algarismos significativos} comunicam a precisão com que um resultado é conhecido. Em multiplicações e divisões, o resultado não deve sugerir maior precisão do que a grandeza menos precisa. Em somas e subtrações, deve atender-se sobretudo à casa decimal menos precisa.

\begin{nota}
Uma calculadora pode apresentar dez casas decimais, mas isso não significa que todas tenham significado experimental. O resultado final deve ser arredondado de forma coerente com os dados e com a incerteza.
\end{nota}

\section{Incerteza de leitura e apresentação de uma medida}
Nenhuma medição é infinitamente precisa. Numa escala analógica, uma regra escolar comum é associar à leitura uma incerteza da ordem de metade da menor divisão da escala, se a leitura puder ser interpolada. Num aparelho digital, é frequente usar como estimativa inicial a unidade do último algarismo apresentado, salvo indicação do fabricante.

Uma medida pode ser apresentada sob a forma
\[
x=(x_{\rm med}\pm u)\,\text{unidade},
\]
onde $u$ é uma estimativa da \textbf{incerteza absoluta}.

\begin{exemplo}[title=Exemplo resolvido -- proveta e pipeta]
Uma proveta graduada de \SI{1}{mL} em \SI{1}{mL} permite ler aproximadamente até \SI{0.5}{mL}. Uma pipeta volumétrica de \SI{10.00}{mL} pode ter uma tolerância muito menor. Para medir \SI{10}{mL} com elevada exatidão, a pipeta é, em geral, a opção mais adequada. A escolha do instrumento faz parte do próprio problema experimental.
\end{exemplo}

\section{Resolução, precisão e exatidão}
A \textbf{resolução} de um instrumento é a menor variação que o instrumento permite distinguir. A \textbf{precisão} está associada à proximidade entre medições repetidas. A \textbf{exatidão} está associada à proximidade entre o resultado e um valor de referência aceite.

Para um conjunto de $N$ medições $x_i$, a média é
\[
\bar{x}=\frac{1}{N}\sum_{i=1}^{N}x_i.
\]
Um desvio individual pode ser escrito como
\[
d_i=|x_i-\bar{x}|.
\]
Em trabalhos escolares, uma estimativa simples da dispersão é o \textbf{desvio absoluto médio}:
\[
\bar d=\frac{1}{N}\sum_{i=1}^{N}|x_i-\bar{x}|.
\]
A incerteza relativa percentual pode ser estimada por
\[
u_r(\%)=\frac{u(x)}{|x|}\times 100.
\]
Para avaliar a exatidão relativamente a um valor de referência $x_{\rm ref}$, é útil o erro relativo percentual:
\[
\varepsilon_r(\%)=\frac{|x_{\rm exp}-x_{\rm ref}|}{|x_{\rm ref}|}\times100.
\]

\begin{exemplo}
Um grupo mede cinco vezes o período de um pêndulo e obtém, em segundos:
\[
1{,}42;\quad1{,}39;\quad1{,}41;\quad1{,}43;\quad1{,}40.
\]
A média é
\[
\bar t=\frac{1{,}42+1{,}39+1{,}41+1{,}43+1{,}40}{5}=\SI{1.41}{s}.
\]
Os desvios absolutos são $0{,}01$, $0{,}02$, $0{,}00$, $0{,}02$ e $0{,}01$ s. Assim,
\[
\bar d=\frac{0{,}06}{5}=\SI{0.012}{s}\approx\SI{0.01}{s}.
\]
Um resultado razoável, neste nível, é $t=(1{,}41\pm0{,}01)\,\mathrm{s}$.
\end{exemplo}

\subsection*{Desvio percentual e precisão}
Quando se pretende avaliar a dispersão de uma medição relativamente à média, pode definir-se, em contexto escolar, um desvio percentual:
\[
d_{\%}=\frac{|x_i-\bar{x}|}{\bar{x}}\times100.
\]
Valores pequenos e pouco dispersos são um indício de boa precisão. A precisão está sobretudo associada a erros aleatórios; a exatidão é particularmente sensível a erros sistemáticos.

\begin{table}[H]
\centering
\caption{Precisão, exatidão e tipos de erro.}
\renewcommand{\arraystretch}{1.18}\begin{tabularx}{\textwidth}{lXX}
\toprule
Conceito & Pergunta orientadora & Exemplo \\
\midrule
Precisão & As medições repetidas são próximas entre si? & leituras 5,01; 5,02; 5,01 g \\
Exatidão & O resultado está próximo do valor de referência? & massa volúmica obtida próxima do valor tabelado \\
Erro aleatório & A dispersão varia de medição para medição? & pequenas diferenças na leitura do menisco \\
Erro sistemático & Há um desvio persistente no mesmo sentido? & balança descalibrada \\
\bottomrule
\end{tabularx}\renewcommand{\arraystretch}{1.0}
\end{table}

\begin{exemplo}[title=Exemplo resolvido -- erro relativo e precisão]
A massa volúmica experimental de um metal é \SI{8.76}{g/cm^3}; o valor de referência é \SI{8.96}{g/cm^3}. Então
\[
\varepsilon_r=\frac{|8{,}76-8{,}96|}{8{,}96}\times100\approx2{,}2\%.
\]
Se várias determinações independentes estiverem muito próximas de \SI{8.76}{g/cm^3}, o procedimento pode ser \emph{preciso} mas apresentar uma exatidão limitada por um erro sistemático.
\end{exemplo}

\section{Gráficos e modelos}
Um gráfico não é decoração: é uma forma de testar relações entre grandezas. Antes de construir um gráfico, importa decidir que grandezas se representam, que incerteza acompanha cada ponto e que modelo pode ser testado.

\begin{figure}[H]
\centering
\begin{tikzpicture}[scale=.78]
\foreach \cx/\titulo in {0/{preciso e exato},5/{preciso, pouco exato},10/{pouco preciso}}{
  \draw[gray!55] (\cx,0) circle (1.55);
  \draw[gray!55] (\cx,0) circle (1.05);
  \draw[gray!55] (\cx,0) circle (.52);
  \draw[eured,thick] (\cx-.18,0)--(\cx+.18,0); \draw[eured,thick] (\cx,-.18)--(\cx,.18);
  \node[font=\small\bfseries] at (\cx,-2.05) {\titulo};
}
\foreach \x/\y in {-.18/.12,.12/.18,.05/-.14,-.1/-.08} \fill[eublue] (\x,\y) circle (2.1pt);
\foreach \x/\y in {5.78/.65,5.60/.77,5.74/.82,5.65/.58} \fill[eublue] (\x,\y) circle (2.1pt);
\foreach \x/\y in {9.2/.7,10.7/.5,9.55/-.9,10.5/-.65,10.1/.95} \fill[eublue] (\x,\y) circle (2.1pt);
\end{tikzpicture}
\caption{Representação qualitativa da diferença entre precisão e exatidão.}
\end{figure}

\begin{conceito}
\textbf{Incerteza e comparação.} Dois resultados experimentais não devem ser comparados apenas pelos seus valores centrais. Quando as incertezas são relevantes, deve analisar-se se os intervalos de valores plausíveis se sobrepõem e se a diferença observada é maior do que a resolução do método.
\end{conceito} Sempre que possível:
\begin{enumerate}
  \item coloca a variável independente no eixo horizontal;
  \item coloca a variável dependente no eixo vertical;
  \item identifica grandezas e unidades;
  \item usa uma escala uniforme;
  \item procura tendências e avalia se um modelo linear é adequado.
\end{enumerate}

\begin{center}
\begin{tikzpicture}
\begin{axis}[
width=10cm,height=6cm,
xlabel={$d$/m},ylabel={$E_{\rm c}$/J},
grid=major,axis lines=left]
\addplot[only marks,mark=*] coordinates{(0.2,0.5)(0.4,1.0)(0.6,1.52)(0.8,2.02)(1.0,2.49)};
\addplot[domain=0:1.05]{2.5*x};
\end{axis}
\end{tikzpicture}
\end{center}

\subsection*{Declive, ordenada na origem e ajuste linear}
Quando um modelo prevê uma relação $y=ax+b$, o declive $a$ e a ordenada na origem $b$ têm significado físico.
 A qualidade do ajuste pode ser apreciada pela proximidade dos pontos à reta e pela existência, ou não, de um padrão nos desvios.

\begin{figure}[H]
\centering
\begin{tikzpicture}
\begin{axis}[width=11cm,height=6.2cm,xlabel={$x$ / unidade},ylabel={$y$ / unidade},xmin=0,xmax=6,ymin=0,ymax=15,grid=major]
\addplot[only marks,mark=*,eublue,error bars/.cd,y dir=both,y explicit] table[row sep=\\,y error=err]{x y err\\ 1 2.8 .35\\ 2 5.0 .40\\ 3 7.2 .45\\ 4 9.3 .35\\ 5 11.7 .50\\};
\addplot[domain=.6:5.4,thick,euorange]{2.2*x+0.5};
\legend{dados experimentais,ajuste linear}
\end{axis}
\end{tikzpicture}
\caption{Exemplo de representação de dados com barras de incerteza e uma reta de ajuste.}
\end{figure}

\begin{leituragrafico}
Se $y=ax+b$, pergunta sempre: \textbf{qual é a unidade de $a$?} O que significa $b$ quando $x=0$? Faz sentido extrapolar a reta para fora do intervalo medido? Uma relação linear observada num intervalo não prova que continue linear em qualquer condição.
\end{leituragrafico} A unidade do declive é a unidade de $y$ dividida pela unidade de $x$. Em dados experimentais, não se deve ligar mecanicamente os pontos: procura-se a reta que melhor representa a tendência global.

\begin{metodo}
\textbf{Checklist experimental:} definir a questão-problema; identificar variável independente, dependente e variáveis de controlo; prever o resultado; escolher instrumentos; repetir medições; registar unidades e incertezas; representar os dados; comparar com o modelo; discutir limitações e propor melhorias.
\end{metodo}

\begin{mundo}
Sensores de um telemóvel medem aceleração, rotação, intensidade luminosa e campo magnético. Uma atividade STEM pode transformar o telemóvel num pequeno laboratório: recolhe-se um conjunto de dados, exporta-se para uma folha de cálculo, constrói-se um gráfico e compara-se o comportamento observado com um modelo físico.
\end{mundo}

\begin{exercicios}
\begin{enumerate}
  \item Converte: (a) \SI{72}{km/h} para \si{m/s}; (b) \SI{3.5}{mg} para kg; (c) \SI{2.4e-6}{m} para \si{\micro m}.
  \item Uma balança apresenta leituras de \SI{52.43}{g}, \SI{52.41}{g}, \SI{52.44}{g} e \SI{52.42}{g}. Determina a média e o desvio absoluto médio.
  \item Um valor experimental para a massa volúmica do alumínio é \SI{2.64}{g/cm^3}; o valor de referência é \SI{2.70}{g/cm^3}. Calcula o erro relativo percentual.
  \item Explica porque é possível ter resultados muito precisos mas pouco exatos.
  \item Num ensaio, a energia cinética varia linearmente com a distância. O declive do gráfico $E_c(d)$ é \SI{3.2}{J/m}. Que unidade equivalente tem esse declive? Que grandeza física sugere?
\end{enumerate}
\end{exercicios}

% =========================================================
\part{Química — Estrutura, propriedades e transformações da matéria}
% =========================================================

\chapter{Átomos, isótopos, quantidade de matéria e composição}

\begin{objetivos}
\begin{itemize}
  \item descrever a constituição de átomos e iões usando $Z$ e $A$;
  \item distinguir isótopos e interpretar massa atómica relativa média;
  \item relacionar número de entidades, quantidade de matéria e massa;
  \item efetuar cálculos de composição de misturas e soluções simples.
\end{itemize}
\end{objetivos}

\section{Do macroscópico ao microscópico}
A matéria é constituída por entidades extremamente pequenas. Um átomo contém um núcleo, onde se encontram protões e neutrões, rodeado por uma nuvem eletrónica.

\begin{center}
\begin{tikzpicture}[scale=1.05]
\fill[eublue!10] (0,0) circle (2.45);
\fill[euorange] (0,0) circle (.52);
\node[white] at (0,0) {núcleo};
\draw[eublue,thick] (0,0) circle (1.4);
\draw[eublue,thick,dashed] (0,0) circle (2.1);
\fill[eucyan] (1.4,0) circle (3pt) node[above right,black] {$e^-$};
\fill[eucyan] (-1.48,1.48) circle (3pt) node[above left,black] {$e^-$};
\node[align=center] at (0,-3.05) {Representação esquemática; não está à escala.};
\end{tikzpicture}
\end{center}

O \textbf{número atómico}, $Z$, é o número de protões. Num átomo neutro é também o número de eletrões. O \textbf{número de massa}, $A$, é a soma do número de protões e neutrões:
\[
A=Z+N.
\]
A notação de um nuclídeo pode escrever-se como $^{A}_{Z}\mathrm{X}$.

\begin{table}[H]
\centering
\caption{Partículas subatómicas no modelo introdutório do átomo.}
\rowcolors{2}{eulightblue}{white}
\begin{tabular}{lccc}
\toprule
Partícula & Carga relativa & Massa aproximada / u & Localização dominante \\
\midrule
Protão & $+1$ & $1$ & núcleo \\
Neutrão & $0$ & $1$ & núcleo \\
Eletrão & $-1$ & $5{,}5\times10^{-4}$ & nuvem eletrónica \\
\bottomrule
\end{tabular}
\end{table}

\begin{conceito}
Quase toda a massa do átomo está concentrada no núcleo, mas quase todo o seu volume corresponde à região ocupada pela nuvem eletrónica. Por isso, um desenho do átomo à escala real seria muito diferente das representações escolares.
\end{conceito}

\section{Escala atómica e nanotecnologia}
A escala atómica situa-se tipicamente entre cerca de $10^{-10}$ e $10^{-9}$ m. Estruturas desta dimensão não são observáveis com microscópios óticos comuns porque a resolução é limitada pelo comprimento de onda da luz visível. Técnicas como microscopia eletrónica ou microscopia de força atómica permitem investigar superfícies e estruturas muito menores.

\begin{center}
\begin{tikzpicture}[x=1.35cm,y=1cm]
\draw[-{Stealth}] (0,0)--(8.72,0) node[above right] {escala};
\foreach \x/\lab in {0.3/$10^{-10}$ m,2.2/$10^{-9}$ m,4.2/$10^{-6}$ m,6.3/$10^{-4}$ m,8.2/$10^{-2}$ m}
  \draw (\x,.12)--(\x,-.12) node[below=4pt] {\lab};
\node[align=center] at (.7,.75) {átomos};
\node[align=center] at (2.4,.75) {nanopartículas};
\node[align=center] at (4.3,.75) {bactérias};
\node[align=center] at (6.4,.75) {cabelo};
\node[align=center] at (8.2,.75) {objetos\\quotidianos};
\end{tikzpicture}
\end{center}

\begin{mundo}
Na nanotecnologia, alterar tamanho e superfície pode modificar propriedades óticas, elétricas ou catalíticas. Nanopartículas aparecem em sensores, materiais funcionais, medicina e tecnologias energéticas. A avaliação de benefícios e riscos exige relacionar propriedades microscópicas com impactos tecnológicos e ambientais.
\end{mundo}


\begin{figure}[H]
\centering
\begin{tikzpicture}[x=1cm,y=1cm]
  \draw[fill=eugrey!70,draw=eudark] (0,0) rectangle (10,.75);
  \node[left] at (0,.38) {fio de cabelo};
  \draw[<->,thick,eublue] (0,-.35)--(10,-.35);
  \node at (5,-.72) {$\approx\SI{90}{\micro m}$};
  \draw[fill=eucyan!65,draw=eucyan!90!black] (0,-1.45) rectangle (.11,-.83);
  \node[left,align=right] at (0,-1.14) {nanotubo de\\carbono};
  \draw[<->,thick,euorange] (0,-1.8)--(1.1,-1.8);
  \node[anchor=west] at (1.3,-1.8) {representação ampliada; diâmetros típicos da ordem dos nm};
\end{tikzpicture}
\caption{Comparação conceptual entre as escalas micrométrica e nanométrica. As larguras desenhadas não estão à mesma escala.}
\end{figure}

\begin{exemplo}[title=Exemplo resolvido -- comparar escalas]
Um fio de cabelo tem diâmetro aproximado de \SI{90}{\micro m} e um nanotubo tem diâmetro de \SI{10}{nm}. Como
\[
\SI{90}{\micro m}=9{,}0\times10^{-5}\,\mathrm{m},\qquad
\SI{10}{nm}=1{,}0\times10^{-8}\,\mathrm{m},
\]
a razão entre os diâmetros é
\[
\frac{9{,}0\times10^{-5}}{1{,}0\times10^{-8}}=9{,}0\times10^3.
\]
O fio de cabelo é, neste exemplo, cerca de \textbf{9000 vezes mais largo} do que o nanotubo.
\end{exemplo}

\section{Isótopos e massa atómica relativa}
A \textbf{unidade de massa atómica}, símbolo $\mathrm{u}$, é definida a partir do isótopo carbono-12:
\[
1\,\mathrm{u}=\frac{1}{12}m(^{12}\mathrm{C})\approx1{,}6605\times10^{-27}\,\mathrm{kg}.
\]
Isótopos são átomos do mesmo elemento, portanto com o mesmo $Z$, mas com diferente número de neutrões. A massa atómica relativa média de um elemento resulta da média ponderada das massas isotópicas pelas respetivas abundâncias naturais.

\begin{figure}[H]
\centering
\begin{tikzpicture}
\begin{axis}[ybar,width=10.5cm,height=5.8cm,ymin=0,ymax=85,bar width=22pt,
ylabel={abundância / \%},xlabel={isótopo},symbolic x coords={$^{35}$X,$^{37}$X},xtick=data,nodes near coords,nodes near coords align={vertical},grid=major]
\addplot[eublue,fill=eublue!45] coordinates {($^{35}$X,75.8) ($^{37}$X,24.2)};
\end{axis}
\end{tikzpicture}
\caption{Abundâncias isotópicas usadas no exemplo. A massa atómica relativa é uma média ponderada.}
\end{figure}

\begin{exemplo}
Considere um elemento com dois isótopos de massas relativas 35 e 37, com abundâncias 75,8\% e 24,2\%. Uma estimativa da massa atómica relativa média é
\[
A_r\approx 35(0{,}758)+37(0{,}242)=35{,}48.
\]
\end{exemplo}

\section{A mole: uma ponte entre átomos e laboratório}
Uma mole contém o número de Avogadro de entidades:
\[
\NA=6{,}022\times10^{23}\ \mathrm{mol^{-1}}.
\]
As relações fundamentais são
\[
N=n\NA,\qquad n=\frac{m}{M},
\]
onde $N$ é o número de entidades, $n$ a quantidade de matéria, $m$ a massa e $M$ a massa molar.

\begin{exemplo}
Quantas moléculas existem em \SI{9.0}{g} de água? Usando $M(\ce{H2O})\approx\SI{18.0}{g/mol}$:
\[
n=\frac{9{,}0}{18{,}0}=\SI{0.50}{mol}
\]
e
\[
N=0{,}50\times6{,}022\times10^{23}\approx3{,}0\times10^{23}\ \text{moléculas}.
\]
\end{exemplo}


\subsection*{Massa molecular relativa e massa molar}
A \textbf{massa molecular relativa}, $M_r$, é a soma das massas atómicas relativas dos átomos que constituem uma molécula. É uma grandeza adimensional. Por exemplo,
\[
M_r(\ce{CO2})=A_r(\ce C)+2A_r(\ce O)\approx12{,}01+2(16{,}00)=44{,}01.
\]
A \textbf{massa molar}, $M$, é a massa de uma mole de entidades e exprime-se habitualmente em \si{g/mol}. Assim,
\[
M(\ce{CO2})\approx\SI{44.01}{g/mol}.
\]
É importante distinguir as duas grandezas: $M_r$ compara massas numa escala relativa microscópica; $M$ liga essa descrição microscópica à massa macroscópica medida no laboratório.

\begin{exemplo}[title={Exemplo resolvido -- massa, quantidade de matéria e entidades}]
Um recipiente contém \SI{0.800}{mol} de sulfureto de hidrogénio, \ce{H2S}. Usando $A_r(\ce H)=1{,}01$ e $A_r(\ce S)=32{,}06$,
\[
M(\ce{H2S})=2(1{,}01)+32{,}06=\SI{34.08}{g/mol}.
\]
A massa da amostra é
\[
m=nM=0{,}800\times34{,}08=\SI{27.3}{g}.
\]
O número de moléculas é
\[
N=nN_A=0{,}800\times6{,}022\times10^{23}=4{,}82\times10^{23}.
\]
Como cada molécula contém dois átomos de hidrogénio e um de enxofre,
\[
n(\ce H)=\SI{1.600}{mol},\qquad n(\ce S)=\SI{0.800}{mol}.
\]
\end{exemplo}

\begin{figure}[H]
\centering
\begin{minipage}{.48\textwidth}
\centering
\begin{tikzpicture}
\begin{axis}[width=\linewidth,height=5.2cm,xlabel={$n$/mol},ylabel={$N/(10^{23})$},xmin=0,xmax=2.1,ymin=0,ymax=13,grid=major]
\addplot[thick,eublue,domain=0:2]{6.022*x};
\addplot[only marks,mark=*,euorange] coordinates{(.5,3.011)(1,6.022)(1.5,9.033)};
\end{axis}
\end{tikzpicture}
\end{minipage}\hfill
\begin{minipage}{.48\textwidth}
\centering
\begin{tikzpicture}
\begin{axis}[width=\linewidth,height=5.2cm,xlabel={$m$/g},ylabel={$n$/mol},xmin=0,xmax=50,ymin=0,ymax=1.2,grid=major]
\addplot[thick,eugreen,domain=0:50]{x/44.0};
\addplot[only marks,mark=square*,euorange] coordinates{(11,.25)(22,.5)(44,1)};
\end{axis}
\end{tikzpicture}
\end{minipage}
\caption{Relações lineares fundamentais: $N=nN_A$ e, para uma substância de massa molar fixa, $n=m/M$. O gráfico da direita usa $M=\SI{44.0}{g/mol}$.}
\end{figure}

\begin{leituragrafico}
No gráfico $N(n)$, o declive é $N_A$. No gráfico $n(m)$, o declive é $1/M$. Assim, um gráfico experimental pode ser usado não apenas para prever valores, mas também para determinar uma constante física ou uma propriedade da substância.
\end{leituragrafico}

\begin{atividade}[title=AL 1.1 -- Volume e número de moléculas de uma gota de água]
\textbf{Questão:} como estimar a massa, o volume e o número de moléculas de uma única gota de água?

\textbf{Ideia experimental.} Uma gota é demasiado pequena para ser medida com boa precisão usando muitos instrumentos escolares. Mede-se, por isso, a massa e o volume de um número elevado de gotas, por exemplo 100, e divide-se pelo número de gotas.

\textbf{Procedimento possível:}
\begin{enumerate}
  \item selecionar uma balança e um recipiente seco; medir a massa do recipiente;
  \item adicionar um número conhecido de gotas e medir a nova massa;
  \item medir o volume do mesmo número de gotas com material volumétrico adequado;
  \item calcular massa e volume médios de uma gota;
  \item usar $n=m/M$ e $N=nN_A$ para estimar o número de moléculas;
  \item apresentar os resultados com algarismos significativos coerentes e discutir incertezas.
\end{enumerate}

\textbf{Variáveis críticas:} tamanho das gotas, água retida na ponta do conta-gotas, evaporação, gotas perdidas e resolução do material volumétrico.
\end{atividade}

\begin{exemplo}[title=Exemplo resolvido -- moléculas numa gota]
Cem gotas de água têm massa total \SI{4.80}{g}. A massa média de uma gota é
\[
m_g=\frac{4{,}80}{100}=\SI{0.0480}{g}.
\]
Como $M(\ce{H2O})=\SI{18.0}{g/mol}$,
\[
n=\frac{0{,}0480}{18{,}0}=2{,}67\times10^{-3}\ \mathrm{mol}.
\]
Logo,
\[
N=nN_A\approx2{,}67\times10^{-3}\times6{,}022\times10^{23}
  \approx1{,}61\times10^{21}
\]
moléculas por gota.
\end{exemplo}

\begin{figure}[H]
\centering
\begin{tikzpicture}[node distance=1.4cm,>=Stealth]
\tikzset{q/.style={rounded corners,draw=eublue,fill=eulightblue,minimum width=3.0cm,minimum height=1cm,align=center}}
\node[q] (mass) {massa $m$};
\node[q,right=of mass] (mol) {quantidade\\$n$};
\node[q,right=of mol] (ent) {entidades\\$N$};
\draw[->,thick,euorange] (mass)-- node[above] {$\div M$} (mol);
\draw[<-,thick,euorange] (mass)-- node[below] {$\times M$} (mol);
\draw[->,thick,eugreen] (mol)-- node[above] {$\times N_A$} (ent);
\draw[<-,thick,eugreen] (mol)-- node[below] {$\div N_A$} (ent);
\end{tikzpicture}
\caption{Mapa de conversões entre massa, quantidade de matéria e número de entidades.}
\end{figure}

\section{Composição de misturas}
A fração mássica de um componente é
\[
w_i=\frac{m_i}{m_{\rm mistura}},
\]
e a percentagem mássica é $100w_i$.

\begin{mundo}
Rótulos alimentares, ligas metálicas, combustíveis e medicamentos utilizam diferentes formas de indicar composição. Saber converter entre massa, quantidade de matéria e percentagem permite interpretar informação técnica e tomar decisões fundamentadas.
\end{mundo}

\begin{stem}
\textbf{Projeto: quantos átomos há numa folha de alumínio?} Mede a massa de uma pequena folha de alumínio e estima o número de átomos presentes. Em seguida, mede comprimento, largura e espessura para estimar o volume e a massa volúmica. Compara com um valor tabelado e discute as principais fontes de incerteza.
\end{stem}

\begin{exercicios}
\begin{enumerate}
  \item Para $^{23}_{11}\ce{Na}$, indica o número de protões, neutrões e eletrões no átomo neutro e no ião \ce{Na+}.
  \item Uma amostra contém \SI{0.250}{mol} de \ce{CO2}. Determina o número de moléculas e o número total de átomos de oxigénio.
  \item Determina a massa correspondente a \SI{0.150}{mol} de \ce{CaCO3}, usando $M=\SI{100.1}{g/mol}$.
  \item Uma liga de \SI{250}{g} contém \SI{18}{g} de níquel. Determina a percentagem mássica de níquel.
  \item Um elemento apresenta dois isótopos, de massas relativas 10,0 e 11,0, com abundâncias 19,9\% e 80,1\%. Calcula a massa atómica relativa média.
  \item \textbf{Integração.} Uma garrafa contém \SI{500}{mL} de uma solução com massa volúmica \SI{1.05}{g/mL} e 8,0\% (m/m) de soluto. Estima a massa de soluto presente.
  \item Ordena por ordem crescente de comprimento: \SI{0.1}{nm}, \SI{2}{\micro m}, \SI{0.08}{mm} e \SI{3}{cm}.
  \item Uma amostra contém $1{,}20\times10^{22}$ átomos de cobre. Determina a quantidade de matéria correspondente.
  \item Cem gotas de água têm um volume total de \SI{4.6}{mL}. Estima o volume de uma gota e indica duas razões para repetir a experiência várias vezes.
  \item Ao escrever uma assinatura com um lápis de grafite consomem-se cerca de \SI{1.0}{mg} de carbono. Admitindo grafite pura e $M(\ce C)=\SI{12.01}{g/mol}$, estima o número de átomos de carbono depositados.
  \item Uma liga usada num contexto dentário contém 72\% (m/m) de prata. Numa amostra de \SI{30}{mg}, determina a quantidade de matéria e o número de átomos de prata. Usa $M(\ce{Ag})=\SI{107.87}{g/mol}$.
  \item Uma botija contém $1{,}807\times10^{25}$ átomos de oxigénio sob a forma de \ce{O2}. Determina a quantidade de matéria de \ce{O2}.
  \item Considera uma gota de água de volume \SI{0.050}{mL}. Sabendo que $\rho(\ce{H2O})=\SI{0.997}{g/mL}$ e $M(\ce{H2O})=\SI{18.0}{g/mol}$, estima o número total de átomos existentes na gota.
  \item Calcula a massa total de uma mistura constituída por $3{,}01\times10^{23}$ moléculas de \ce{CO2}, \SI{0.500}{mol} de \ce{N2} e \SI{1.200}{mol} de \ce{O2}. Usa $M(\ce{CO2})=\SI{44.0}{g/mol}$, $M(\ce{N2})=\SI{28.0}{g/mol}$ e $M(\ce{O2})=\SI{32.0}{g/mol}$.
\end{enumerate}
\end{exercicios}

\chapter{Energia dos eletrões, luz e Tabela Periódica}

\begin{objetivos}
\begin{itemize}
  \item relacionar energia de um fotão com frequência e comprimento de onda;
  \item interpretar espectros de emissão e absorção como assinaturas atómicas;
  \item usar o modelo de Bohr como modelo histórico para interpretar o átomo de hidrogénio, reconhecendo os seus limites;
  \item relacionar transições eletrónicas, séries espetrais e energia de ionização;
  \item compreender níveis, subníveis e orbitais no modelo quântico;
  \item interpretar qualitativamente informação de espetroscopia fotoeletrónica (PES);
  \item escrever configurações eletrónicas até $Z=23$ e aplicar os princípios de Aufbau, Pauli e Hund;
  \item relacionar configuração eletrónica com grupos, períodos, blocos e propriedades periódicas.
\end{itemize}
\end{objetivos}

\section{Radiação eletromagnética e energia}
As radiações eletromagnéticas formam um espetro contínuo. Diferem na frequência $f$ e no comprimento de onda $\lambda$, relacionados por
\[
c=\lambda f.
\]
Quanto maior a frequência, maior a energia de cada fotão.

\subsection*{Comprimento de onda, frequência e período}
Uma onda periódica pode ser caracterizada por:
\begin{itemize}
  \item \textbf{comprimento de onda} $\lambda$: distância entre dois pontos consecutivos no mesmo estado de vibração;
  \item \textbf{período} $T$: tempo necessário para completar um ciclo;
  \item \textbf{frequência} $f$: número de ciclos por segundo, com $f=1/T$.
\end{itemize}
Num meio em que a onda se propaga com velocidade $v$,
\[
v=\lambda f.
\]
No vazio, para radiação eletromagnética, $v=c$.

\begin{figure}[H]
\centering
\begin{tikzpicture}
\begin{axis}[width=11cm,height=5.4cm,xlabel={$x$},ylabel={amplitude},xmin=0,xmax=12.6,ymin=-1.35,ymax=1.35,xtick=\empty,ytick={-1,0,1},grid=none,axis lines=middle]
\addplot[domain=0:12.6,samples=180,thick,eublue]{sin(deg(x))};
\draw[<->,very thick,euorange] (axis cs:1.57,1.12)--(axis cs:7.85,1.12);
\node[euorange,fill=white,inner sep=2pt] at (axis cs:4.71,1.12) {$\lambda$};
\end{axis}
\end{tikzpicture}
\caption{O comprimento de onda mede uma repetição espacial completa da onda.}
\end{figure}

\begin{exemplo}[title={Exemplo resolvido -- frequência, período e energia}]
A sensibilidade do olho humano é elevada perto de $\lambda=\SI{555}{nm}$. A frequência correspondente é
\[
f=\frac{c}{\lambda}=\frac{3{,}00\times10^8}{555\times10^{-9}}
 \approx5{,}41\times10^{14}\,\mathrm{Hz}.
\]
O período é
\[
T=\frac1f\approx1{,}85\times10^{-15}\,\mathrm{s},
\]
e a energia de um fotão é
\[
E=hf\approx3{,}58\times10^{-19}\,\mathrm{J}.
\]
\end{exemplo}

\begin{table}[H]
\centering
\caption{Regiões do espetro eletromagnético e aplicações típicas.}
\renewcommand{\arraystretch}{1.18}\begin{tabularx}{\textwidth}{lXX}
\toprule
Região & Ordem de grandeza de $\lambda$ & Exemplos de aplicação \\
\midrule
Ondas de rádio & m a km & comunicações \\
Micro-ondas & mm a cm & comunicações, radar, aquecimento \\
Infravermelho & $\mu$m a mm & termografia, comandos, sensores \\
Visível & cerca de 400--700 nm & visão, iluminação, espectroscopia \\
Ultravioleta & dezenas a centenas de nm & fluorescência, processos fotoquímicos \\
Raios X & aproximadamente $10^{-11}$--$10^{-8}$ m & imagiologia e análise de materiais \\
\bottomrule
\end{tabularx}\renewcommand{\arraystretch}{1.0}
\end{table}

\section{Luz como radiação eletromagnética}
A luz transporta energia em quantidades discretas chamadas fotões. A energia de um fotão é
\[
E=hf=\frac{hc}{\lambda},
\]
onde $h=6{,}626\times10^{-34}\,\mathrm{J\,s}$ e $c\approx3{,}00\times10^8\,\mathrm{m\,s^{-1}}$.

\begin{center}
\begin{tikzpicture}[x=1.3cm,y=1cm]
\draw[-{Stealth}] (0,0)--(9,0) node[right] {$\lambda$ aumenta};
\foreach \x/\txt in {0.5/UV,2.3/Violeta,3.5/Azul,4.8/Verde,6.1/Amarelo,7.3/Vermelho,8.5/IV}
  \node[rotate=35,anchor=west] at (\x,.1) {\txt};
\draw[eublue,very thick] (1.6,-.25)--(7.9,-.25);
\node at (4.8,-.75) {região visível (esquema qualitativo)};
\end{tikzpicture}
\end{center}

\begin{figure}[H]
\centering
\begin{tikzpicture}
\begin{axis}[width=10.8cm,height=5.8cm,xlabel={$f/(10^{14}\,\mathrm{Hz})$},ylabel={$E/(10^{-19}\,\mathrm{J})$},xmin=0,xmax=9,ymin=0,ymax=6.2,grid=major]
\addplot[domain=0:9,thick,eublue]{0.6626*x};
\node[anchor=west,font=\small] at (axis cs:5.1,2.35) {$E=hf$};
\end{axis}
\end{tikzpicture}
\caption{Energia de um fotão em função da frequência. A proporcionalidade é direta.}
\end{figure}

\begin{exemplo}[title=Exemplo resolvido -- energia de um fotão]
Para uma radiação verde de \SI{530}{nm},
\[
E=\frac{hc}{\lambda}
 =\frac{(6{,}626\times10^{-34})(3{,}00\times10^8)}
 {530\times10^{-9}}
 \approx3{,}75\times10^{-19}\ \mathrm{J}.
\]
Uma mole destes fotões transportaria aproximadamente
\[
E_m=EN_A\approx\SI{226}{kJ/mol}.
\]
\end{exemplo}

\section{Espectros e identificação de elementos}
Um \textbf{espetro contínuo} contém uma sequência praticamente ininterrupta de comprimentos de onda. Num gás atómico excitado, surgem frequentemente \textbf{espetros de emissão descontínuos}: apenas certas riscas são observadas. Se uma radiação contínua atravessar um gás frio, podem aparecer riscas escuras de absorção nas frequências que o gás absorve.
Quando um eletrão transita de um nível de energia superior para um inferior, o átomo pode emitir um fotão cuja energia corresponde à diferença entre níveis:
\[
\Delta E=E_{\rm superior}-E_{\rm inferior}=hf.
\]
Como cada elemento possui um conjunto próprio de níveis de energia, o seu espectro funciona como uma assinatura.

\begin{center}
\begin{tikzpicture}[x=1.2cm,y=.8cm]
\draw[thick] (0,0) rectangle (8,1.2);
\fill[black] (0,0) rectangle (8,1.2);
\foreach \x/\c in {1.0/eucyan,2.5/eublue,4.4/eugreen,6.2/euorange,7.1/red}
  \draw[\c,line width=2.5pt] (\x,.08)--(\x,1.12);
\node at (4,-.45) {Espectro de emissão esquemático};
\end{tikzpicture}
\end{center}


\begin{figure}[H]
\centering
\begin{tabular}{c}
\textbf{Emissão do hidrogénio}\\[2mm]
\pgfspectra[element=H,width=.82\textwidth,height=1.1cm,axis,axis step=100,begin=380,end=750,axis ticks=4,axis unit=nm]\\[4mm]
\textbf{Absorção do hidrogénio}\\[2mm]
\pgfspectra[element=H,absorption,width=.82\textwidth,height=1.1cm,axis,axis step=100,begin=380,end=750,axis ticks=4,axis unit=nm]
\end{tabular}
\caption{Riscas de emissão e de absorção do mesmo elemento aparecem nas mesmas posições espetrais.}
\end{figure}

\begin{conceito}
Um gás atómico pode absorver precisamente fotões cujas energias correspondem a transições permitidas entre níveis eletrónicos. Por isso, as frequências das riscas escuras de absorção coincidem com frequências que o mesmo elemento pode emitir quando excitado. Esta correspondência permite identificar elementos em atmosferas estelares e em amostras laboratoriais.
\end{conceito}

\subsection*{Níveis de energia no átomo de hidrogénio}
A quantização significa que o eletrão ligado ao átomo não pode ter qualquer energia: existem níveis permitidos. Uma transição descendente liberta um fotão; uma transição ascendente exige absorção de energia.

\begin{center}
\begin{tikzpicture}[x=1cm,y=.9cm]
\foreach \y/\lab in {0/$n=1$,1.5/$n=2$,2.5/$n=3$,3.2/$n=4$}{
  \draw[eublue,thick] (0,\y)--(5,\y);
  \node[right] at (5,\y) {\lab};
}
\draw[eured,very thick,-{Stealth}] (1.4,3.2)--(1.4,1.5);
\node[left] at (1.4,2.35) {emissão};
\draw[eugreen,very thick,-{Stealth}] (3.5,0)--(3.5,2.5);
\node[right] at (3.5,1.25) {absorção};
\end{tikzpicture}
\end{center}


\subsection*{O modelo de Bohr: um passo histórico decisivo}
O modelo proposto por Niels Bohr em 1913 não corresponde ao modelo atómico atualmente aceite, mas desempenhou um papel essencial na interpretação do espetro do hidrogénio. No modelo de Bohr, o eletrão só pode ocupar determinados \textbf{estados estacionários}, associados a valores bem definidos de energia. A passagem entre dois estados é acompanhada pela absorção ou emissão de energia.

Para o átomo de hidrogénio, a energia do eletrão no nível de número quântico principal $n$ é dada por
\[
E_n=-\frac{2{,}18\times10^{-18}}{n^2}\ \mathrm{J},\qquad n=1,2,3,\ldots
\]
O valor $E=0$ corresponde, por convenção, ao eletrão livre infinitamente afastado do protão. Por isso, todos os estados ligados apresentam energia negativa. O estado $n=1$ é o \textbf{estado fundamental}; os estados $n>1$ são \textbf{estados excitados}.

\begin{figure}[H]
\centering
\begin{tikzpicture}[scale=.82,>=Stealth]
% Painel absorção
\begin{scope}[xshift=-4.8cm]
\node[font=\bfseries\color{eublue}] at (0,3.4) {Absorção};
\fill[eublue!75] (0,0) circle (.28);
\node[white,font=\scriptsize\bfseries] at (0,0) {$+$};
\draw[eublue!35,dashed] (0,0) circle (1.25);
\draw[eublue!65,dashed] (0,0) circle (2.15);
\fill[eured] (1.25,0) circle (.13);
\node[font=\scriptsize] at (1.25,0) {$-$};
\draw[very thick,eugreen,-{Stealth}] (1.32,.16) to[out=15,in=165] (2.03,.30);
\draw[euorange,decorate,decoration={snake,segment length=5pt,amplitude=2pt},-{Stealth}] (-3.0,1.7)--(-1.05,.75);
\node[align=center,font=\scriptsize] at (-2.25,2.25) {fotão\\absorvido};
\node[align=center,font=\scriptsize] at (0,-2.8) {o eletrão passa para\\um nível de maior energia};
\end{scope}
% Painel emissão
\begin{scope}[xshift=4.8cm]
\node[font=\bfseries\color{eublue}] at (0,3.4) {Emissão};
\fill[eublue!75] (0,0) circle (.28);
\node[white,font=\scriptsize\bfseries] at (0,0) {$+$};
\draw[eublue!35,dashed] (0,0) circle (1.25);
\draw[eublue!65,dashed] (0,0) circle (2.15);
\fill[eured] (2.15,0) circle (.13);
\node[font=\scriptsize] at (2.15,0) {$-$};
\draw[very thick,eugreen,-{Stealth}] (2.03,.30) to[out=165,in=15] (1.32,.16);
\draw[euorange,decorate,decoration={snake,segment length=5pt,amplitude=2pt},-{Stealth}] (1.05,.75)--(3.0,1.7);
\node[align=center,font=\scriptsize] at (2.45,2.25) {fotão\\emitido};
\node[align=center,font=\scriptsize] at (0,-2.8) {o eletrão passa para\\um nível de menor energia};
\end{scope}
\end{tikzpicture}
\caption{Representação histórica, segundo o modelo de Bohr, dos processos de absorção e emissão. As ``órbitas'' servem aqui apenas para visualizar a quantização da energia.}
\end{figure}

Numa transição eletrónica,
\[
\Delta E=E_f-E_i.
\]
Se $\Delta E>0$, o átomo absorveu energia; se $\Delta E<0$, libertou energia. A energia do fotão envolvido é sempre positiva:
\[
E_{\text{fotão}}=|\Delta E|=hf=\frac{hc}{\lambda}.
\]

\begin{figure}[H]
\centering
\begin{tikzpicture}[x=1.25cm,y=.82cm,>=Stealth]
\draw[->,eudark] (0,-2.65)--(0,1.15) node[above] {$E$};
\foreach \y/\lab/\val in {-2.18/$n=1$/-2.18,-0.545/$n=2$/-0.545,-0.242/$n=3$/-0.242,-0.136/$n=4$/-0.136,-0.087/$n=5$/-0.087,0/$n=\infty$/0}{
  \draw[eublue,thick] (.25,{1.32*\y})--(5.5,{1.32*\y});
  \node[right,font=\small] at (5.5,{1.32*\y}) {\lab};
  \node[left,font=\scriptsize] at (-.05,{1.32*\y}) {\val};
}
\draw[eured,very thick,-{Stealth}] (1.25,{1.32*(-.242)})--(1.25,{1.32*(-.545)});
\node[font=\scriptsize,anchor=west] at (1.35,-.55) {$3\rightarrow2$};
\draw[eugreen,very thick,-{Stealth}] (3.3,{1.32*(-2.18)})--(3.3,{1.32*(-.242)});
\node[font=\scriptsize,anchor=west] at (3.4,-1.55) {$1\rightarrow3$};
\node[rotate=90,font=\scriptsize] at (-1.15,-.9) {$E/(10^{-18}\,\mathrm J)$};
\end{tikzpicture}
\caption{Níveis de energia do átomo de hidrogénio. Os níveis aproximam-se progressivamente de $E=0$ à medida que $n$ aumenta.}
\end{figure}

\begin{leituragrafico}
O espaçamento entre níveis não é constante. Por exemplo, a diferença entre $n=1$ e $n=2$ é muito maior do que entre $n=4$ e $n=5$. Esta característica explica por que as riscas de uma mesma série espetral ficam progressivamente mais próximas.
\end{leituragrafico}

\subsection*{Séries espetrais do hidrogénio}
As transições do eletrão do hidrogénio podem agrupar-se em séries de acordo com o nível final. Historicamente, estas séries ajudaram a reconhecer a regularidade do espetro antes de existir uma explicação quântica completa.

\begin{table}[H]
\centering
\caption{Principais séries espetrais do átomo de hidrogénio.}
\begin{tabular}{lll}
\toprule
Série & Nível final & Região predominante do espetro \\
\midrule
Lyman & $n_f=1$ & ultravioleta \\
Balmer & $n_f=2$ & visível e limite próximo do UV \\
Paschen & $n_f=3$ & infravermelho \\
Brackett & $n_f=4$ & infravermelho \\
\bottomrule
\end{tabular}
\end{table}

\begin{exemplo}[title={Exemplo resolvido -- a risca $H_\alpha$ da série de Balmer}]
Considere a transição $n_i=3\rightarrow n_f=2$ no átomo de hidrogénio.
\[
E_3=-\frac{2{,}18\times10^{-18}}{9}=-2{,}42\times10^{-19}\ \mathrm J,
\qquad
E_2=-5{,}45\times10^{-19}\ \mathrm J.
\]
Logo,
\[
E_{\text{fotão}}=|E_2-E_3|\approx3{,}03\times10^{-19}\ \mathrm J.
\]
O comprimento de onda é
\[
\lambda=\frac{hc}{E_{\text{fotão}}}
\approx\frac{(6{,}626\times10^{-34})(3{,}00\times10^8)}{3{,}03\times10^{-19}}
\approx6{,}56\times10^{-7}\ \mathrm m=\SI{656}{nm}.
\]
Esta radiação situa-se na região vermelha do visível e pertence à série de Balmer.
\end{exemplo}

\subsection*{Energia de ionização do hidrogénio}
Ionizar o átomo significa levar o eletrão de um estado ligado até $n=\infty$, onde $E=0$. Assim, a energia mínima de ionização a partir do nível $n$ é
\[
E_{\text{ion}}=0-E_n=-E_n.
\]
A partir do estado fundamental, $E_{\text{ion}}=2{,}18\times10^{-18}\,\mathrm J$ por átomo.

\begin{nota}
O modelo de Bohr descreve com sucesso muitos aspetos do hidrogénio e de espécies com um único eletrão, mas não descreve adequadamente átomos polieletrónicos. No modelo quântico atual não se atribui ao eletrão uma trajetória circular definida.
\end{nota}

\section{Átomos polieletrónicos e energia de remoção}
Nos átomos com vários eletrões existem simultaneamente atração eletrão--núcleo e repulsões eletrão--eletrão. Os eletrões interiores reduzem parcialmente a atração nuclear sentida pelos eletrões mais externos, efeito que pode ser descrito qualitativamente como \textbf{blindagem}.

A \textbf{energia de remoção} é a energia mínima necessária para remover um eletrão. Valores experimentais de energia de remoção permitem inferir a existência de níveis e subníveis: eletrões que ocupam regiões energéticas diferentes não são igualmente fáceis de remover.

\begin{exemplo}[title=Exemplo resolvido -- leitura qualitativa de energias de remoção]
Se um conjunto de eletrões exige energias de remoção muito superiores às dos restantes, esses eletrões pertencem, em geral, a níveis mais internos e estão mais fortemente ligados ao núcleo. Um salto grande entre energias sucessivas é um indício de mudança de camada eletrónica.
\end{exemplo}


\subsection*{Espetroscopia fotoeletrónica e energias de remoção}
A \textbf{espetroscopia fotoeletrónica} (PES, do inglês \emph{photoelectron spectroscopy}) fornece informação experimental sobre as energias necessárias para remover eletrões de um átomo. Num espetro simplificado, cada conjunto de eletrões pertencente a um subnível pode originar um sinal. A posição do sinal relaciona-se com a energia de remoção e a sua intensidade relativa relaciona-se com o número de eletrões desse subnível.

\begin{figure}[H]
\centering
\begin{tikzpicture}
\begin{axis}[
width=11.5cm,height=5.9cm,
xlabel={energia de remoção (escala qualitativa)},
ylabel={intensidade relativa},
x dir=reverse,xmin=0,xmax=12,ymin=0,ymax=7.2,
ytick=\empty,xtick=\empty,axis lines=left,grid=none]
\addplot[ycomb,very thick,eublue,mark=*] coordinates {(10.6,2) (6.8,2) (4.3,6) (1.2,1)};
\node[font=\scriptsize] at (axis cs:10.6,2.6) {$1s^2$};
\node[font=\scriptsize] at (axis cs:6.8,2.6) {$2s^2$};
\node[font=\scriptsize] at (axis cs:4.3,6.6) {$2p^6$};
\node[font=\scriptsize] at (axis cs:1.2,1.6) {$3s^1$};
\end{axis}
\end{tikzpicture}
\caption{Espetro PES esquemático de um átomo com configuração $1s^2 2s^2 2p^6 3s^1$. O gráfico é qualitativo e não representa valores experimentais à escala.}
\end{figure}

\begin{conceito}
A menor energia de remoção corresponde, em geral, ao eletrão mais fácil de retirar. A \textbf{primeira energia de ionização} de um átomo gasoso no estado fundamental corresponde à energia necessária para remover esse eletrão mais fracamente ligado.
\end{conceito}

\begin{exemplo}[title={Exemplo resolvido -- interpretar um espetro PES qualitativo}]
Num espetro aparecem quatro grupos de sinais com intensidades relativas compatíveis com $2:2:6:1$, sendo o sinal de menor energia de remoção o menos intenso. Uma configuração coerente é
\[
1s^2\,2s^2\,2p^6\,3s^1.
\]
O átomo possui 11 eletrões, logo, se for neutro, tem $Z=11$ e corresponde ao sódio. O eletrão $3s$ é o mais fácil de remover.
\end{exemplo}

\section{Nuvem eletrónica e orbitais}
Em átomos, não se descrevem eletrões como pequenas esferas em órbitas planetárias rígidas. Utiliza-se um modelo probabilístico. Uma \textbf{orbital} corresponde a uma distribuição espacial de probabilidade associada a um eletrão.

As configurações eletrónicas obedecem, no nível estudado, a três ideias essenciais:
\begin{itemize}
  \item princípio da construção: ocupam-se primeiro orbitais de menor energia;
  \item princípio da exclusão de Pauli: cada orbital pode conter no máximo dois eletrões com estados de spin opostos;
  \item regra de Hund: em orbitais degeneradas, maximiza-se inicialmente o número de eletrões desemparelhados.
\end{itemize}

\begin{center}
\begin{tikzpicture}[scale=.9]
% s orbital
\fill[eublue!15] (0,0) circle (.8);
\draw[eublue] (0,0) circle (.8);
\node at (0,-1.25) {orbital $s$};
% p orbitals
\fill[eucyan!20] (3.1,.55) ellipse (.45 and .8);
\fill[eucyan!20] (3.1,-.55) ellipse (.45 and .8);
\draw[eucyan] (3.1,.55) ellipse (.45 and .8);
\draw[eucyan] (3.1,-.55) ellipse (.45 and .8);
\node at (3.1,-1.65) {uma orbital $p$};
% d schematic four lobes
\foreach \a in {45,135,225,315}{
  \begin{scope}[rotate around={\a:(6.4,0)}]
  \fill[euorange!20] (6.4,.65) ellipse (.32 and .7);
  \draw[euorange] (6.4,.65) ellipse (.32 and .7);
  \end{scope}}
\node at (6.4,-1.65) {orbital $d$ (esquema)};
\end{tikzpicture}
\end{center}

Num mesmo subnível, orbitais com a mesma energia dizem-se \textbf{degeneradas}. O subnível $p$ contém três orbitais e o subnível $d$ contém cinco.


\subsection*{Do modelo de Bohr ao modelo quântico}
A evolução dos modelos atómicos mostra como a ciência reformula representações quando surgem novos dados experimentais. O modelo quântico substitui a ideia de órbita por uma descrição probabilística da distribuição eletrónica.

\begin{figure}[H]
\centering
\begin{tikzpicture}[>=Stealth,node distance=.45cm,scale=.82,every node/.style={transform shape}]
\tikzset{model/.style={rounded corners,draw=eublue!70,fill=eulightblue,minimum width=1.95cm,minimum height=1.25cm,align=center,font=\small}}
\node[model] (d) {Dalton\\átomo indivisível};
\node[model,right=of d] (t) {Thomson\\eletrões no átomo};
\node[model,right=of t] (r) {Rutherford\\núcleo pequeno};
\node[model,right=of r] (b) {Bohr\\energia quantizada};
\node[model,right=of b] (q) {Modelo quântico\\orbitais};
\draw[->,thick,eucyan] (d)--(t);
\draw[->,thick,eucyan] (t)--(r);
\draw[->,thick,eucyan] (r)--(b);
\draw[->,thick,eucyan] (b)--(q);
\end{tikzpicture}
\caption{Evolução muito simplificada de alguns modelos atómicos. Cada modelo respondeu a evidências disponíveis no seu tempo.}
\end{figure}

No modelo quântico, uma orbital não é uma trajetória. É uma descrição da distribuição espacial associada ao eletrão. As orbitais $s$ apresentam simetria aproximadamente esférica; as orbitais $p$ apresentam dois lóbulos orientados segundo direções diferentes do espaço.

\begin{figure}[H]
\centering
\begin{tikzpicture}[scale=.95]
% s
\shade[ball color=eublue!35,opacity=.75] (0,0) circle (.9);
\fill[eured] (0,0) circle (2.3pt);
\node[font=\bfseries] at (0,-1.4) {$s$};
% px
\fill[eucyan!30,opacity=.9] (3.4,.62) ellipse (.42 and .82);
\fill[eucyan!55,opacity=.8] (3.4,-.62) ellipse (.42 and .82);
\fill[eured] (3.4,0) circle (2.2pt);
\node[font=\bfseries] at (3.4,-1.7) {$p_z$};
% py rotated
\begin{scope}[rotate around={90:(6.8,0)}]
\fill[eugreen!30,opacity=.9] (6.8,.62) ellipse (.42 and .82);
\fill[eugreen!55,opacity=.8] (6.8,-.62) ellipse (.42 and .82);
\end{scope}
\fill[eured] (6.8,0) circle (2.2pt);
\node[font=\bfseries] at (6.8,-1.7) {$p_x$};
% axes hint
\draw[->,gray!65] (8.7,-.65)--(9.7,-.65) node[right,font=\scriptsize] {$x$};
\draw[->,gray!65] (8.7,-.65)--(8.7,.35) node[above,font=\scriptsize] {$z$};
\end{tikzpicture}
\caption{Representações qualitativas de uma orbital $s$ e de duas orbitais $p$. As superfícies desenhadas delimitam regiões de elevada probabilidade, não ``paredes'' físicas.}
\end{figure}

\subsection*{Preenchimento das orbitais}
Para construir configurações eletrónicas até $Z=23$, é útil conjugar três regras:
\begin{enumerate}
  \item \textbf{Princípio de Aufbau ou da energia mínima:} ocupam-se primeiro as orbitais disponíveis de menor energia;
  \item \textbf{Princípio de Pauli:} cada orbital pode conter no máximo dois eletrões, representados com spins opostos;
  \item \textbf{Regra de Hund:} orbitais degeneradas são ocupadas primeiro por eletrões desemparelhados antes de ocorrer emparelhamento.
\end{enumerate}

\begin{figure}[H]
\centering
\begin{tikzpicture}[x=1.28cm,y=.72cm,>=Stealth]
\foreach \y/\name/\nboxes in {0/1s/1,1/2s/1,2/2p/3,3/3s/1,4/3p/3,5/4s/1,6/3d/5}{
  \node[anchor=east,font=\small\bfseries] at (-.15,\y) {\name};
  \foreach \i in {1,...,\nboxes}{\draw[eudark] ({\i-.85},\y-.28) rectangle ({\i-.25},\y+.28);}
}
\draw[->,euorange,very thick] (4.9,-.1)--(.15,1.9);
\draw[->,euorange,very thick] (4.9,1.3)--(.15,3.3);
\draw[->,euorange,very thick] (4.9,2.7)--(.15,4.7);
\draw[->,euorange,very thick] (4.9,4.1)--(.15,6.1);
\node[font=\scriptsize,align=center] at (6.2,3.0) {ordem de energia\\relevante até $Z=23$};
\end{tikzpicture}
\caption{Diagrama de caixas e lembrete gráfico da ordem de preenchimento das orbitais estudadas.}
\end{figure}

\begin{exemplo}[title=Exemplo resolvido -- configuração eletrónica]
Para o cloro, $Z=17$:
\[
\ce{Cl}: 1s^2\,2s^2\,2p^6\,3s^2\,3p^5.
\]
Possui sete eletrões de valência, pelo que tende a ganhar um eletrão em muitos contextos químicos, formando \ce{Cl-}.
\end{exemplo}

\section{Tabela Periódica: construção histórica e organização}
A Tabela Periódica é resultado de um processo histórico. Mendeleev organizou elementos de modo a evidenciar periodicidades químicas e deixou espaços para elementos ainda desconhecidos. Mais tarde, o trabalho de Moseley mostrou que a organização fundamental deve ser feita por \textbf{número atómico}, e não simplesmente por massa atómica.

A tabela atual organiza os elementos em \textbf{períodos} (linhas), \textbf{grupos} (colunas) e blocos $s$, $p$, $d$ e $f$, relacionados com o tipo de subnível que está a ser preenchido.

\begin{center}
\begin{tikzpicture}[x=.55cm,y=.55cm]
\draw[step=1,gray!50] (0,0) grid (18,4);
\fill[eublue!18] (0,0) rectangle (2,4);
\fill[eugreen!16] (12,0) rectangle (18,4);
\fill[euorange!18] (2,0) rectangle (12,2);
\node at (1,2) {bloco $s$};
\node at (15,2) {bloco $p$};
\node at (7,1) {bloco $d$};
\end{tikzpicture}
\end{center}


\subsection*{Grupos, períodos, blocos e famílias}
Os \textbf{grupos} são as colunas da Tabela Periódica (1 a 18) e os \textbf{períodos} são as linhas (1 a 7). Elementos do mesmo grupo apresentam frequentemente configurações eletrónicas de valência semelhantes e, por isso, propriedades químicas relacionadas.

\begin{figure}[H]
\centering
\begin{tikzpicture}[x=.5cm,y=.5cm]
\draw[step=1,gray!45] (0,0) grid (18,7);
\fill[eured!25] (0,0) rectangle (1,7);
\fill[euorange!25] (1,0) rectangle (2,7);
\fill[eugreen!25] (16,0) rectangle (17,7);
\fill[eublue!25] (17,0) rectangle (18,7);
\fill[eupurple!13] (2,0) rectangle (12,4);
\node[rotate=90,font=\scriptsize\bfseries] at (.5,3.5) {grupo 1 -- alcalinos};
\node[rotate=90,font=\scriptsize\bfseries] at (1.5,3.5) {grupo 2 -- alcalinoterrosos};
\node[rotate=90,font=\scriptsize\bfseries] at (16.5,3.5) {grupo 17 -- halogéneos};
\node[rotate=90,font=\scriptsize\bfseries] at (17.5,3.5) {grupo 18 -- gases nobres};
\node[font=\scriptsize,align=center] at (7,2) {metais de transição\\(bloco $d$, esquema)};
\end{tikzpicture}
\caption{Algumas famílias e zonas da Tabela Periódica. O desenho é esquemático e não pretende reproduzir as posições especiais do H, He e bloco $f$.}
\end{figure}

\begin{table}[H]
\centering
\caption{Configuração eletrónica de valência de alguns grupos representativos.}
\begin{tabular}{lll}
\toprule
Grupo & Configuração de valência típica & Consequência frequente \\
\midrule
1 & $ns^1$ & tendência para formar catiões $+1$ \\
2 & $ns^2$ & tendência para formar catiões $+2$ \\
16 & $ns^2np^4$ & seis eletrões de valência \\
17 & $ns^2np^5$ & tendência para formar aniões $-1$ \\
18 & $ns^2np^6$ (exceto He) & elevada estabilidade em muitas condições \\
\bottomrule
\end{tabular}
\end{table}

A configuração eletrónica permite localizar muitos elementos representativos:
\begin{itemize}
  \item o maior valor de $n$ indica o \textbf{período};
  \item o subnível que está a ser preenchido indica o \textbf{bloco};
  \item o número de eletrões de valência ajuda a identificar o \textbf{grupo} nos blocos $s$ e $p$.
\end{itemize}

\begin{figure}[H]
\centering
\begin{tikzpicture}[>=Stealth,node distance=7mm]
\node[rounded corners,draw=eublue,fill=eulightblue,minimum width=5.4cm,minimum height=1cm] (c) {$1s^2\,2s^2\,2p^6\,3s^2\,3p^4$};
\node[rounded corners,draw=eucyan,fill=eulightcyan,below left=of c,align=center] (p) {maior $n=3$\\$\Rightarrow$ período 3};
\node[rounded corners,draw=euorange,fill=eulightorange,below=of c,align=center] (b) {termina em $p$\\$\Rightarrow$ bloco $p$};
\node[rounded corners,draw=eugreen,fill=eulightgreen,below right=of c,align=center] (g) {$3s^23p^4$: 6 e$^-$ de valência\\$\Rightarrow$ grupo 16};
\draw[->,thick] (c)--(p); \draw[->,thick] (c)--(b); \draw[->,thick] (c)--(g);
\end{tikzpicture}
\caption{Exemplo de localização de um elemento a partir da configuração eletrónica: o enxofre.}
\end{figure}

\begin{exemplo}[title={Exemplo resolvido -- localizar um elemento na Tabela Periódica}]
A configuração $1s^2 2s^2 2p^6 3s^2 3p^4$ contém 16 eletrões, logo corresponde ao enxofre neutro, $Z=16$. O maior valor de $n$ é 3, portanto pertence ao 3.º período. Termina em $p$, logo pertence ao bloco $p$. Possui seis eletrões de valência, pelo que se encontra no grupo 16.
\end{exemplo}

\section{Tendências periódicas}
A Tabela Periódica organiza elementos por número atómico crescente. As tendências não são regras mecânicas: resultam da competição entre carga nuclear, distância ao núcleo, blindagem eletrónica e configuração eletrónica.

Três ideias ajudam a justificar muitas tendências periódicas:
\begin{table}[H]
\centering
\caption{Fatores qualitativos que influenciam a atração núcleo--eletrão.}
\begin{tabularx}{\textwidth}{lXX}
\toprule
Fator & Significado & Efeito qualitativo sobre os eletrões de valência \\
\midrule
Distância ao núcleo & camadas com maior $n$ ficam, em média, mais afastadas & maior distância tende a enfraquecer a atração \\
Carga nuclear efetiva & atração líquida exercida pelo núcleo depois de considerar a presença dos outros eletrões & maior carga efetiva tende a reforçar a atração \\
Blindagem & redução da atração sentida pelos eletrões externos devido aos eletrões interiores & maior blindagem tende a enfraquecer a atração \\
\bottomrule
\end{tabularx}
\end{table}

\begin{figure}[H]
\centering
\begin{tikzpicture}[x=1cm,y=1cm,>=Stealth]
\draw[rounded corners,fill=eugrey!60,draw=gray!55] (0,0) rectangle (9,4.5);
\foreach \x in {0.5,1,...,8.5} \draw[gray!25] (\x,.35)--(\x,4.15);
\foreach \y in {.75,1.75,2.75,3.75} \draw[gray!25] (.35,\y)--(8.65,\y);
\draw[->,very thick,eured] (1.0,1.0)--(7.8,1.0) node[midway,below=3pt,black]{energia de ionização tende a aumentar};
\draw[->,very thick,eured] (7.8,1.25)--(7.8,3.8);
\draw[->,very thick,eublue] (7.7,3.5)--(1.0,3.5) node[midway,above=3pt,black]{raio atómico tende a aumentar};
\draw[->,very thick,eublue] (1.0,3.3)--(1.0,.8);
\end{tikzpicture}
\caption{Tendências periódicas qualitativas mais importantes no 10.º ano.}
\end{figure} Elementos do mesmo grupo apresentam semelhanças de configuração eletrónica de valência e, por isso, propriedades químicas relacionadas.

De forma qualitativa, para elementos representativos:
\begin{itemize}
  \item o raio atómico tende a diminuir ao longo de um período, da esquerda para a direita;
  \item o raio atómico tende a aumentar ao descer num grupo;
  \item a primeira energia de ionização tende a aumentar ao longo de um período e a diminuir ao descer num grupo, com exceções explicáveis pela estrutura eletrónica.
\end{itemize}

\begin{table}[H]
\centering
\caption{Relação entre posição, eletrões de valência e tendência de formação de iões.}
\begin{tabular}{ccc}
\toprule
Grupo (simplificado) & Eletrões de valência & Tendência comum \\
\midrule
1 & 1 & formar catiões $+1$ \\
2 & 2 & formar catiões $+2$ \\
17 & 7 & formar aniões $-1$ \\
18 & camada externa completa & baixa reatividade \\
\bottomrule
\end{tabular}
\end{table}

Os \textbf{gases nobres} apresentam baixa reatividade em muitas condições porque a sua camada de valência está completa. Esta estabilidade não significa que sejam absolutamente incapazes de formar compostos, mas explica a sua reduzida tendência para formar iões ou ligações nas situações mais comuns estudadas neste nível.


\subsection*{Raio de iões e séries isoeletrónicas}
Ao formar um catião, um átomo perde eletrões. A nuvem eletrónica tende a contrair-se porque diminuem as repulsões eletrão--eletrão e, em muitos casos, desaparece a camada de valência mais externa. Assim, um catião é geralmente \textbf{menor} do que o átomo que lhe deu origem. Um anião, pelo contrário, recebe eletrões e é geralmente \textbf{maior} do que o átomo neutro correspondente.

\begin{figure}[H]
\centering
\begin{tikzpicture}[scale=.95]
% cation
\fill[eublue!16] (0,0) circle (1.15); \draw[eublue,thick] (0,0) circle (1.15);
\fill[euorange!28] (3.0,0) circle (.76); \draw[euorange,thick] (3.0,0) circle (.76);
\node at (0,-1.55) {átomo}; \node at (3,-1.55) {catião};
\draw[->,very thick,eudark] (1.35,0)--(2.1,0) node[midway,above,font=\scriptsize] {$-e^-$};
% anion
\fill[eublue!16] (6.1,0) circle (1.0); \draw[eublue,thick] (6.1,0) circle (1.0);
\fill[eugreen!18] (9.2,0) circle (1.35); \draw[eugreen,thick] (9.2,0) circle (1.35);
\node at (6.1,-1.55) {átomo}; \node at (9.2,-1.55) {anião};
\draw[->,very thick,eudark] (7.3,0)--(8.0,0) node[midway,above,font=\scriptsize] {$+e^-$};
\end{tikzpicture}
\caption{Tendência geral da dimensão ao formar catiões e aniões. O desenho não está à escala.}
\end{figure}

Em espécies \textbf{isoeletrónicas}, isto é, com o mesmo número de eletrões, a dimensão tende a diminuir quando aumenta o número de protões. Por exemplo, na série $\ce{O^{2-}}$, $\ce{F^-}$, $\ce{Na+}$, $\ce{Mg^{2+}}$ e $\ce{Al^{3+}}$, todas as espécies têm 10 eletrões; o aumento de $Z$ reforça a atração sobre a mesma nuvem eletrónica.

\begin{exemplo}[title={Exemplo resolvido -- comparar raios de espécies isoeletrónicas}]
Considere $\ce{O^{2-}}$ e $\ce{Mg^{2+}}$. Ambas possuem 10 eletrões. O magnésio tem 12 protões, enquanto o oxigénio tem 8. A nuvem de $\ce{Mg^{2+}}$ é atraída por uma carga nuclear maior, pelo que $\ce{Mg^{2+}}$ apresenta raio menor que $\ce{O^{2-}}$.
\end{exemplo}

\begin{atividade}[title=AL 1.2 -- Teste de chama]
\textbf{Objetivo:} identificar elementos metálicos em amostras de sais a partir da cor da chama e, quando disponível, do respetivo espetro de emissão.

\textbf{Fundamento:} o aquecimento excita eletrões. No regresso a estados de menor energia, são emitidos fotões cujas energias dependem da estrutura eletrónica do elemento.

\textbf{Controlo experimental:} usar material limpo para evitar contaminação, observar amostras conhecidas antes das desconhecidas e manter as condições de observação o mais constantes possível.

\textbf{Segurança:} óculos de proteção, pequenas quantidades, cabelo preso e manipulação de fontes de chama apenas segundo as instruções do professor.
\end{atividade}

\begin{atividade}[title=AL 1.3 -- Densidade relativa de metais por picnometria]
A densidade relativa $d$ de um metal pode ser determinada comparando a massa do metal com a massa de igual volume de água. Num procedimento comum, mede-se:
\begin{itemize}
  \item $m$: massa do metal;
  \item $M$: massa do picnómetro cheio de água mais o metal colocado ao lado;
  \item $M_3$: massa do picnómetro com o metal no interior e completado com água.
\end{itemize}
A massa de água deslocada pelo metal é $M-M_3$, pelo que, numa aproximação adequada,
\[
d\simeq\frac{m}{M-M_3}.
\]
Devem controlar-se bolhas de ar, secagem exterior do picnómetro, temperatura e pureza do metal. O resultado pode ser comparado com um valor de referência por meio do erro relativo percentual.
\end{atividade}

\begin{exemplo}[title=Exemplo resolvido -- picnometria]
Se $m=\SI{8.70}{g}$, $M=\SI{107.56}{g}$ e $M_3=\SI{102.53}{g}$,
\[
d=\frac{8{,}70}{107{,}56-102{,}53}\approx1{,}73.
\]
Um valor desta ordem é compatível, por exemplo, com um metal pouco denso como o magnésio.
\end{exemplo}

\begin{atividade}[title=Extensão -- Teste de chama]
\textbf{Teste de chama.} Sais de alguns metais produzem cores características quando aquecidos. A observação qualitativa pode ser interpretada em termos de excitação eletrónica e emissão de fotões. Devem ser usadas pequenas quantidades e procedimentos seguros definidos pelo professor.
\end{atividade}

\begin{mundo}
A espectroscopia permite identificar elementos presentes em estrelas sem recolher amostras. A mesma ideia de ``assinatura espectral'' é usada em química forense, controlo ambiental, análise de materiais e astronomia.
\end{mundo}

\begin{exercicios}
\begin{enumerate}
  \item Calcula a frequência de uma radiação com $\lambda=\SI{600}{nm}$.
  \item Determina a energia de um fotão de frequência $5{,}0\times10^{14}\,\mathrm{Hz}$.
  \item Explica por que razão um espectro de emissão apresenta linhas em vez de um contínuo.
  \item Escreve as configurações eletrónicas de O, Na, Cl, Ca e V.
  \item Compara qualitativamente os raios atómicos de Na e Cl e justifica.
  \item Qual dos elementos, Mg ou Ca, deverá ter maior primeira energia de ionização? Justifica com base na posição na Tabela Periódica.
  \item \textbf{Integração.} Uma lâmpada de descarga emite luz de \SI{589}{nm}. Calcula a energia de um fotão e de uma mole desses fotões.
  \item Ordena por energia crescente fotões de \SI{900}{nm}, \SI{600}{nm} e \SI{300}{nm}.
  \item Explica por que razão os eletrões de camadas internas apresentam, em geral, maior energia de remoção do que eletrões de valência.
  \item Num ensaio de picnometria obtêm-se $m=\SI{5.00}{g}$, $M=\SI{80.25}{g}$ e $M_3=\SI{78.40}{g}$. Estima a densidade relativa do metal.
  \item Para radiação de comprimento de onda \SI{555}{nm}, determina (a) a frequência; (b) o período; (c) a energia de um fotão.
  \item Uma radiação vermelha tem $\lambda=\SI{700}{nm}$ e uma violeta $\lambda=\SI{400}{nm}$. Calcula a razão $E_{\rm violeta}/E_{\rm vermelha}$.
  \item Explica por que as riscas de absorção de um elemento surgem nos mesmos comprimentos de onda que algumas das suas riscas de emissão.
  \item Ordena vermelho, verde e violeta por ordem crescente de frequência e por ordem crescente de energia por fotão.
  \item No átomo de hidrogénio, calcula $E_2$ e $E_4$. Determina a energia do fotão emitido na transição $4\rightarrow2$.
  \item Uma transição do hidrogénio termina em $n=1$. A que série espetral pertence? Em que região do espetro se encontra predominantemente?
  \item Calcula a energia mínima necessária para ionizar um átomo de hidrogénio inicialmente no nível $n=2$.
  \item Um espetro PES qualitativo apresenta sinais com intensidades relativas $2:2:6:2:5$. Propõe uma configuração eletrónica coerente e identifica o elemento neutro.
  \item Localiza na Tabela Periódica um elemento cuja configuração termina em $3s^23p^3$: indica período, bloco e grupo.
  \item Ordena $\ce{O^{2-}}$, $\ce{F^-}$, $\ce{Na+}$ e $\ce{Mg^{2+}}$ por raio crescente e justifica sabendo que são espécies isoeletrónicas.
\end{enumerate}
\end{exercicios}

\chapter{Ligação química, geometria molecular e forças intermoleculares}

\begin{objetivos}
\begin{itemize}
  \item reconhecer átomos, moléculas e iões como unidades estruturais de diferentes substâncias;
  \item interpretar uma curva de energia potencial em função da distância internuclear;
  \item distinguir ligações covalente, iónica e metálica;
  \item representar eletrões de valência e usar estruturas de Lewis em moléculas simples;
  \item relacionar ordem de ligação, comprimento de ligação e energia de ligação;
  \item prever geometrias moleculares simples e polaridade;
  \item distinguir hidrocarbonetos saturados e insaturados;
  \item reconhecer grupos funcionais básicos;
  \item distinguir ligações intramoleculares de interações intermoleculares e relacioná-las com propriedades macroscópicas.
\end{itemize}
\end{objetivos}


\section{Unidades estruturais da matéria}
As substâncias podem ser descritas, a uma escala microscópica, por diferentes unidades estruturais. Em substâncias moleculares, a unidade estrutural é a \textbf{molécula}; em muitos cristais iónicos consideram-se \textbf{iões}; nos metais e em algumas redes covalentes é útil considerar \textbf{átomos} organizados numa estrutura extensa.

\begin{figure}[H]
\centering
\begin{tikzpicture}[scale=1]
\draw[rounded corners,fill=eulightblue,draw=eublue] (-.5,-1.45) rectangle (3.0,1.55);
\node[font=\bfseries\color{eublue}] at (1.25,1.18) {Moléculas};
\fill[eured!85] (.8,0) circle (.42); \fill[white,draw=eudark!50] (.15,.35) circle (.24); \fill[white,draw=eudark!50] (1.45,.35) circle (.24);
\draw[thick] (.35,.25)--(.55,.12); \draw[thick] (1.05,.12)--(1.25,.25);
\node at (2.25,0) {$\ce{H2O}$};

\draw[rounded corners,fill=eulightgreen,draw=eugreen] (3.5,-1.45) rectangle (7.0,1.55);
\node[font=\bfseries\color{eugreen}] at (5.25,1.18) {Iões};
\fill[euorange!45] (4.55,0) circle (.40); \node at (4.55,0) {$+$};
\fill[eugreen!35] (5.75,0) circle (.52); \node at (5.75,0) {$-$};
\node at (6.35,-.75) {$\ce{Na+}$, $\ce{Cl-}$};

\draw[rounded corners,fill=eulightorange,draw=euorange] (7.5,-1.45) rectangle (11.0,1.55);
\node[font=\bfseries\color{euorange}] at (9.25,1.18) {Átomos};
\foreach \x/\y in {8.2/.2,8.9/.2,9.6/.2,10.3/.2,8.55/-.45,9.25/-.45,9.95/-.45}{\fill[euorange!55] (\x,\y) circle (.22);}
\node at (9.25,-.95) {ex.: metal};
\end{tikzpicture}
\caption{Três formas de representar unidades estruturais de substâncias.}
\end{figure}

Uma \textbf{ligação química} mantém átomos ou iões associados numa substância. As interações entre moléculas, embora muito importantes para as propriedades macroscópicas, não devem ser confundidas com as ligações que mantêm unidos os átomos dentro de cada molécula.

\section{Porque se formam ligações?}
A formação de uma ligação é favorecida quando conduz a um sistema de menor energia e maior estabilidade. Num modelo simples de dois átomos, existe uma distância internuclear de equilíbrio para a qual a energia potencial é mínima.

\begin{figure}[H]
\centering
\begin{tikzpicture}
\begin{axis}[
width=11.5cm,height=6.5cm,
xlabel={distância internuclear},ylabel={energia potencial},
axis lines=left,xtick=\empty,ytick=\empty,ymin=-3.7,ymax=2.2,xmin=.55,xmax=5.2,grid=none]
\path[fill=eured!8] (axis cs:.55,-3.7) rectangle (axis cs:.92,2.2);
\path[fill=eugreen!10] (axis cs:.92,-3.7) rectangle (axis cs:1.35,2.2);
\path[fill=euyellow!65] (axis cs:1.35,-3.7) rectangle (axis cs:2.85,2.2);
\path[fill=euorange!8] (axis cs:2.85,-3.7) rectangle (axis cs:5.2,2.2);
\addplot[domain=.68:5.05,samples=250,very thick,eublue]{4*((1/x)^12-(1/x)^6)};
\draw[dashed,eured] (axis cs:1.122,-3.2)--(axis cs:1.122,0);
\node[font=\scriptsize,anchor=west] at (axis cs:1.18,-3.05) {distância de equilíbrio};
\node[font=\scriptsize,align=center] at (axis cs:.74,1.55) {repulsão\\dominante};
\node[font=\scriptsize,align=center] at (axis cs:1.13,-2.25) {máxima\\estabilidade};
\node[font=\scriptsize,align=center] at (axis cs:2.0,.75) {atração\\dominante};
\node[font=\scriptsize,align=center] at (axis cs:4.0,.65) {átomos muito\\afastados};
\end{axis}
\end{tikzpicture}
\caption{Curva de estabilidade de uma molécula diatómica. A ligação corresponde ao mínimo de energia potencial.}
\end{figure}

\begin{leituragrafico}
Quando os átomos estão muito afastados, a interação é pequena e toma-se frequentemente $E_p\approx0$. Ao aproximarem-se, predominam inicialmente atrações e a energia diminui. No mínimo existe uma distância de equilíbrio. A distâncias ainda menores, as repulsões núcleo--núcleo e eletrão--eletrão tornam-se muito intensas e a energia aumenta rapidamente.
\end{leituragrafico}

\subsection*{Energia de ligação, comprimento e ordem de ligação}
Quando dois átomos se aproximam, atrações e repulsões variam com a distância internuclear. A distância de equilíbrio corresponde ao mínimo da energia potencial. A profundidade do ``poço'' energético relaciona-se com a energia necessária para separar os átomos.

\begin{figure}[H]
\centering
\begin{tikzpicture}
\begin{axis}[width=11cm,height=6.3cm,xlabel={distância internuclear / unidade arbitrária},ylabel={energia potencial / unidade arbitrária},xmin=.55,xmax=4.5,ymin=-3.4,ymax=5.5,axis lines=left,grid=major]
\addplot[domain=.72:4.4,samples=180,thick,eublue]{4*((1/x)^12-(1/x)^6)};
\draw[dashed,euorange] (axis cs:1.122,-3.2)--(axis cs:1.122,0);
\node[anchor=west,font=\small] at (axis cs:1.18,-2.65) {distância de equilíbrio};
\draw[<->,eugreen,thick] (axis cs:2.9,0)--(axis cs:2.9,-3.0);
\node[anchor=west,font=\small] at (axis cs:2.95,-1.5) {energia de ligação};
\end{axis}
\end{tikzpicture}
\caption{Curva qualitativa de energia potencial para dois átomos ligados.}
\end{figure}
O mínimo da curva de energia potencial corresponde à distância internuclear de equilíbrio. A profundidade do mínimo relaciona-se com a energia necessária para separar os átomos. Para ligações entre os mesmos elementos, uma ligação múltipla é, em geral, \textbf{mais curta e mais forte} do que uma ligação simples.

\begin{table}[H]
\centering
\caption{Tendência qualitativa para ligações C--C.}
\begin{tabular}{ccc}
\toprule
Ligação & Comprimento relativo & Energia de ligação relativa \\
\midrule
C--C & maior & menor \\
C=C & intermédio & intermédia \\
C$\equiv$C & menor & maior \\
\bottomrule
\end{tabular}
\end{table}


\subsection*{Eletrões de valência e notação de Lewis}
Os eletrões de valência participam diretamente na formação de muitas ligações químicas. A \textbf{notação de Lewis} representa o símbolo químico e os eletrões de valência por pontos colocados à sua volta. Esta representação ajuda a visualizar a regra do octeto, que é útil para muitas espécies dos elementos representativos, embora tenha exceções.

\begin{table}[H]
\centering
\caption{Exemplos de notação de Lewis para átomos.}
\begin{tabular}{ccc}
\toprule
Átomo & Eletrões de valência & Representação esquemática \\
\midrule
H & 1 & $\mathrm{H}\!\boldsymbol{\cdot}$ \\
Mg & 2 & $\boldsymbol{\cdot}\!\mathrm{Mg}\!\boldsymbol{\cdot}$ \\
N & 5 & $\begin{array}{c}\boldsymbol{\cdot\cdot}\\[-2mm]\boldsymbol{\cdot}\,\mathrm{N}\,\boldsymbol{\cdot}\\[-2mm]\boldsymbol{\cdot}\end{array}$ \\
Ne & 8 & $\begin{array}{c}\boldsymbol{\cdot\cdot}\\[-2mm]\boldsymbol{\cdot\cdot}\,\mathrm{Ne}\,\boldsymbol{\cdot\cdot}\\[-2mm]\boldsymbol{\cdot\cdot}\end{array}$ \\
\bottomrule
\end{tabular}
\end{table}

\begin{nota}
A regra do octeto é um modelo simples, não uma lei universal. No 10.º ano é particularmente útil para compreender muitas moléculas e iões formados por elementos representativos.
\end{nota}

\section{Três modelos de ligação}
\begin{description}
  \item[Ligação covalente:] envolve partilha de pares eletrónicos entre átomos, típica de muitas substâncias moleculares.
  \item[Ligação iónica:] resulta da atração eletrostática entre iões de carga oposta numa rede cristalina.
  \item[Ligação metálica:] descreve catiões metálicos num retículo associados a eletrões deslocalizados, ajudando a explicar condutividade e maleabilidade.
\end{description}

\begin{figure}[H]
\centering
\begin{tikzpicture}[scale=.92]
% covalent
\node[font=\bfseries\color{eublue}] at (1.3,2.2) {covalente};
\fill[eublue!22] (.6,1.0) circle (.65); \fill[eublue!22] (2.0,1.0) circle (.65);
\draw[eublue,thick] (.6,1.0) circle (.65); \draw[eublue,thick] (2.0,1.0) circle (.65);
\fill[eured] (1.18,1.08) circle (.07); \fill[eured] (1.42,.92) circle (.07);
\node[font=\scriptsize,align=center] at (1.3,-.05) {par eletrónico\\partilhado};
% ionic
\node[font=\bfseries\color{euorange}] at (5.1,2.2) {iónica};
\fill[euorange!25] (4.45,1.0) circle (.55); \fill[eugreen!25] (5.75,1.0) circle (.7);
\node at (4.45,1.0) {$+$}; \node at (5.75,1.0) {$-$};
\draw[<->,thick,eudark] (4.98,1.0)--(5.08,1.0);
\node[font=\scriptsize,align=center] at (5.1,-.05) {atração eletrostática\\entre iões};
% metallic
\node[font=\bfseries\color{eupurple}] at (9.2,2.2) {metálica};
\foreach \x/\y in {8.4/.7,9.2/.7,10/.7,8.8/1.4,9.6/1.4}{\fill[eupurple!30] (\x,\y) circle (.28); \node[font=\scriptsize] at (\x,\y) {$+$};}
\foreach \x/\y in {8.35,1.7,8.8,.3,9.3,1.05,9.75,.25,10.15,1.65}{\fill[eublue] (\x,\y) circle (.055);}
\node[font=\scriptsize,align=center] at (9.2,-.05) {eletrões de valência\\deslocalizados};
\end{tikzpicture}
\caption{Representações muito simplificadas dos três modelos de ligação estudados.}
\end{figure}

\section{Estruturas de Lewis e geometria}
Estruturas de Lewis representam eletrões de valência. Um procedimento útil é:
\begin{enumerate}
  \item contar o número total de eletrões de valência;
  \item escolher o átomo central, quando existe;
  \item ligar os átomos inicialmente por ligações simples;
  \item completar duetos de H e octetos dos átomos periféricos;
  \item colocar os eletrões restantes no átomo central;
  \item se o átomo central não completar o octeto, considerar ligações múltiplas.
\end{enumerate}

\begin{exemplo}[title=Exemplo resolvido -- estrutura de Lewis do dióxido de carbono]
O carbono contribui com 4 eletrões de valência e cada oxigénio com 6: no total, 16. A estrutura que completa os octetos é
\[
\ce{O=C=O}.
\]
O átomo central tem duas regiões eletrónicas de ligação, pelo que a molécula é linear.
\end{exemplo}

Estruturas de Lewis representam eletrões de valência. A geometria pode ser prevista qualitativamente considerando a repulsão entre regiões eletrónicas em torno do átomo central.

\begin{center}
\begin{tikzpicture}[scale=1.05]
% CO2
\node (C) at (0,0) {C};
\node (O1) at (-1.5,0) {O};
\node (O2) at (1.5,0) {O};
\draw[double] (O1)--(C); \draw[double] (C)--(O2);
\node at (0,-.75) {\ce{CO2}: linear, globalmente apolar};
% H2O
\node (Ow) at (6,0) {O};
\node (H1) at (5,-1) {H};
\node (H2) at (7,-1) {H};
\draw (Ow)--(H1); \draw (Ow)--(H2);
\node at (6,-1.75) {\ce{H2O}: angular, polar};
\end{tikzpicture}
\end{center}

A polaridade global depende da polaridade das ligações e da geometria. Ligações polares podem cancelar-se vetorialmente numa molécula simétrica.

\begin{table}[H]
\centering
\caption{Geometrias moleculares simples a reconhecer.}
\begin{tabular}{llll}
\toprule
Molécula & Regiões em torno do átomo central & Geometria & Polaridade global \\
\midrule
\ce{CO2} & 2 & linear & apolar \\
\ce{CH4} & 4 & tetraédrica & apolar \\
\ce{NH3} & 4 (1 par não ligante) & piramidal trigonal & polar \\
\ce{H2O} & 4 (2 pares não ligantes) & angular & polar \\
\bottomrule
\end{tabular}
\end{table}

\begin{exemplo}[title=Exemplo resolvido -- polaridade molecular]
Em \ce{CO2}, cada ligação C--O é polar. Contudo, a molécula é linear e os dois vetores de dipolo têm igual módulo e sentidos opostos: anulam-se. Em \ce{H2O}, a geometria angular impede essa anulação, pelo que a molécula é polar.
\end{exemplo}

\begin{figure}[H]
\centering
\begin{tikzpicture}[scale=.95]
\node[font=\bfseries] at (0,2.4) {linear};
\fill[eublue] (0,1.2) circle (5pt); \fill[euorange] (-1.2,1.2) circle (4pt); \fill[euorange] (1.2,1.2) circle (4pt);
\draw[thick] (-1.15,1.2)--(-.1,1.2); \draw[thick] (.1,1.2)--(1.15,1.2);
\node[font=\small] at (0,.45) {$180^\circ$};
\node[font=\bfseries] at (4,2.4) {angular};
\fill[eublue] (4,1.2) circle (5pt); \fill[euorange] (3.0,.55) circle (4pt); \fill[euorange] (5.0,.55) circle (4pt);
\draw[thick] (3.1,.62)--(3.92,1.15); \draw[thick] (4.08,1.15)--(4.9,.62);
\node[font=\small] at (4,.25) {ex.: $\ce{H2O}$};
\node[font=\bfseries] at (8,2.4) {tetraédrica};
\fill[eublue] (8,1.2) circle (5pt);
\foreach \a in {30,150,230,310}{\draw[thick] (8,1.2)--++(\a:1.0); \fill[euorange] ($(8,1.2)+(\a:1.0)$) circle (4pt);}
\node[font=\small] at (8,-.05) {ex.: $\ce{CH4}$};
\end{tikzpicture}
\caption{Geometrias moleculares frequentes representadas de modo esquemático.}
\end{figure}

\begin{conceito}
A geometria molecular depende da organização tridimensional das regiões de densidade eletrónica em torno do átomo central. A polaridade da molécula depende simultaneamente da polaridade das ligações e da geometria: ligações polares podem originar uma molécula globalmente apolar se os dipolos se compensarem por simetria.
\end{conceito}

\section{Hidrocarbonetos e grupos funcionais}
Hidrocarbonetos contêm apenas carbono e hidrogénio. Alcanos são saturados; alcenos e alcinos possuem ligações múltiplas e são insaturados.

\begin{table}[H]
\centering
\caption{Alguns grupos funcionais a reconhecer.}
\begin{tabular}{lll}
\toprule
Família & Grupo característico & Exemplo \\
\midrule
Álcool & $-\mathrm{OH}$ & etanol \\
Aldeído & $-\mathrm{CHO}$ & etanal \\
Cetona & $>\mathrm{C{=}O}$ & propanona \\
Ácido carboxílico & $-\mathrm{COOH}$ & ácido etanoico \\
Amina & $-\mathrm{NH_2}$ & metilamina \\
\bottomrule
\end{tabular}
\end{table}

\section{Interações intermoleculares}
Entre moléculas podem existir forças de dispersão de London, interações dipolo--dipolo e, em situações apropriadas, pontes de hidrogénio. Estas interações ajudam a explicar pontos de ebulição, viscosidade, tensão superficial e miscibilidade.

\begin{exemplo}[title=Exemplo resolvido -- forças intermoleculares]
Etano e etanol têm massas molares da mesma ordem de grandeza, mas o etanol possui um grupo O--H e estabelece pontes de hidrogénio. Por isso, as atrações intermoleculares são mais intensas e o seu ponto de ebulição é muito superior ao do etano.
\end{exemplo}

\begin{atividade}[title=AL 2.1 -- Miscibilidade de líquidos]
Selecionam-se pares de líquidos, por exemplo água, etanol e um hidrocarboneto adequado, e formulam-se previsões com base na polaridade molecular. Misturam-se pequenos volumes, agita-se e observa-se se se forma uma única fase ou duas fases. A conclusão deve relacionar estrutura molecular, polaridade e intensidade/tipo de interações intermoleculares.

\textbf{Segurança:} usar apenas substâncias autorizadas pelo professor, evitar fontes de ignição quando existirem líquidos inflamáveis e encaminhar resíduos segundo as regras do laboratório.
\end{atividade}

\begin{mundo}
A água tem propriedades invulgares devido à sua polaridade e às pontes de hidrogénio. Isso influencia a regulação térmica dos seres vivos, a estrutura de biomoléculas e o transporte de substâncias em ecossistemas.
\end{mundo}

\begin{atividade}[title=Extensão investigativa -- miscibilidade]
\textbf{Miscibilidade.} Compare água, etanol e um óleo alimentar em pequenos volumes. Registe que pares formam uma fase e quais formam duas fases. Interprete os resultados com base na polaridade e nas interações intermoleculares. Não ingerir qualquer material usado em laboratório.
\end{atividade}

\begin{exercicios}
\begin{enumerate}
  \item Classifica, de forma predominante, a ligação em \ce{NaCl}, \ce{O2} e Cu(s).
  \item Desenha estruturas de Lewis para \ce{H2O}, \ce{NH3}, \ce{CO2} e \ce{CH4}.
  \item Indica a geometria de \ce{CO2}, \ce{CH4}, \ce{NH3} e \ce{H2O}.
  \item Explica por que razão \ce{CO2} é globalmente apolar, embora cada ligação C--O seja polar.
  \item Distingue etano, eteno e etino quanto à saturação.
  \item Identifica o grupo funcional principal em etanol, ácido etanoico e propanona.
  \item Prevê qual deverá apresentar maior ponto de ebulição, etano ou etanol, e justifica qualitativamente.
  \item Compara qualitativamente o comprimento e a energia das ligações C--C, C=C e C$\equiv$C.
  \item Indica se \ce{BF3} (trigonal planar e simétrica) deverá ser polar ou apolar, sabendo que as ligações B--F são polares.
  \item Entre \ce{CH4}, \ce{NH3} e \ce{H2O}, identifica as moléculas capazes de estabelecer pontes de hidrogénio entre si.
  \item Numa curva $E_p(r)$, explica o significado físico (a) da região onde $E_p\to0$; (b) do mínimo; (c) do crescimento muito rápido de $E_p$ a distâncias pequenas.
  \item Representa a notação de Lewis dos átomos H, O, Na e Cl e indica quantos eletrões de valência possui cada um.
  \item Explica, usando um modelo de partículas, por que um cristal iónico sólido não deve ser representado como um conjunto de moléculas isoladas de \ce{NaCl}.
  \item A ligação H--H tem comprimento de equilíbrio cerca de \SI{74}{pm}. Explica o que aconteceria energeticamente se os dois núcleos fossem aproximados para uma distância muito inferior.
\end{enumerate}
\end{exercicios}

\chapter{Gases, atmosfera, dispersões e soluções}

\begin{objetivos}
\begin{itemize}
  \item compreender a lei de Avogadro e o conceito de volume molar;
  \item relacionar massa, massa molar, volume e massa volúmica de gases;
  \item descrever a composição geral da troposfera e discutir poluentes;
  \item efetuar cálculos de concentração em soluções;
  \item preparar soluções por dissolução e por diluição.
\end{itemize}
\end{objetivos}

\section{Lei de Avogadro e volume molar}
À mesma temperatura e pressão, volumes iguais de gases ideais contêm o mesmo número de moléculas. O volume molar $V_m$ depende das condições de temperatura e pressão:
\[
V_m=\frac{V}{n}.
\]
Não se deve usar um único valor de $V_m$ sem indicar as condições a que se refere.

\begin{exemplo}[title=Exemplo resolvido -- volume molar]
Uma amostra de gás ocupa \SI{6.00}{L} em condições nas quais $V_m=\SI{24.0}{L/mol}$. Então
\[
n=\frac{V}{V_m}=\frac{6{,}00}{24{,}0}=\SI{0.250}{mol}.
\]
Se a temperatura ou a pressão mudarem, o valor de $V_m$ também muda; por isso, não deve ser usado sem especificar as condições.
\end{exemplo}

\section{Massa volúmica de um gás}
Como $n=m/M$ e $V=nV_m$,
\[
\rho=\frac{m}{V}=\frac{M}{V_m}.
\]
Isto permite relacionar propriedades macroscópicas com massa molar.

\section{Misturas gasosas e fração molar}
A \textbf{fração molar} de um componente $i$ é
\[
x_i=\frac{n_i}{n_{\rm total}}.
\]
Numa mistura gasosa ideal, a fração molar é também igual à fração de moléculas desse componente. Percentagens em quantidade correspondem a $100x_i$.

\begin{exemplo}[title=Exemplo resolvido -- composição gasosa]
Uma mistura contém \SI{0.78}{mol} de \ce{N2}, \SI{0.21}{mol} de \ce{O2} e \SI{0.01}{mol} de outros gases. Como $n_{\rm total}=\SI{1.00}{mol}$,
\[
x_{\ce{O2}}=0{,}21,
\]
ou seja, 21\% em quantidade.
\end{exemplo}

\section{Atmosfera terrestre}
A atmosfera é uma mistura gasosa dinâmica. No ar seco junto à superfície, o azoto e o oxigénio dominam largamente a composição; vapor de água é muito variável e, por isso, costuma ser tratado separadamente.

\begin{figure}[H]
\centering
\begin{tikzpicture}
\begin{axis}[ybar,width=11cm,height=6cm,ymin=0,ymax=85,bar width=20pt,ylabel={fração em volume / \%},symbolic x coords={$\ce{N2}$,$\ce{O2}$,Ar,$\ce{CO2}$},xtick=data,grid=major,nodes near coords]
\addplot[eublue,fill=eublue!45] coordinates {($\ce{N2}$,78.0) ($\ce{O2}$,21.0) (Ar,0.93) ($\ce{CO2}$,0.04)};
\end{axis}
\end{tikzpicture}
\caption{Composição aproximada do ar seco. Os valores servem para leitura de escala e comparação; componentes vestigiais não estão representados.}
\end{figure}
A troposfera é constituída maioritariamente por azoto e oxigénio, com argão, dióxido de carbono, vapor de água e outros componentes em menores proporções. A distinção entre um constituinte natural e um poluente depende da concentração, do local e dos efeitos produzidos.

\begin{table}[H]
\centering
\caption{Exemplos de poluentes atmosféricos e fontes possíveis.}
\renewcommand{\arraystretch}{1.18}\begin{tabularx}{\textwidth}{lXX}
\toprule
Poluente & Fontes & Efeitos/questões \\
\midrule
\ce{NO_x} & combustão a alta temperatura & smog fotoquímico, acidificação \\
\ce{SO2} & combustíveis com enxofre, processos industriais & irritação, chuva ácida \\
CO & combustão incompleta & toxicidade por interferência no transporte de oxigénio \\
Partículas & tráfego, indústria, incêndios & qualidade do ar e saúde respiratória \\
\ce{O3} troposférico & reações fotoquímicas secundárias & poluente oxidante \\
\bottomrule
\end{tabularx}\renewcommand{\arraystretch}{1.0}
\end{table}

\begin{table}[H]
\centering
\caption{Composição aproximada do ar seco perto da superfície.}
\begin{tabular}{lc}
\toprule
Componente & Fração em volume aproximada \\
\midrule
\ce{N2} & 78\% \\
\ce{O2} & 21\% \\
Ar & 0,93\% \\
\ce{CO2} & cerca de 0,04\% e variável \\
Outros & vestígios \\
\bottomrule
\end{tabular}
\end{table}

Os gases com efeito de estufa absorvem radiação em determinadas regiões do infravermelho. O efeito de estufa natural é essencial para a temperatura média do planeta; alterações na composição atmosférica podem modificar o balanço radiativo.

\section{Ozono: protetor e poluente}
Na estratosfera, o ozono absorve radiação ultravioleta nociva. Na troposfera, concentrações elevadas de ozono podem ser prejudiciais. O mesmo composto pode, portanto, ter papéis diferentes conforme a região atmosférica.

\begin{figure}[H]
\centering
\begin{tikzpicture}[scale=.95]
\draw[thick,eublue] (0,0) rectangle (3.5,4.3);
\draw[thick,eublue] (6,0) rectangle (9.5,4.3);
\foreach \i in {1,...,18}{\pgfmathsetmacro{\x}{.25+mod(37*\i,29)/10}\pgfmathsetmacro{\y}{.25+mod(19*\i,36)/10}\fill[eucyan] (\x,\y) circle (2pt);}
\foreach \i in {1,...,55}{\pgfmathsetmacro{\x}{6.25+mod(23*\i,29)/10}\pgfmathsetmacro{\y}{.25+mod(31*\i,36)/10}\fill[euorange] (\x,\y) circle (1.8pt);}
\node at (1.75,-.45) {menos concentrada};
\node at (7.75,-.45) {mais concentrada};
\end{tikzpicture}
\caption{Representação microscópica qualitativa: concentração maior significa mais quantidade de soluto por unidade de volume, não partículas ``maiores''.}
\end{figure}

\section{Soluções e concentração}
A concentração molar é
\[
c=\frac{n}{V},
\]
com $V$ em litros quando se pretende $c$ em \si{mol/L}.
A concentração mássica pode escrever-se como
\[
c_m=\frac{m_{\mathrm{soluto}}}{V_{\text{solução}}}.
\]
Também se usam outras formas de expressar composição:
\[
w=\frac{m_{\mathrm{soluto}}}{m_{\text{solução}}},
\qquad
\%\,m/m=100w,
\]
\[
\%\,V/V=100\,\frac{V_{\mathrm{soluto}}}{V_{\text{solução}}}.
\]
Para concentrações muito pequenas podem surgir unidades como mg/L ou ppm, desde que se interprete claramente a definição usada.

Numa diluição, a quantidade de soluto retirada mantém-se, levando à relação
\[
c_iV_i=c_fV_f.
\]

\begin{table}[H]
\centering
\caption{Formas frequentes de expressar composição de soluções.}
\renewcommand{\arraystretch}{1.18}\begin{tabularx}{\textwidth}{lXX}
\toprule
Grandeza & Expressão & Unidade típica \\
\midrule
Concentração molar & $c=n/V$ & mol/L \\
Concentração mássica & $c_m=m_{\mathrm{soluto}}/V$ & g/L \\
Percentagem mássica & $100m_{\mathrm{soluto}}/m_{\text{solução}}$ & \% (m/m) \\
Percentagem volúmica & $100V_{\mathrm{soluto}}/V_{\text{solução}}$ & \% (V/V) \\
\bottomrule
\end{tabularx}\renewcommand{\arraystretch}{1.0}
\end{table}

\begin{exemplo}
Pretende-se preparar \SI{250}{mL} de uma solução de \ce{NaCl} \SI{0.200}{mol/L}. A quantidade necessária é
\[
n=cV=0{,}200\times0{,}250=\SI{0.0500}{mol}.
\]
Usando $M(\ce{NaCl})=\SI{58.44}{g/mol}$,
\[
m=nM=\SI{2.92}{g}.
\]
\end{exemplo}

\begin{exemplo}[title=Exemplo resolvido -- diluição]
Pretendem-se preparar \SI{200}{mL} de solução \SI{0.150}{mol/L} a partir de uma solução-mãe \SI{1.00}{mol/L}. Como
\[
c_iV_i=c_fV_f,
\]
\[
V_i=\frac{0{,}150\times200}{1{,}00}=\SI{30.0}{mL}.
\]
Devem medir-se \SI{30.0}{mL} da solução-mãe com material adequado, transferir para um balão volumétrico de \SI{200}{mL}, completar com solvente até ao traço e homogeneizar.
\end{exemplo}

\begin{atividade}[title=AL 2.2 -- Solução aquosa a partir de um soluto sólido]
\textbf{Etapas essenciais:} calcular a massa necessária; pesar o sólido; dissolver num volume reduzido de água; transferir quantitativamente para o balão volumétrico; lavar o recipiente de dissolução; completar até ao traço com o menisco corretamente posicionado; tapar e homogeneizar.

\textbf{Erros frequentes:} perda de soluto na transferência, soluto não completamente dissolvido, ultrapassar o traço de aferição e leitura incorreta do menisco.
\end{atividade}

\begin{atividade}[title=AL 2.3 -- Diluição de uma solução]
Mede-se uma alíquota da solução-mãe com pipeta adequada, transfere-se para um balão volumétrico e completa-se com solvente. A quantidade de soluto na alíquota não se altera durante a diluição:
\[
n_{\rm antes}=n_{\rm depois}\Rightarrow c_iV_i=c_fV_f.
\]
A segurança depende da substância usada; quando se diluem soluções corrosivas, seguem-se as regras específicas do laboratório.
\end{atividade}

\begin{atividade}[title=Extensão -- preparação de uma solução]
\textbf{Preparação de uma solução.} Determinar a massa de soluto, dissolver em pequeno volume de água destilada, transferir quantitativamente para um balão volumétrico, lavar o recipiente de dissolução para o balão, completar até ao traço e homogeneizar. Avaliar criticamente erros como leitura incorreta do menisco, perda de soluto ou enchimento acima do traço.
\end{atividade}

\begin{stem}
\textbf{Qualidade do ar na escola.} Construir um estudo com dados públicos de qualidade do ar, temperatura e tráfego. Formular uma pergunta, recolher dados durante vários dias, construir gráficos e discutir correlações sem confundir correlação com causalidade.
\end{stem}

\begin{exercicios}
\begin{enumerate}
  \item Um gás ocupa \SI{4.80}{L} nas condições em que $V_m=\SI{24.0}{L/mol}$. Determina $n$.
  \item Uma amostra de \ce{CO2} contém \SI{0.300}{mol}. Qual o volume se $V_m=\SI{24.5}{L/mol}$?
  \item Calcula a concentração molar de uma solução obtida dissolvendo \SI{5.85}{g} de \ce{NaCl} em água até \SI{500}{mL}. Usa $M=\SI{58.5}{g/mol}$.
  \item Que volume de uma solução \SI{1.00}{mol/L} é necessário para preparar \SI{250}{mL} de solução \SI{0.100}{mol/L}?
  \item Explica por que motivo o ozono pode ser simultaneamente considerado protetor e poluente.
  \item \textbf{Integração.} Um cilindro contém \SI{12.0}{L} de um gás, nas condições em que $V_m=\SI{24.0}{L/mol}$. Se a massa é \SI{22.0}{g}, estima a massa molar do gás.
  \item Uma mistura gasosa contém \SI{2.0}{mol} de \ce{N2} e \SI{0.50}{mol} de \ce{O2}. Calcula a fração molar de cada componente.
  \item Uma solução contém \SI{12}{g} de soluto em \SI{300}{g} de solução. Determina a percentagem mássica.
  \item Pretendem-se preparar \SI{500}{mL} de solução \SI{0.0200}{mol/L} a partir de uma solução \SI{0.500}{mol/L}. Determina o volume de solução-mãe necessário.
\end{enumerate}
\end{exercicios}

\chapter{Transformações químicas e energia}

\begin{objetivos}
\begin{itemize}
  \item interpretar reações em termos de quebra e formação de ligações;
  \item distinguir processos exotérmicos e endotérmicos;
  \item interpretar o sinal de $\Delta H$;
  \item relacionar qualitativamente $\Delta H$ com energias de ligação;
  \item compreender reações fotoquímicas e o papel de radicais livres.
\end{itemize}
\end{objetivos}

\section{Reações químicas e energia}
Uma reação química reorganiza átomos. Algumas ligações são quebradas e outras são formadas. A quebra de ligações requer energia; a formação de ligações liberta energia.

A variação de entalpia é definida por
\[
\Delta H=H_{\rm produtos}-H_{\rm reagentes}.
\]
A pressão constante, o calor trocado pelo sistema relaciona-se com $\Delta H$.

\begin{center}
\begin{tikzpicture}[x=1.25cm,y=1cm]
\draw[-{Stealth}] (0,0)--(7,0) node[right] {progresso da reação};
\draw[-{Stealth}] (0,0)--(0,4) node[above] {$H$};
\draw[thick,eublue] (.6,3.0)--(2.2,3.0);
\draw[thick,eublue] (4.8,1.2)--(6.5,1.2);
\draw[euorange,thick,-{Stealth}] (3.7,2.9)--(3.7,1.3);
\node[left] at (3.7,2.1) {$\Delta H<0$};
\node at (1.4,3.35) {reagentes};
\node at (5.65,1.55) {produtos};
\end{tikzpicture}
\end{center}

Se $\Delta H<0$, a reação é exotérmica; se $\Delta H>0$, é endotérmica.

\begin{figure}[H]
\centering
\begin{tikzpicture}
\begin{axis}[width=11cm,height=6.2cm,xlabel={coordenada de reação},ylabel={energia potencial},xmin=0,xmax=10,ymin=0,ymax=8,xtick=\empty,ytick=\empty,axis lines=left]
\addplot[domain=.5:9.5,samples=150,thick,eublue]{2.0+4.7*exp(-((x-4.8)^2)/2.1)-0.16*x};
\draw[dashed,gray] (axis cs:.7,1.9)--(axis cs:9.2,1.9);
\draw[<->,euorange,thick] (axis cs:2.0,1.75)--(axis cs:2.0,6.15);
\node[anchor=west,font=\small] at (axis cs:2.15,4.15) {$E_a$};
\node[font=\small] at (axis cs:1.3,1.35) {reagentes};
\node[font=\small] at (axis cs:8.6,.65) {produtos};
\node[font=\small] at (axis cs:4.8,7.15) {estado de transição};
\end{axis}
\end{tikzpicture}
\caption{Perfil energético qualitativo de uma reação exotérmica, incluindo a energia de ativação.}
\end{figure}

\begin{conceito}
Uma reação energeticamente favorável não tem de ser rápida. A \textbf{energia de ativação} ajuda a explicar por que algumas transformações necessitam de aquecimento, luz ou catalisadores para ocorrer a uma velocidade observável. Um catalisador altera o caminho reacional e reduz a energia de ativação, mas não altera o valor global de $\Delta H$ entre reagentes e produtos.
\end{conceito}

\begin{center}
\begin{tikzpicture}[x=1cm,y=.75cm]
\node at (1.5,4.4) {\textbf{exotérmica}};
\draw[-{Stealth}] (0,0)--(0,4);
\draw[thick,eublue] (.4,3.1)--(2.2,3.1);
\draw[thick,eugreen] (.4,1.4)--(2.2,1.4);
\draw[eured,thick,-{Stealth}] (2.5,3)--(2.5,1.5);
\node[right] at (2.5,2.25) {$\Delta H<0$};

\node at (6.0,4.4) {\textbf{endotérmica}};
\draw[-{Stealth}] (4.3,0)--(4.3,4);
\draw[thick,eublue] (4.7,1.4)--(6.5,1.4);
\draw[thick,eugreen] (4.7,3.1)--(6.5,3.1);
\draw[eured,thick,-{Stealth}] (6.8,1.5)--(6.8,3);
\node[right] at (6.8,2.25) {$\Delta H>0$};
\end{tikzpicture}
\end{center}

\begin{exemplo}[title=Exemplo resolvido -- interpretação de $\Delta H$]
Se uma reação liberta \SI{85}{kJ} por mole de reação, a entalpia dos produtos é inferior à dos reagentes e
\[
\Delta H=-\SI{85}{kJ/mol}.
\]
O sinal negativo descreve o sistema: energia foi transferida para a vizinhança.
\end{exemplo}

\section{Estimativa com energias de ligação}
Uma aproximação útil é
\[
\Delta H\approx\sum E(\text{ligações quebradas})-\sum E(\text{ligações formadas}).
\]
Esta abordagem ajuda a ligar o nível microscópico ao balanço energético macroscópico.

\begin{exemplo}[title=Exemplo resolvido -- energias de ligação]
Numa reação hipotética, quebram-se ligações com energia total \SI{950}{kJ/mol} e formam-se ligações cuja soma é \SI{1120}{kJ/mol}. Então
\[
\Delta H\approx950-1120=-\SI{170}{kJ/mol},
\]
pelo que a reação é exotérmica.
\end{exemplo}

\section{Fotoquímica}
Numa reação fotoquímica, a luz fornece energia para iniciar transformações. A energia do fotão deve ser compatível com a transição ou processo envolvido. A reação de escurecimento de certos sais de prata é um exemplo clássico explorável em contexto laboratorial.

\begin{atividade}[title=AL 2.4 -- Reação fotoquímica do cloreto de prata]
O cloreto de prata escurece quando exposto à luz. Numa descrição simplificada, a radiação desencadeia a transformação fotoquímica
\[
\ce{2AgCl(s) ->[luz] 2Ag(s) + Cl2(g)}.
\]
Em contexto escolar, comparam-se amostras sujeitas a diferentes condições de iluminação e uma amostra de controlo protegida da luz. O objetivo é avaliar o efeito da radiação, não produzir quantidades apreciáveis de produtos.

\textbf{Segurança e resíduos:} sais de prata devem ser usados em pequena escala, com luvas/óculos quando indicado, e todos os resíduos devem ser recolhidos segundo o procedimento definido pelo professor.
\end{atividade}

\section{Ozono e fotoquímica atmosférica}
Uma sequência simplificada para a formação de ozono estratosférico é
\[
\ce{O2} + h\nu \longrightarrow 2\ce{O},
\qquad
\ce{O + O2 -> O3}.
\]
O ozono também absorve ultravioleta:
\[
\ce{O3} + h\nu \longrightarrow \ce{O2 + O}.
\]
Assim, existe um ciclo fotoquímico dinâmico. Na troposfera, o ozono pode formar-se em redes de reações envolvendo óxidos de azoto e compostos orgânicos sob luz solar, tornando-se um poluente oxidante.

\section{Radicais livres}
Radicais são espécies com eletrões desemparelhados. Podem ser muito reativos, participando em mecanismos atmosféricos e em processos biológicos. A palavra ``radical'' descreve uma característica eletrónica e não significa, por si só, ``substância perigosa''.

Na estratosfera, certos radicais podem participar em ciclos catalíticos de destruição do ozono. Em sistemas biológicos, espécies radicalares podem iniciar reações de oxidação em cadeia; os organismos possuem, por isso, mecanismos antioxidantes de controlo. A química dos radicais ajuda a compreender processos atmosféricos, envelhecimento de materiais e algumas transformações em seres vivos.

\begin{mundo}
Bolsas térmicas instantâneas, combustão, fotossíntese, cozinhar alimentos e armazenamento químico de energia são exemplos em que transformações químicas e energia estão ligadas. A engenharia química procura controlar velocidade, rendimento, segurança e eficiência energética desses processos.
\end{mundo}

\begin{exercicios}
\begin{enumerate}
  \item Classifica como exo ou endotérmico um processo com $\Delta H=-\SI{125}{kJ/mol}$.
  \item Uma reação tem produtos com entalpia \SI{40}{kJ/mol} acima dos reagentes. Qual o sinal e valor de $\Delta H$?
  \item Estima $\Delta H$ para uma reação hipotética que quebra ligações totalizando \SI{620}{kJ/mol} e forma ligações totalizando \SI{760}{kJ/mol}.
  \item Explica por que razão quebrar ligações requer energia.
  \item Calcula a energia de um fotão UV de \SI{300}{nm}.
  \item \textbf{Integração.} Uma reação absorve \SI{2.0}{kJ}. Se essa energia fosse fornecida exclusivamente por fotões de \SI{400}{nm}, estima o número de fotões necessários.
  \item Uma reação quebra ligações totalizando \SI{1520}{kJ/mol} e forma ligações totalizando \SI{1380}{kJ/mol}. Estima $\Delta H$ e classifica o processo.
  \item Explica a função da amostra protegida da luz na experiência com cloreto de prata.
  \item Escreve uma frase que distinga o papel do ozono estratosférico do ozono troposférico.
\end{enumerate}
\end{exercicios}

% =========================================================
\part{Física — Energia e sua conservação}
% =========================================================

\chapter{Energia e movimentos}

\begin{objetivos}
\begin{itemize}
  \item interpretar transferências de energia por trabalho;
  \item aplicar o teorema da energia cinética;
  \item usar energia potencial gravítica e conservação da energia mecânica;
  \item distinguir forças conservativas e não conservativas;
  \item interpretar perdas de energia mecânica como transformação/transferência de energia.
\end{itemize}
\end{objetivos}

\section{Sistema, fronteira e energia}
Antes de resolver um problema, define o \textbf{sistema}. A escolha do sistema determina que interações são internas e quais representam transferências de energia através da fronteira.

\subsection*{Sistema, vizinhança e fronteira}
A \textbf{vizinhança} é tudo o que não pertence ao sistema. A \textbf{fronteira} é o limite real ou imaginário que separa o sistema da vizinhança. Em termodinâmica é útil distinguir:
\begin{itemize}
  \item \textbf{sistema aberto}: pode trocar matéria e energia com a vizinhança;
  \item \textbf{sistema fechado}: pode trocar energia, mas não matéria;
  \item \textbf{sistema isolado}: idealmente não troca nem matéria nem energia.
\end{itemize}
Uma transformação de energia ocorre no interior do sistema; uma transferência de energia atravessa a fronteira.

\begin{figure}[H]
\centering
\begin{tikzpicture}[>=Stealth,node distance=.9cm]
\tikzset{sys/.style={rounded corners=3mm,minimum width=3.5cm,minimum height=2cm,draw=eublue,fill=eulightblue,align=center,font=\bfseries}}
\node[sys] (a) {aberto};
\node[sys,right=1.1cm of a] (f) {fechado};
\node[sys,right=1.1cm of f] (i) {isolado};
\draw[<->,thick,euorange] ($(a.west)+(-.9,.35)$)--(a.west) node[midway,above,font=\scriptsize]{energia};
\draw[<->,thick,eugreen] ($(a.west)+(-.9,-.35)$)--(a.west) node[midway,below,font=\scriptsize]{matéria};
\draw[<->,thick,euorange] ($(f.west)+(-.9,0)$)--(f.west) node[midway,above,font=\scriptsize]{energia};
\draw[thick,eured] ($(i.west)+(-.55,-.35)$)--($(i.west)+(-.05,.35)$);
\draw[thick,eured] ($(i.west)+(-.55,.35)$)--($(i.west)+(-.05,-.35)$);
\end{tikzpicture}
\caption{Classificação idealizada de sistemas segundo as trocas através da fronteira.}
\end{figure}

A energia cinética de um corpo de massa $m$ e rapidez $v$ é
\[
E_c=\frac12 mv^2.
\]
Perto da superfície terrestre, a energia potencial gravítica pode escrever-se, relativamente a um nível de referência,
\[
E_{p,g}=mgh.
\]

\subsection*{Modelo do centro de massa}
Em muitas situações, um corpo extenso pode ser modelado como uma partícula localizada no seu centro de massa quando apenas interessa descrever o movimento global. Este modelo não é adequado se a rotação ou a deformação forem essenciais para o problema.

\begin{exemplo}[title=Exemplo resolvido -- energia cinética]
Um ciclista e bicicleta têm massa total \SI{80}{kg} e movem-se a \SI{7.0}{m/s}. A energia cinética do sistema, modelado pelo seu centro de massa, é
\[
E_c=\frac12(80)(7{,}0)^2=\SI{1.96e3}{J}.
\]
Se a rapidez duplicar, a energia cinética quadruplica.
\end{exemplo}

\begin{figure}[H]
\centering
\begin{tikzpicture}
\begin{axis}[width=10.8cm,height=5.8cm,xlabel={rapidez $v$ / m s$^{-1}$},ylabel={$E_c$ / kJ},xmin=0,xmax=30,ymin=0,ymax=550,grid=major]
\addplot[domain=0:30,samples=100,thick,eublue]{0.6*x^2};
\node[anchor=west,font=\small] at (axis cs:18,120) {$E_c\propto v^2$};
\end{axis}
\end{tikzpicture}
\caption{Energia cinética em função da rapidez para um corpo de massa constante (curva ilustrativa).}
\end{figure}

\begin{leituragrafico}
Duplicar a rapidez multiplica a energia cinética por quatro. Esta relação quadrática é central na análise de segurança rodoviária: a energia que os travões e a deformação do veículo têm de transformar cresce muito rapidamente com a velocidade.
\end{leituragrafico}


\begin{figure}[H]
\centering
\begin{tikzpicture}
\begin{axis}[width=10.8cm,height=5.8cm,xlabel={$v^2$ / m$^2$ s$^{-2}$},ylabel={$E_c$ / kJ},xmin=0,xmax=900,ymin=0,ymax=550,grid=major]
\addplot[domain=0:900,thick,eugreen]{0.6*x};
\node[anchor=west,font=\small] at (axis cs:430,180) {declive $=m/2$};
\end{axis}
\end{tikzpicture}
\caption{Para massa constante, $E_c$ é linear em função de $v^2$. O declive do gráfico permite determinar a massa.}
\end{figure}

\begin{exemplo}[title=Exemplo resolvido -- massa a partir de um gráfico]
Num gráfico $E_c$ em função de $v^2$, o declive é \SI{40}{kg}. Como
\[
E_c=\frac12mv^2,
\]
o declive é $m/2$. Portanto,
\[
m=2\times40=\SI{80}{kg}.
\]
\end{exemplo}

\section{Trabalho de uma força}
Para uma força constante, o trabalho é
\[
W=Fd\cos\theta,
\]
onde $\theta$ é o ângulo entre a força e o deslocamento.

\begin{figure}[H]
\centering
\begin{tikzpicture}[>=Stealth,scale=1]
\coordinate (A) at (0,0); \coordinate (B) at (7,0); \coordinate (C) at (7,3.6);
\draw[thick,eudark] (A)--(B)--(C)--cycle;
\draw[fill=eublue!20,rotate around={27:(3.8,1.95)}] (3.35,1.65) rectangle (4.25,2.25);
\draw[->,very thick,euorange] (3.8,1.95)--++(1.8,0.92) node[above] {$\vec F$};
\draw[->,very thick,eublue] (3.8,1.95)--++(1.7,.88) node[pos=.55,below] {$\vec d$};
\draw[->,thick,eured] (3.8,1.95)--++(0,-1.7) node[right] {$\vec P$};
\draw[eudark,->] (.9,0) arc (0:27:.9);
\node[eudark] at (1.18,.23) {$\theta$};
\end{tikzpicture}
\caption{Num plano inclinado, o trabalho depende da componente da força paralela ao deslocamento.}
\end{figure}

O teorema da energia cinética afirma
\[
W_{\rm resultante}=\Delta E_c.
\]

\begin{center}
\begin{tikzpicture}[scale=1]
\draw[thick] (-.5,0)--(8,0);
\draw[fill=eublue!15] (1,.05) rectangle (2.2,.85);
\node at (1.6,.45) {$m$};
\draw[-{Stealth},very thick,euorange] (2.2,.6)--(4.4,.6) node[midway,above] {$\vect F$};
\draw[-{Stealth},thick] (1.6,-.35)--(5.8,-.35) node[midway,below] {$\vect d$};
\draw[-{Stealth},thick,eured] (1.0,.25)--(-.2,.25) node[midway,above] {atrito};
\end{tikzpicture}
\end{center}

\begin{exemplo}
Um carrinho de \SI{2.0}{kg}, inicialmente em repouso, é puxado por uma força horizontal resultante constante de \SI{6.0}{N} ao longo de \SI{4.0}{m}.
\[
W_R=Fd=6{,}0\times4{,}0=\SI{24}{J}.
\]
Logo, $\Delta E_c=\SI{24}{J}$ e
\[
\frac12 mv^2=24 \Rightarrow v=\sqrt{24}=\SI{4.9}{m/s}.
\]
\end{exemplo}

\subsection*{Sinal do trabalho e força gravítica}
O trabalho é positivo se a força tiver componente no sentido do deslocamento, negativo se tiver componente oposta e nulo se for perpendicular ao deslocamento.

Para a força gravítica, perto da superfície terrestre,
\[
W_P=-\Delta E_{p,g}.
\]
Esta relação exprime o caráter conservativo do peso: o trabalho depende apenas das posições inicial e final, não do caminho seguido.

\begin{exemplo}[title=Exemplo resolvido -- subida e trabalho do peso]
Uma mochila de \SI{5.0}{kg} sobe \SI{2.0}{m}. A variação de energia potencial é
\[
\Delta E_{p,g}=mgh=5{,}0\times9{,}8\times2{,}0=\SI{98}{J}.
\]
Logo, o trabalho realizado pelo peso é \SI{-98}{J}.
\end{exemplo}

\section{Energia mecânica}
Define-se
\[
E_m=E_c+E_{p,g}.
\]
Quando apenas forças conservativas realizam trabalho relevante,
\[
E_{m,i}=E_{m,f}.
\]
Quando há atrito, a energia mecânica pode diminuir, mas a energia total não desaparece: é transferida para energia interna, som, deformações, etc.

\subsection*{Forças não conservativas e variação da energia mecânica}
Quando forças não conservativas, como o atrito, realizam trabalho,
\[
W_{\rm nc}=\Delta E_m.
\]
Se $W_{\rm nc}<0$, a energia mecânica diminui, mas a energia total não desaparece: aumenta, por exemplo, a energia interna dos corpos e do meio envolvente.

\begin{exemplo}[title=Exemplo resolvido -- plano inclinado com atrito]
Um bloco de \SI{2.0}{kg} parte do repouso e desce uma rampa com desnível de \SI{3.0}{m}. Percorre \SI{7.0}{m} sobre a rampa, sujeito a uma força média de atrito de \SI{5.0}{N}. Tomando o ponto final como referência de energia potencial,
\[
E_{m,i}=mgh=2{,}0\times9{,}8\times3{,}0=\SI{58.8}{J}.
\]
O trabalho do atrito é
\[
W_{\rm atr}=-F_{\rm atr}d=-5{,}0\times7{,}0=\SI{-35}{J}.
\]
Logo,
\[
E_{m,f}=58{,}8-35=\SI{23.8}{J}=\frac12mv^2,
\]
e $v\approx\SI{4.9}{m/s}$. A energia total conservou-se; \SI{35}{J} de energia mecânica foram transformados sobretudo em energia interna.
\end{exemplo}

\subsection*{Potência mecânica}
A potência média mede a rapidez de transferência de energia:
\[
P_{\rm med}=\frac{W}{\Delta t}.
\]
Para uma força constante e velocidade paralela à força, pode escrever-se $P=Fv$ num dado instante.

\begin{atividade}[title=AL 1.1 -- Movimento num plano inclinado: $\Delta E_c$ e distância]
Larga-se um carrinho de uma posição fixa numa rampa e mede-se a rapidez em diferentes posições, por exemplo com fotocélula ou vídeo. Para cada distância $d$, calcula-se a energia cinética. Representa-se $E_c$ ou $\Delta E_c$ em função de $d$.

Se a resultante ao longo do plano for aproximadamente constante,
\[
\Delta E_c=F_R d,
\]
pelo que o gráfico deve ser aproximadamente linear e o declive tem unidade \si{J/m}=\si{N}. A repetição de medições permite estimar dispersão e avaliar o ajuste.
\end{atividade}

\begin{figure}[H]
\centering
\begin{tikzpicture}
\begin{axis}[width=10.5cm,height=5.8cm,xlabel={altura de queda $h_q$ / m},ylabel={altura de ressalto $h_r$ / m},xmin=0,xmax=2.2,ymin=0,ymax=1.6,grid=major]
\addplot[only marks,mark=*,eublue] coordinates{(.4,.25)(.8,.52)(1.2,.76)(1.6,1.02)(2.0,1.27)};
\addplot[domain=0:2.1,thick,euorange]{0.64*x};
\node[font=\small,anchor=west] at (axis cs:1.15,.42) {$h_r\approx0{,}64h_q$};
\end{axis}
\end{tikzpicture}
\caption{Dados ilustrativos de uma experiência de queda e ressalto.}
\end{figure}

\begin{atividade}[title=AL 1.2 -- Queda e ressalto de uma bola]
Larga-se uma bola de diferentes alturas $h_q$ e mede-se a primeira altura de ressalto $h_r$. Para uma mesma bola e superfície, é frequente obter aproximadamente
\[
h_r=a\,h_q.
\]
O declive $a<1$ indica a fração aproximada de energia mecânica que permanece após a colisão, pois imediatamente antes e depois da colisão as energias associadas às alturas são proporcionais a $mgh$.

A percentagem de energia mecânica dissipada na colisão pode estimar-se por
\[
\%E_{\rm diss}\approx(1-a)\times100.
\]
\end{atividade}

\begin{exemplo}[title=Exemplo resolvido -- ressalto]
Se o gráfico $h_r(h_q)$ tiver declive $0{,}64$, cerca de 64\% da energia mecânica associada à altura é recuperada após a colisão e cerca de
\[
(1-0{,}64)\times100=36\%
\]
é transformada noutras formas durante o impacto.
\end{exemplo}

\begin{atividade}[title=Extensão -- energia cinética e distância]
\textbf{Energia cinética e distância.} Usar um carrinho, pista e sistema de aquisição de dados para investigar a relação entre $\Delta E_c$ e a distância percorrida sob uma força resultante aproximadamente constante. Representar $\Delta E_c$ em função de $d$. O declive esperado corresponde à força resultante.
\end{atividade}

\begin{stem}
\textbf{Travagem segura.} Investiga como varia a distância de travagem com a velocidade inicial. Usa um modelo energético simples: o trabalho da força média de travagem iguala a diminuição da energia cinética. Discute por que duplicar a velocidade não duplica a distância de travagem quando a força média se mantém aproximadamente constante.
\end{stem}


% ---------- Enriquecimento a partir das fichas Fis_FT01--Fis_FT06 ----------
\section{Ler a energia através de gráficos}
As fichas de trabalho de energia e movimento mostram uma ideia especialmente útil: muitas relações físicas tornam-se mais claras quando se escolhem variáveis que transformam a relação num gráfico linear. Isso permite não só prever tendências, mas também \textbf{determinar grandezas físicas a partir do declive}.

\subsection*{Energia cinética: de uma parábola a uma reta}
Como
\[
E_c=\frac12mv^2,
\]
um gráfico $E_c(v)$ é parabólico, mas um gráfico $E_c(v^2)$ é uma reta de declive $m/2$.

\begin{figure}[H]
\centering
\begin{minipage}{0.48\textwidth}
\centering
\begin{tikzpicture}
\begin{axis}[width=\linewidth,height=5.1cm,xlabel={$v$ / \si{m.s^{-1}}},ylabel={$E_c$ / J},xmin=0,xmax=8,ymin=0,ymax=135,grid=major]
\addplot[domain=0:8,samples=100,very thick,eublue]{2*x^2};
\end{axis}
\end{tikzpicture}
\small $m=\SI{4.0}{kg}$: relação quadrática.
\end{minipage}\hfill
\begin{minipage}{0.48\textwidth}
\centering
\begin{tikzpicture}
\begin{axis}[width=\linewidth,height=5.1cm,xlabel={$v^2$ / \si{m^2.s^{-2}}},ylabel={$E_c$ / J},xmin=0,xmax=64,ymin=0,ymax=135,grid=major]
\addplot[domain=0:64,samples=2,very thick,eugreen]{2*x};
\node[font=\scriptsize,anchor=west] at (axis cs:34,82) {declive $=m/2$};
\end{axis}
\end{tikzpicture}
\small A linearização permite obter a massa.
\end{minipage}
\caption{Duas representações dos mesmos dados de energia cinética.}
\end{figure}

\begin{exemplo}[title=Exemplo resolvido -- determinar a massa pelo declive]
Num ensaio, a reta $E_c=f(v^2)$ passa pelos pontos $(\SI{9.0}{m^2/s^2},\SI{13.5}{J})$ e $(\SI{25.0}{m^2/s^2},\SI{37.5}{J})$. O declive é
\[
a=\frac{37{,}5-13{,}5}{25{,}0-9{,}0}=\SI{1.50}{kg}.
\]
Como $a=m/2$, conclui-se
\[
m=2a=\SI{3.00}{kg}.
\]
Esta estratégia é mais robusta do que determinar a massa a partir de um único ponto experimental.
\end{exemplo}

\subsection*{Trabalho como área e como declive}
Para uma força constante cuja componente na direção do deslocamento é $F_{\parallel}$,
\[
W=F_{\parallel}d.
\]
Assim, num gráfico $W=f(d)$, o declive representa $F_{\parallel}$. Se a força faz um ângulo $\theta$ com o deslocamento, então $F_{\parallel}=F\cos\theta$.

\begin{figure}[H]
\centering
\begin{tikzpicture}
\begin{axis}[width=11cm,height=5.8cm,xlabel={$d$ / m},ylabel={$W$ / J},xmin=0,xmax=8,ymin=-55,ymax=130,grid=major,legend pos=north west]
\addplot[domain=0:8,very thick,eugreen]{15*x};
\addlegendentry{$F_{\parallel}=\SI{15}{N}$}
\addplot[domain=0:8,very thick,eublue]{7.5*x};
\addlegendentry{$F_{\parallel}=\SI{7.5}{N}$}
\addplot[domain=0:8,very thick,eured]{-6*x};
\addlegendentry{força resistente}
\end{axis}
\end{tikzpicture}
\caption{O declive de $W(d)$ é a componente da força na direção do deslocamento.}
\end{figure}

\begin{leituragrafico}
Uma reta horizontal em $W(d)$ teria declive nulo: a força não teria componente segundo o deslocamento e, nesse modelo, não realizaria trabalho. Uma reta com declive negativo representa trabalho resistente.
\end{leituragrafico}

\subsection*{Energia ao longo de uma trajetória}
Num movimento sem dissipação, a energia mecânica conserva-se. A energia potencial pode transformar-se em energia cinética e vice-versa.

\begin{figure}[H]
\centering
\begin{tikzpicture}
\begin{axis}[width=11cm,height=6cm,xlabel={posição ao longo do percurso / u.a.},ylabel={energia / J},xmin=0,xmax=10,ymin=0,ymax=105,grid=major,legend style={at={(0.5,-0.22)},anchor=north,legend columns=3}]
\addplot[very thick,eublue,smooth] coordinates{(0,95)(2,75)(4,48)(6,25)(8,12)(10,5)};
\addlegendentry{$E_{p,g}$}
\addplot[very thick,euorange,smooth] coordinates{(0,5)(2,25)(4,52)(6,75)(8,88)(10,95)};
\addlegendentry{$E_c$}
\addplot[very thick,eugreen,dashed] coordinates{(0,100)(10,100)};
\addlegendentry{$E_m$}
\end{axis}
\end{tikzpicture}
\caption{Transformação de energia potencial gravítica em cinética num modelo sem dissipação.}
\end{figure}

\begin{mundo}
\textbf{Segurança rodoviária.} Como $E_c\propto v^2$, passar de \SI{50}{km/h} para \SI{100}{km/h} não duplica a energia cinética: quadruplica-a. Em travagem, essa energia tem de ser transferida para os travões, pneus, pavimento, ar e deformações do veículo. É por isso que a velocidade tem um efeito tão forte nas distâncias e consequências de travagem.
\end{mundo}

\section{Potência, tempo e rendimento em situações reais}
A ficha dedicada à potência reforça que a mesma quantidade de energia pode ser transferida em tempos muito diferentes. A potência média é
\[
P=\frac{\Delta E}{\Delta t}.
\]
Quando se compara um dispositivo real com a energia que recebe, usa-se o rendimento
\[
\eta=\frac{E_{\text{útil}}}{E_{\text{fornecida}}}\times100\%
=\frac{P_{\text{útil}}}{P_{\text{fornecida}}}\times100\%.
\]

\begin{figure}[H]
\centering
\begin{tikzpicture}[>=Stealth,node distance=1.4cm]
\node[rounded corners=2mm,draw=eublue,fill=eulightblue,minimum width=3.1cm,minimum height=1cm] (in) {energia fornecida};
\node[rounded corners=2mm,draw=eugreen,fill=eulightgreen,minimum width=3.1cm,minimum height=1cm,right=3.6cm of in] (use) {energia útil};
\node[rounded corners=2mm,draw=eured,fill=red!4,minimum width=3.1cm,minimum height=1cm,below=1.4cm of use] (dis) {energia dissipada};
\draw[->,very thick,eugreen] (in)--(use) node[midway,above,font=\small] {$E_{\text{útil}}$};
\draw[->,very thick,eured] (in.south east)--(dis.west) node[midway,below left,font=\small] {$E_{\mathrm{dissipada}}$};
\end{tikzpicture}
\caption{Balanço energético simplificado de um dispositivo real.}
\end{figure}

\begin{exemplo}[title=Exemplo resolvido -- elevador e potência]
Um elevador transporta uma carga total de \SI{650}{kg} por um desnível de \SI{12}{m} em \SI{18}{s}. Desprezando variações de energia cinética,
\[
E_{\text{útil}}=mgh=650\times9{,}8\times12\approx\SI{7.64e4}{J}.
\]
A potência útil média é
\[
P_{\text{útil}}\approx\frac{7{,}64\times10^4}{18}=\SI{4.24}{kW}.
\]
Se o motor receber \SI{5.30}{kW}, o rendimento é aproximadamente
\[
\eta=\frac{4{,}24}{5{,}30}\times100\approx\SI{80}{\percent}.
\]
\end{exemplo}

\begin{exercicios}
\begin{enumerate}
  \item Calcula a energia cinética de um automóvel de \SI{1200}{kg} a \SI{20}{m/s}.
  \item Um objeto de \SI{0.50}{kg} sobe \SI{3.0}{m}. Calcula a variação de energia potencial gravítica, usando $g=\SI{9.8}{m/s^2}$.
  \item Uma força de \SI{10}{N} atua ao longo de \SI{5}{m} formando \ang{60} com o deslocamento. Calcula o trabalho.
  \item Um corpo de \SI{1.5}{kg} passa de \SI{2.0}{m/s} para \SI{6.0}{m/s}. Determina o trabalho da resultante.
  \item Uma bola é largada de \SI{5.0}{m}. Desprezando a resistência do ar, calcula a rapidez imediatamente antes do solo.
  \item Um bloco desce uma rampa e perde \SI{8}{J} de energia mecânica devido ao atrito. Interpreta fisicamente essa ``perda''.
  \item \textbf{Integração.} Um elevador de massa total \SI{900}{kg} sobe \SI{12}{m} em \SI{20}{s}. Estima o aumento de energia potencial e a potência média mínima necessária, ignorando perdas.
  \item Uma bola é largada de \SI{1.20}{m} e ressalta até \SI{0.75}{m}. Estima a percentagem de energia mecânica dissipada na colisão.
  \item O gráfico $\Delta E_c(d)$ de um carrinho tem declive \SI{2.8}{J/m}. Interpreta fisicamente o valor.
  \item Um motor realiza \SI{36}{kJ} de trabalho útil em \SI{12}{s}. Determina a potência média.
  \item Uma força de \SI{50}{N} atua durante um deslocamento de \SI{4.0}{m}, formando \ang{30} com o deslocamento. Determina o trabalho realizado pela força.
  \item Num gráfico $E_c(v^2)$, o declive é \SI{40}{kg}. Determina a massa do sistema.
  \item Um bloco de \SI{2.0}{kg} desce, a partir do repouso, um desnível de \SI{3.0}{m}. A força média de atrito é \SI{5.0}{N} e o percurso sobre a rampa tem \SI{7.0}{m}. Determina a energia cinética final e a rapidez final.
  \item Um dispositivo transfere \SI{500}{J} de energia útil em \SI{4.0}{s}. Determina a potência média.
  \item Justifica, por um argumento energético, por que razão duplicar a rapidez de um automóvel pode aproximadamente quadruplicar a distância de travagem se a força média de travagem se mantiver.
\end{enumerate}
\end{exercicios}

\chapter{Energia e fenómenos elétricos}

\begin{objetivos}
\begin{itemize}
  \item interpretar corrente elétrica e diferença de potencial;
  \item montar e analisar circuitos simples em série e em paralelo;
  \item aplicar relações entre tensão, corrente, resistência, potência e energia;
  \item compreender a função de um gerador real de forma qualitativa;
  \item explicar o efeito Joule e balanços energéticos em circuitos.
\end{itemize}
\end{objetivos}

\section{Corrente e diferença de potencial}
A corrente elétrica mede a taxa de passagem de carga:
\[
I=\frac{\Delta Q}{\Delta t}.
\]
A diferença de potencial entre dois pontos relaciona a energia transferida por unidade de carga:
\[
U=\frac{\Delta E}{\Delta Q}.
\]

Os portadores de carga dependem do meio material. Num metal, a corrente está associada ao movimento orientado de eletrões de condução; numa solução eletrolítica, participam iões positivos e negativos que se deslocam em sentidos opostos; num gás ionizado, podem contribuir iões e eletrões. Esta distinção liga o estudo dos circuitos à estrutura microscópica da matéria.

\begin{table}[H]
\centering
\caption{Portadores de carga em diferentes meios condutores.}
\begin{tabularx}{\textwidth}{lX}
\toprule
Meio & Portadores de carga dominantes \\
\midrule
Metal & eletrões de condução \\
Eletrólito em solução ou fundido & catiões e aniões \\
Gás ionizado/plasma & eletrões e iões \\
\bottomrule
\end{tabularx}
\end{table}

\subsection*{Sentido convencional e medição em circuitos}
O sentido convencional da corrente é definido como o sentido do movimento de cargas positivas. Num metal, os eletrões de condução deslocam-se, em média, no sentido oposto. Um amperímetro ideal é ligado \textbf{em série}; um voltímetro ideal é ligado \textbf{em paralelo} com o componente cuja diferença de potencial se pretende medir.

\section{Lei de Ohm e resistência}
Para um condutor óhmico a temperatura aproximadamente constante,
\[
U=RI.
\]

\begin{center}
\begin{circuitikz}
\draw (0,0) to[battery1,l=$\varepsilon$] (0,3)
      to[switch] (2,3)
      to[R,l=$R$] (5,3)
      -- (5,0) -- (0,0);
\draw (2.6,3) node[above] {$I$};
\end{circuitikz}
\end{center}

\subsection*{Condutores óhmicos e não óhmicos}
Num condutor óhmico a temperatura aproximadamente constante, o gráfico $U(I)$ é linear e passa pela origem.
 O declive do gráfico $U(I)$ corresponde à resistência quando a relação é óhmica.

\begin{figure}[H]
\centering
\begin{tikzpicture}
\begin{axis}[width=11cm,height=6.2cm,xlabel={$I$ / A},ylabel={$U$ / V},xmin=0,xmax=1.25,ymin=0,ymax=12.5,grid=major,legend pos=north west]
\addplot[domain=0:1.2,thick,eublue]{8*x};
\addplot[domain=0:1.2,samples=100,thick,euorange]{4.0*x+4.2*x^2};
\legend{resistor óhmico,lâmpada de filamento (qualitativo)}
\end{axis}
\end{tikzpicture}
\caption{Características elétricas qualitativas de um componente óhmico e de um componente não óhmico.}
\end{figure} Componentes como lâmpadas de filamento, díodos ou termístores podem apresentar relações não lineares.

\subsection*{Resistência de condutores filiformes e resistividade}
Para um fio homogéneo, a resistência relaciona-se com o material e a geometria por
\[
R=\rho\frac{L}{A},
\]
onde $\rho$ é a \textbf{resistividade}, $L$ o comprimento e $A$ a área da secção transversal. A unidade SI da resistividade é \si{\ohm.m}. Para um mesmo material e temperatura, $R$ cresce linearmente com $L$ e diminui quando a área da secção aumenta.

\begin{figure}[H]
\centering
\begin{minipage}{.48\textwidth}
\centering
\begin{tikzpicture}
\begin{axis}[width=\linewidth,height=5.2cm,xlabel={$L$/m},ylabel={$R$/\si{\ohm}},xmin=0,xmax=10,ymin=0,ymax=5.5,grid=major,legend style={at={(.03,.97)},anchor=north west}]
\addplot[thick,eublue,domain=0:10]{0.25*x};
\addplot[thick,euorange,domain=0:10]{0.48*x};
\legend{material X,material Y}
\end{axis}
\end{tikzpicture}
\end{minipage}\hfill
\begin{minipage}{.48\textwidth}
\centering
\begin{tikzpicture}
\begin{axis}[width=\linewidth,height=5.2cm,xlabel={$1/A$ / unidade arbitrária},ylabel={$R$/\si{\ohm}},xmin=0,xmax=5,ymin=0,ymax=5.5,grid=major]
\addplot[thick,eugreen,domain=0:5]{x};
\end{axis}
\end{tikzpicture}
\end{minipage}
\caption{Para temperatura e material fixos, $R\propto L$ e $R\propto1/A$. Em $R(L)$, o declive é $\rho/A$.}
\end{figure}

A resistividade também pode depender da temperatura. Em metais, aumenta geralmente com a temperatura no intervalo de utilização comum; em semicondutores, o comportamento pode ser oposto e é fortemente dependente do material e da temperatura. Ligas como constantan ou manganina são escolhidas em certas aplicações por apresentarem variações relativamente pequenas de resistência em intervalos específicos.

\begin{mundo}
Cabos de potência precisam de baixa resistência para reduzir aquecimento por efeito Joule; por isso, materiais bons condutores e secções maiores são usados quando a corrente é elevada. Em sensores, pelo contrário, pode ser útil escolher materiais cuja resistência varie de forma previsível com a temperatura.
\end{mundo}

\begin{exemplo}[title=Exemplo resolvido -- resistência de um fio de cobre]
Um fio de cobre tem $L=\SI{20}{m}$, área $A=\SI{1.5}{mm^2}=1{,}5\times10^{-6}\,\mathrm{m^2}$ e $\rho=1{,}67\times10^{-8}\,\si{\ohm.m}$. Então,
\[
R=\rho\frac{L}{A}
=1{,}67\times10^{-8}\frac{20}{1{,}5\times10^{-6}}
\approx\SI{0.223}{\ohm}.
\]
Se a corrente máxima admissível for \SI{14.5}{A}, a queda de tensão no fio será
\[
U=RI\approx0{,}223\times14{,}5\approx\SI{3.23}{V}.
\]
\end{exemplo}

\section{Associações de resistências}

\begin{figure}[H]
\centering
\begin{circuitikz}[scale=.85,transform shape]
\draw (0,0) to[battery1,l=$U$] (0,2.5) to[R,l=$R_1$] (3,2.5) to[R,l=$R_2$] (6,2.5) -- (6,0)--(0,0);
\node[font=\bfseries] at (3,-.5) {série};
\begin{scope}[xshift=8cm]
\draw (0,0) to[battery1,l=$U$] (0,3) -- (1.2,3) -- (5,3) -- (5,0) -- (0,0);
\draw (1.5,3) to[R,l=$R_1$] (1.5,0);
\draw (4.0,3) to[R,l=$R_2$] (4.0,0);
\node[font=\bfseries] at (2.5,-.5) {paralelo};
\end{scope}
\end{circuitikz}
\caption{Duas associações fundamentais de resistências.}
\end{figure}

Em série, a corrente é a mesma nos componentes e
\[
R_{eq}=R_1+R_2+\cdots.
\]
Em paralelo, a diferença de potencial é a mesma nos ramos e
\[
\frac1{R_{eq}}=\frac1{R_1}+\frac1{R_2}+\cdots.
\]

\begin{exemplo}[title=Exemplo resolvido -- série e paralelo]
Dois resistores de \SI{6}{\ohm} e \SI{3}{\ohm}:
\[
R_{\text{série}}=6+3=\SI{9}{\ohm},
\]
enquanto em paralelo
\[
\frac{1}{R_{\rm p}}=\frac16+\frac13=\frac12
\Rightarrow R_{\rm p}=\SI{2}{\ohm}.
\]
A resistência equivalente em paralelo é inferior à menor resistência do conjunto.
\end{exemplo}

\section{Potência e efeito Joule}
A potência elétrica é
\[
P=UI.
\]
Para um resistor óhmico,
\[
P=RI^2=\frac{U^2}{R}.
\]
A energia transferida num intervalo $\Delta t$ é
\[
E=P\Delta t.
\]

\begin{exemplo}
Um aquecedor de \SI{1500}{W} funciona durante \SI{20}{min}. A energia transferida é
\[
E=P\Delta t=1500\times(20\times60)=1{,}8\times10^6\ \mathrm{J}=\SI{1.8}{MJ}.
\]
Em quilowatt-hora, $E=\SI{0.50}{kWh}$.
\end{exemplo}

\section{Geradores reais}
Uma pilha não mantém necessariamente uma tensão terminal constante quando a corrente aumenta. Um modelo simples usa força eletromotriz $\varepsilon$ e resistência interna $r$:
\[
U=\varepsilon-rI.
\]
Esta relação permite obter $\varepsilon$ a partir da ordenada na origem e $r$ a partir do módulo do declive num gráfico $U(I)$.

\begin{center}
\begin{tikzpicture}
\begin{axis}[
width=9cm,height=5.5cm,
xlabel={$I$/A},ylabel={$U$/V},
xmin=0,xmax=1.6,ymin=0,ymax=1.8,
grid=major]
\addplot[domain=0:1.5,thick]{1.55-0.55*x};
\addplot[only marks] coordinates{(0.2,1.44)(0.5,1.28)(0.8,1.11)(1.1,.95)(1.35,.81)};
\node at (axis cs:.45,.45) {$U=\varepsilon-rI$};
\end{axis}
\end{tikzpicture}
\end{center}

A força eletromotriz representa energia fornecida pelo gerador por unidade de carga:
\[
\varepsilon=\frac{E_{\rm gerada}}{Q}.
\]
A potência produzida pelo gerador é $\varepsilon I$; a potência útil entregue ao circuito externo é $UI$; a diferença $rI^2$ corresponde à dissipação interna no modelo resistivo.

\begin{exemplo}[title=Exemplo resolvido -- gerador real]
Uma pilha tem $\varepsilon=\SI{1.50}{V}$ e $r=\SI{0.40}{\ohm}$. Se fornece \SI{0.50}{A},
\[
U=1{,}50-0{,}40\times0{,}50=\SI{1.30}{V}.
\]
A potência útil é $UI=\SI{0.65}{W}$ e a potência dissipada internamente é
\[
rI^2=0{,}40(0{,}50)^2=\SI{0.10}{W}.
\]
\end{exemplo}

\begin{atividade}[title=AL 2.1 -- Características de uma pilha]
Monta-se um circuito com uma pilha, amperímetro e resistência externa variável; o voltímetro mede a tensão nos terminais da pilha. Para vários valores da resistência externa registam-se pares $(I,U)$.

Representa-se $U$ em função de $I$ e ajusta-se uma reta:
\[
U=\varepsilon-rI.
\]
A ordenada na origem fornece $\varepsilon$ e o módulo do declive fornece $r$. Deve evitar-se corrente excessiva e nunca se deve curto-circuitar a pilha.
\end{atividade}

\begin{atividade}[title=Extensão -- característica de uma pilha]
\textbf{Característica de uma pilha.} Medir pares $(I,U)$ para diferentes resistências externas adequadas. Construir o gráfico $U$ em função de $I$. Ajustar uma reta e interpretar os parâmetros. Não curto-circuitar a pilha.
\end{atividade}

\begin{mundo}
A eletrificação permite iluminação, comunicações, mobilidade, refrigeração e medicina tecnológica, mas a produção, transporte, armazenamento e consumo de energia têm custos económicos e ambientais. Uma análise STEM completa deve considerar simultaneamente potência, eficiência, materiais, origem da energia e ciclo de vida dos equipamentos.

A potência indicada num eletrodoméstico ajuda a estimar o consumo energético, mas o custo depende também do tempo de utilização. A eficiência energética liga Física, economia doméstica e sustentabilidade.
\end{mundo}


% ---------- Enriquecimento a partir das fichas Fis_FT07--Fis_FT12 ----------
\section{Carga elétrica, corrente e leitura de circuitos}
Nos metais, os eletrões de condução são os principais portadores de carga. O \textbf{sentido convencional da corrente} é definido como o sentido do movimento de cargas positivas e, por isso, é oposto ao movimento médio dos eletrões num fio metálico.

A intensidade da corrente mede a carga que atravessa uma secção do circuito por unidade de tempo:
\[
I=\frac{|\Delta q|}{\Delta t}.
\]
A diferença de potencial elétrico mede a energia transferida por unidade de carga:
\[
U=\frac{\Delta E}{|\Delta q|}.
\]

\begin{figure}[H]
\centering
\begin{tikzpicture}[>=Stealth]
\draw[rounded corners=2mm,fill=eugrey!45,draw=eudark] (0,0) rectangle (11,1.3);
\foreach \x in {1,2.2,3.4,4.6,5.8,7,8.2,9.4,10.2}{\fill[eublue] (\x,.65) circle (2.4pt);}
\draw[->,very thick,eublue] (8.8,1.7)--(5.7,1.7) node[midway,above,font=\small]{movimento médio dos eletrões};
\draw[->,very thick,euorange] (2.2,-.4)--(5.3,-.4) node[midway,below,font=\small]{corrente convencional $I$};
\node[font=\small] at (5.5,.65) {condutor metálico};
\end{tikzpicture}
\caption{Sentido convencional da corrente e movimento médio dos eletrões num metal.}
\end{figure}

\begin{nota}
O amperímetro liga-se em \textbf{série} com o ramo cuja corrente se pretende medir. O voltímetro liga-se em \textbf{paralelo} aos terminais do componente cuja diferença de potencial se pretende medir.
\end{nota}

\section{Curvas características: do gráfico à lei física}
Para um condutor óhmico a temperatura aproximadamente constante,
\[
U=RI.
\]
Se se representar $U$ em função de $I$, obtém-se uma reta que passa aproximadamente na origem e cujo declive é a resistência $R$. Um filamento de lâmpada pode apresentar uma curva não linear porque a sua temperatura e, consequentemente, a sua resistência variam com a corrente.

\begin{figure}[H]
\centering
\begin{tikzpicture}
\begin{axis}[width=11cm,height=6cm,xlabel={$I$ / A},ylabel={$U$ / V},xmin=0,xmax=2.2,ymin=0,ymax=12.5,grid=major,legend pos=north west]
\addplot[domain=0:2,samples=2,very thick,eublue]{5*x};
\addlegendentry{resistor óhmico, $R=\SI{5.0}{\ohm}$}
\addplot[very thick,euorange,smooth] coordinates{(0,0)(.3,.65)(.6,1.55)(.9,2.8)(1.2,4.5)(1.5,6.8)(1.8,9.7)(2.0,12)};
\addlegendentry{filamento de lâmpada}
\end{axis}
\end{tikzpicture}
\caption{Curvas características $U(I)$ de dois recetores.}
\end{figure}

\begin{exemplo}[title=Exemplo resolvido -- resistência a partir de uma curva $U(I)$]
Uma reta experimental passa pelos pontos $(\SI{0.40}{A},\SI{2.4}{V})$ e $(\SI{1.20}{A},\SI{7.2}{V})$. O declive é
\[
R=\frac{\Delta U}{\Delta I}=\frac{7{,}2-2{,}4}{1{,}20-0{,}40}=\SI{6.0}{\ohm}.
\]
A linearidade e a passagem próxima da origem são compatíveis com um condutor óhmico nas condições do ensaio.
\end{exemplo}

\section{Resistividade: material, comprimento, área e temperatura}
As fichas de resistividade distinguem a \textbf{resistência}, propriedade do condutor concreto, da \textbf{resistividade} $\rho$, propriedade característica do material para determinadas condições. Para um fio uniforme,
\[
R=\rho\frac{L}{A}.
\]
Assim, a resistência aumenta com o comprimento e diminui quando a área da secção aumenta.

\begin{figure}[H]
\centering
\begin{minipage}{0.48\textwidth}
\centering
\begin{tikzpicture}
\begin{axis}[width=\linewidth,height=5.2cm,xlabel={$L$ / m},ylabel={$R$ / $\Omega$},xmin=0,xmax=5,ymin=0,ymax=5.5,grid=major]
\addplot[domain=0:5,very thick,eublue]{1.0*x};
\node[font=\scriptsize,anchor=west] at (axis cs:2.2,3.4) {declive $=\rho/A$};
\end{axis}
\end{tikzpicture}
\small Área e material constantes.
\end{minipage}\hfill
\begin{minipage}{0.48\textwidth}
\centering
\begin{tikzpicture}
\begin{axis}[width=\linewidth,height=5.2cm,xlabel={$1/A$ / u.a.},ylabel={$R$ / $\Omega$},xmin=0,xmax=5,ymin=0,ymax=5.5,grid=major]
\addplot[domain=0:5,very thick,eugreen]{.8*x};
\node[font=\scriptsize,anchor=west] at (axis cs:2.1,3.2) {declive $=\rho L$};
\end{axis}
\end{tikzpicture}
\small Comprimento e material constantes.
\end{minipage}
\caption{Linearizações úteis da expressão $R=\rho L/A$.}
\end{figure}

Em metais, a resistividade tende a aumentar com a temperatura. Em semicondutores, o comportamento pode ser diferente, porque o número e a mobilidade dos portadores de carga dependem fortemente da temperatura.

\begin{mundo}
\textbf{Porque têm os cabos de potência secções grandes?} Para o mesmo material e comprimento, aumentar a área reduz a resistência. Isso reduz a potência dissipada por efeito Joule, $P=RI^2$, e o aquecimento do cabo para a mesma corrente.
\end{mundo}

\section{Efeito Joule e gráficos de potência}
Num recetor puramente resistivo,
\[
P=UI=RI^2=\frac{U^2}{R}.
\]
A forma mais adequada depende das grandezas conhecidas. Se $R$ for constante, um gráfico $P=f(I^2)$ é linear e o declive é $R$.

\begin{figure}[H]
\centering
\begin{tikzpicture}
\begin{axis}[width=10.8cm,height=5.8cm,xlabel={$I^2$ / A$^2$},ylabel={$P$ / W},xmin=0,xmax=9,ymin=0,ymax=38,grid=major]
\addplot[domain=0:9,very thick,eured]{4*x};
\node[font=\small,anchor=west] at (axis cs:4.4,23) {declive $=R=\SI{4}{\ohm}$};
\end{axis}
\end{tikzpicture}
\caption{Para um resistor óhmico, $P$ é diretamente proporcional a $I^2$.}
\end{figure}

\begin{exemplo}[title=Exemplo resolvido -- potência dissipada]
Um resistor de \SI{8.0}{\ohm} é atravessado por uma corrente de \SI{1.5}{A}. A potência transformada em energia interna é
\[
P=RI^2=8{,}0\times1{,}5^2=\SI{18}{W}.
\]
Em \SI{5.0}{min}, a energia transferida é
\[
E=P\Delta t=18\times300=\SI{5.4e3}{J}.
\]
\end{exemplo}

\section{Geradores reais: interpretar a reta $U(I)$}
Num gerador real, a tensão nos terminais diminui com a corrente devido à resistência interna $r$:
\[
U=\varepsilon-rI.
\]
A ordenada na origem é a força eletromotriz $\varepsilon$ e o módulo do declive é a resistência interna $r$.

\begin{figure}[H]
\centering
\begin{tikzpicture}
\begin{axis}[width=11cm,height=6cm,xlabel={$I$ / A},ylabel={$U$ / V},xmin=0,xmax=4.5,ymin=0,ymax=6.6,grid=major]
\addplot[domain=0:4,samples=2,very thick,eublue]{6-1.2*x};
\addplot[dashed,gray] coordinates{(0,6)(4,6)};
\fill[eublue] (axis cs:0,6) circle (2.5pt);
\node[font=\small,anchor=west] at (axis cs:.1,6.1) {$\varepsilon=\SI{6.0}{V}$};
\node[font=\small,anchor=west] at (axis cs:2.2,2.4) {declive $=-r$};
\end{axis}
\end{tikzpicture}
\caption{Curva característica idealizada de um gerador real.}
\end{figure}

\begin{exemplo}[title=Exemplo resolvido -- força eletromotriz e resistência interna]
A reta de um gerador passa pelos pontos $(\SI{0}{A},\SI{9.0}{V})$ e $(\SI{2.0}{A},\SI{7.4}{V})$. Logo, $\varepsilon=\SI{9.0}{V}$ e
\[
r=\left|\frac{7{,}4-9{,}0}{2{,}0-0}\right|=\SI{0.80}{\ohm}.
\]
Para $I=\SI{1.5}{A}$, a tensão útil será
\[
U=9{,}0-0{,}80\times1{,}5=\SI{7.8}{V}.
\]
O rendimento elétrico do gerador, nestas condições, é $\eta=U/\varepsilon\approx87\%$.
\end{exemplo}

\begin{stem}
\textbf{Dimensionar um cabo.} Compara dois fios de cobre com o mesmo comprimento mas secções diferentes. Calcula a resistência e a potência dissipada para uma corrente escolhida. Depois propõe critérios de engenharia que também teriam de ser considerados num projeto real: custo, massa, aquecimento admissível, isolamento e segurança.
\end{stem}

\begin{exercicios}
\begin{enumerate}
  \item Uma corrente de \SI{0.40}{A} circula durante \SI{2.0}{min}. Calcula a carga que atravessa uma secção do condutor.
  \item Um resistor de \SI{12}{\ohm} é ligado a \SI{6.0}{V}. Determina a corrente e a potência.
  \item Calcula a resistência equivalente de \SI{10}{\ohm} e \SI{20}{\ohm} em série e em paralelo.
  \item Uma lâmpada de \SI{9}{W} funciona \SI{5}{h}. Calcula a energia em Wh e em J.
  \item Para uma pilha, o ajuste experimental é $U=1{,}55-0{,}80I$ em unidades SI. Identifica $\varepsilon$ e $r$.
  \item Explica por que razão os fios de uma instalação podem aquecer quando a corrente é elevada.
  \item \textbf{Integração.} Um painel fornece \SI{18}{V} e \SI{2.5}{A} durante \SI{3.0}{h}. Estima a energia elétrica produzida em Wh e J.
  \item Uma pilha tem $\varepsilon=\SI{9.0}{V}$ e $r=\SI{1.0}{\ohm}$. Determina a tensão terminal quando fornece \SI{0.50}{A}.
  \item Dois resistores idênticos de \SI{12}{\ohm} são ligados em paralelo a uma fonte de \SI{6.0}{V}. Determina a corrente em cada ramo e a corrente total.
  \item Um fio tem o dobro do comprimento e a mesma área de secção de outro fio do mesmo material. Compara as resistências.
  \item Indica os principais portadores de carga numa corrente elétrica em (a) um metal; (b) uma solução de cloreto de sódio; (c) um gás ionizado.
  \item Um fio de cobre tem $L=\SI{20}{m}$, $A=\SI{1.5}{mm^2}$ e $\rho=1{,}67\times10^{-8}\,\si{\ohm.m}$. Determina a resistência e a tensão nos seus terminais quando $I=\SI{14.5}{A}$.
  \item Uma lâmpada recebe \SI{216}{J} em \SI{60}{s}, sendo percorrida por uma corrente de \SI{0.300}{A}. Determina a potência, a tensão nos terminais e a resistência instantânea correspondente.
  \item Por uma secção de um fio passam $5{,}0\times10^{18}$ eletrões por minuto. Estima o módulo da corrente elétrica, usando $e=1{,}602\times10^{-19}\,\mathrm{C}$.
  \item Dois fios X e Y têm a mesma área e o mesmo comprimento. No gráfico $R(L)$, Y teria maior declive. O que se pode concluir sobre as resistividades de X e Y?
\end{enumerate}
\end{exercicios}

\chapter{Energia, fenómenos térmicos e radiação}

\begin{objetivos}
\begin{itemize}
  \item distinguir temperatura, energia interna e calor e reconhecer a condição de equilíbrio térmico;
  \item relacionar qualitativamente temperatura com a energia cinética média das partículas;
  \item reconhecer condução, convecção e radiação;
  \item aplicar a Primeira Lei da Termodinâmica em situações simples;
  \item usar $Q=mc\Delta T$ e interpretar capacidade térmica mássica;
  \item compreender mudanças de estado e entalpia de fusão;
  \item avaliar rendimento e degradação de energia;
  \item investigar fatores que afetam o rendimento de painéis fotovoltaicos.
\end{itemize}
\end{objetivos}

\section{Temperatura, equilíbrio térmico, calor e energia interna}
Um \textbf{sistema termodinâmico} é um sistema em que a energia interna é relevante para a descrição. A fronteira pode permitir ou impedir, em diferentes graus, transferências de matéria e energia. Na prática, uma garrafa térmica não é perfeitamente isolada, mas pode ser modelada como aproximadamente isolada durante um intervalo limitado.

Temperatura não é ``quantidade de calor''. A \textbf{temperatura} é a propriedade que permite prever o sentido da transferência espontânea de energia por calor entre dois sistemas em contacto térmico. Quando dois sistemas atingem a mesma temperatura e deixa de existir transferência líquida de energia por calor entre eles, encontram-se em \textbf{equilíbrio térmico}.

Num modelo microscópico simples, o aumento da temperatura corresponde, em geral, a um aumento da \textbf{energia cinética média das partículas}. Isto não significa que todas as partículas tenham a mesma energia nem que a energia interna seja apenas energia cinética: a energia interna inclui diferentes contribuições microscópicas, dependentes do sistema.

\textbf{Calor} é energia transferida espontaneamente devido a uma diferença de temperatura; não é uma substância ``contida'' num corpo. A \textbf{energia interna} $U$ é uma propriedade do sistema associada às energias microscópicas das suas partículas.

\begin{exemplo}[title=Exemplo resolvido -- equilíbrio térmico]
Um bloco metálico a \SI{80}{\celsius} é colocado em água a \SI{20}{\celsius} num recipiente termicamente isolado. Antes do equilíbrio, a energia é transferida espontaneamente do bloco para a água. No equilíbrio, ambos têm a mesma temperatura final $T_f$, compreendida entre \SI{20}{\celsius} e \SI{80}{\celsius}, se não ocorrer mudança de estado. A energia total do sistema isolado conserva-se, mas as energias internas de cada subsistema variam.
\end{exemplo}

\subsection*{Lei Zero e escalas de temperatura}
A \textbf{Lei Zero da Termodinâmica} estabelece que, se A está em equilíbrio térmico com C e B também está em equilíbrio térmico com C, então A e B estão em equilíbrio térmico entre si. Esta ideia fundamenta o uso de termómetros: o termómetro atinge equilíbrio térmico com o sistema e a sua propriedade termométrica permite atribuir uma temperatura.

A unidade SI de temperatura termodinâmica é o kelvin (K). Para converter temperaturas Celsius,
\[
T(\mathrm K)=\theta(^{\circ}\mathrm C)+273{,}15.
\]
Uma \textbf{diferença} de \SI{1}{\celsius} tem o mesmo tamanho que uma diferença de \SI{1}{K}; o que muda entre as escalas é a origem.

\begin{figure}[H]
\centering
\begin{tikzpicture}[y=.035cm]
\draw[very thick,eublue] (0,0)--(0,380);
\draw[very thick,euorange] (5,0)--(5,380);
\foreach \y/\c/\k in {0/-273{,}15/0,273.15/0/273{,}15,373.15/100/373{,}15}{
  \draw (-.18,\y)--(.18,\y); \node[left] at (-.25,\y) {$\c\,^{\circ}\mathrm C$};
  \draw (4.82,\y)--(5.18,\y); \node[right] at (5.25,\y) {$\k\,\mathrm K$};
}
\node[above,font=\bfseries,eublue] at (0,385) {Celsius};
\node[above,font=\bfseries,euorange] at (5,385) {kelvin};
\draw[<->,dashed] (.4,273.15)--(4.6,273.15);
\draw[<->,dashed] (.4,373.15)--(4.6,373.15);
\end{tikzpicture}
\caption{Correspondência entre alguns pontos das escalas Celsius e kelvin.}
\end{figure}

\begin{exemplo}[title=Exemplo resolvido -- temperatura e variação de temperatura]
Uma amostra passa de \SI{18.0}{\celsius} para \SI{43.0}{\celsius}. As temperaturas absolutas são $291{,}15$ K e $316{,}15$ K, mas
\[
\Delta T=316{,}15-291{,}15=\SI{25.0}{K}.
\]
A mesma variação é \SI{25.0}{\celsius} quando expressa como diferença de temperaturas.
\end{exemplo}

\section{Mecanismos de transferência}
A energia pode atravessar a fronteira de um sistema por diferentes mecanismos. Em situações térmicas quotidianas, condução, convecção e radiação podem ocorrer simultaneamente, embora uma delas possa dominar.

\begin{figure}[H]
\centering
\begin{tikzpicture}[>=Stealth,scale=.92]
% conduction
\draw[fill=eugrey] (0,0) rectangle (3.2,.45); \fill[eured] (.25,.225) circle (5pt);
\draw[->,thick,eured] (.55,.225)--(2.75,.225); \node at (1.6,-.45) {condução};
% convection vessel
\draw[thick] (4.7,-.1)--(4.7,2.0)--(7.5,2.0)--(7.5,-.1);
\fill[eucyan!12] (4.72,0) rectangle (7.48,1.65);
\draw[->,thick,euorange] (5.35,.25)..controls(5.05,1.0)and(5.4,1.55)..(6.05,1.5);
\draw[->,thick,eublue] (6.85,1.5)..controls(7.25,.9)and(7.0,.3)..(6.5,.2);
\node at (6.1,-.45) {convecção};
% radiation
\fill[eured] (9.5,.9) circle (.35);
\foreach \r in {.65,.95,1.25}{\draw[euorange,thick] (9.5,0.9) ++(-30:\r) arc (-30:30:\r);}
\draw[fill=eublue!15] (12.0,.1) rectangle (12.4,1.7);
\node at (10.9,-.45) {radiação};
\end{tikzpicture}
\caption{Representações esquemáticas de condução, convecção e radiação.}
\end{figure}
\begin{description}
  \item[Condução:] transferência associada a interações microscópicas no material, sem transporte macroscópico de matéria.
  \item[Convecção:] transferência associada ao movimento de fluidos.
  \item[Radiação:] transferência por ondas eletromagnéticas; não necessita de meio material.
\end{description}

\begin{center}
\begin{tikzpicture}[scale=1]
\draw[fill=euorange!15] (0,0) rectangle (2.1,2.2);
\node at (1.05,1.1) {corpo quente};
\draw[-{Stealth},very thick,eured] (2.2,1.7)--(4.2,1.7) node[midway,above] {radiação};
\draw[-{Stealth},very thick,eublue] (2.2,1.0)--(4.2,1.0) node[midway,above] {condução};
\draw[-{Stealth},very thick,eugreen] (2.2,.3)--(4.2,.3) node[midway,below] {convecção};
\draw[fill=eublue!8] (4.3,0) rectangle (6.4,2.2);
\node at (5.35,1.1) {meio envolvente};
\end{tikzpicture}
\end{center}

\subsection*{Radiação térmica e infravermelho}
Todos os corpos com temperatura acima do zero absoluto emitem radiação eletromagnética. À temperatura ambiente, a emissão de objetos comuns ocorre predominantemente no \textbf{infravermelho}. Câmaras térmicas convertem diferenças de radiação infravermelha em imagens úteis na inspeção de edifícios, medicina, indústria e combate a incêndios.

\begin{sabias}
Um objeto pode emitir e absorver radiação simultaneamente. O seu estado térmico depende do balanço energético líquido, não apenas do facto de estar a emitir radiação.
\end{sabias}


\subsection*{Absorção, emissão e propriedades das superfícies}
A fração de radiação absorvida depende do material, do acabamento superficial e do comprimento de onda. Em equilíbrio térmico, bons absorvedores numa determinada faixa espetral tendem também a ser bons emissores nessa faixa. A \textbf{emissividade} é uma grandeza adimensional entre 0 e 1 que compara a emissão térmica de uma superfície com a de um corpo negro ideal à mesma temperatura.

\begin{nota}
A cor observada no visível, por si só, \textbf{não determina} a emissividade no infravermelho. Uma superfície branca pode, dependendo do material e do acabamento, emitir infravermelho com elevada emissividade. Em aplicações térmicas reais interessa conhecer propriedades espetrais, não apenas a cor visível.
\end{nota}

\begin{figure}[H]
\centering
\begin{tikzpicture}[>=Stealth]
\draw[fill=eudark!90] (0,0) rectangle (3.5,2.4); \node[white] at (1.75,1.2) {superfície A};
\draw[fill=eugrey!45] (8,0) rectangle (11.5,2.4); \node at (9.75,1.2) {superfície B};
\foreach \y in {.45,1.05,1.65}{\draw[->,very thick,euorange] (3.7,\y)--(5.4,\y);}
\foreach \y in {.65,1.25}{\draw[->,thick,eucyan] (7.8,\y)--(6.5,\y);}
\node[align=center] at (5.75,2.85) {emissão e absorção dependem\\do material e de $\lambda$};
\end{tikzpicture}
\caption{Representação qualitativa: emissão e absorção radiativas são propriedades do sistema superfície--radiação.}
\end{figure}

\section{Primeira Lei da Termodinâmica}
Adotando a convenção de que $W$ é o trabalho realizado \textit{sobre} o sistema,
\[
\Delta U=Q+W.
\]
É essencial declarar a convenção usada, pois alguns manuais usam o sinal oposto para o trabalho.

A formulação da conservação da energia em termodinâmica foi consolidada no século XIX através de resultados experimentais sobre equivalência entre trabalho mecânico e aquecimento, associados, entre outros, a Joule. Esta evolução ilustra como medições quantitativas podem ligar fenómenos que inicialmente pareciam distintos.

\begin{exemplo}[title=Exemplo resolvido -- Primeira Lei]
Um gás recebe \SI{300}{J} por calor e realiza \SI{80}{J} de trabalho sobre a vizinhança. Na convenção deste livro, o trabalho \emph{sobre} o sistema é $W=-\SI{80}{J}$:
\[
\Delta U=300-80=\SI{220}{J}.
\]
\end{exemplo}

\section{Aquecimento e capacidade térmica mássica}
Sem mudança de estado, uma aproximação comum é
\[
Q=mc\Delta T.
\]
Materiais com grande $c$ requerem mais energia, por unidade de massa, para a mesma variação de temperatura.

\begin{exemplo}
Que energia é necessária para aquecer \SI{0.50}{kg} de água de \SI{20}{\celsius} para \SI{60}{\celsius}, usando $c=\SI{4180}{J/(kg.K)}$?
\[
Q=0{,}50\times4180\times40\approx\SI{8.36e4}{J}.
\]
\end{exemplo}

\subsection*{Capacidade térmica e capacidade térmica mássica}
A capacidade térmica de um corpo é $C=mc$ e indica a energia necessária para aumentar a sua temperatura em um kelvin:
\[
Q=C\Delta T=mc\Delta T.
\]
A capacidade térmica mássica $c$ é uma propriedade do material, enquanto $C$ depende também da massa do corpo.

\subsection*{Sinal da variação da energia interna}
Num aquecimento sem trabalho mecânico relevante, $\Delta T>0$ e, num modelo simples, $\Delta U>0$. Num arrefecimento, $\Delta T<0$ e $\Delta U<0$. Durante uma transição de fase, a temperatura pode permanecer constante e, ainda assim, a energia interna variar. Por isso, a expressão $mc\Delta T$ não descreve, por si só, todas as situações térmicas.

\begin{exemplo}[title=Exemplo resolvido -- aquecimento e sinal de $\Delta U$]
Uma amostra de alumínio de \SI{0.250}{kg} aquece de \SI{20}{\celsius} para \SI{50}{\celsius}. Admitindo $c=\SI{900}{J/(kg.K)}$ e trabalho mecânico desprezável,
\[
\Delta U\approx mc\Delta T=0{,}250\times900\times30=\SI{6.75e3}{J}.
\]
Como a temperatura aumentou, $\Delta U$ é positiva nesta aproximação.
\end{exemplo}

\section{Mudanças de estado}
Durante uma mudança de estado a temperatura pode manter-se aproximadamente constante enquanto há transferência de energia. Na fusão,
\[
Q=mL_f,
\]
onde $L_f$ é a entalpia específica de fusão (muitas vezes designada calor latente específico de fusão em linguagem escolar).

\begin{figure}[H]
\centering
\begin{tikzpicture}
\begin{axis}[width=11cm,height=6.2cm,xlabel={energia fornecida / unidade arbitrária},ylabel={temperatura / $^\circ$C},xmin=0,xmax=10,ymin=-25,ymax=115,grid=major]
\addplot[thick,eublue] coordinates{(0,-20)(2,0)(4,0)(7,100)(9,100)(10,112)};
\node[font=\small,rotate=25] at (axis cs:1,-8) {sólido};
\node[font=\small] at (axis cs:3,8) {fusão};
\node[font=\small,rotate=24] at (axis cs:5.7,58) {líquido};
\node[font=\small] at (axis cs:8,107) {ebulição};
\node[font=\small,rotate=30] at (axis cs:9.55,107) {gás};
\end{axis}
\end{tikzpicture}
\caption{Curva de aquecimento idealizada para uma substância pura, mostrando patamares de mudança de estado.}
\end{figure}
Durante a transição de fase de uma substância pura, a energia fornecida é usada sobretudo na reorganização microscópica das partículas; por isso, a temperatura pode manter-se aproximadamente constante.

\section{Rendimento}
O rendimento energético pode ser expresso por
\[
\eta=\frac{E_{\text{útil}}}{E_{\rm fornecida}}\times100\%.
\]
Nenhum processo real transforma toda a energia fornecida na forma desejada. A energia conserva-se, mas a sua capacidade de realizar determinadas transformações úteis degrada-se.

\section{Irradiância e painéis fotovoltaicos}
A \textbf{irradiância} é a potência radiante incidente por unidade de área:
\[
G=\frac{P_{\rm inc}}{A}=\frac{E_{\rm inc}}{A\,\Delta t}.
\]
Exprime-se em \si{W/m^2}. Não deve ser confundida com energia: uma irradiância elevada mantida durante pouco tempo pode fornecer menos energia total do que uma irradiância menor mantida durante muitas horas.

\begin{figure}[H]
\centering
\begin{tikzpicture}[>=Stealth]
\fill[euorange] (0,1.5) circle (.42); \node[below=5pt] at (0,1.08) {fonte};
\draw[fill=eublue!15,draw=eublue] (7,.25) rectangle (7.5,2.75);
\foreach \y in {.5,.9,1.3,1.7,2.1,2.5}{\draw[->,euorange!90,thick] (.6,1.5)--(6.8,\y);}
\draw[decorate,decoration={brace,amplitude=5pt},eublue] (7.8,.25)--(7.8,2.75) node[midway,right=7pt] {área $A$};
\node at (4.1,3.05) {potência radiante incidente $P_{\rm inc}$};
\end{tikzpicture}
\caption{A irradiância relaciona a potência radiante recebida com a área exposta.}
\end{figure}

A potência radiante incidente numa área $A$ pode relacionar-se com a irradiância $G$:
\[
P_{\rm inc}=GA.
\]
Se a potência elétrica útil é $P_{el}$,
\[
\eta=\frac{P_{el}}{GA}.
\]

\begin{figure}[H]
\centering
\begin{tikzpicture}
\begin{axis}[width=11cm,height=6cm,xlabel={$U$ / V},ylabel={potência elétrica $P$ / W},xmin=0,xmax=7,ymin=0,ymax=4.2,grid=major]
\addplot[thick,eublue,smooth] coordinates{(0,0)(1,1.0)(2,1.9)(3,2.75)(4,3.35)(4.8,3.55)(5.5,3.2)(6.0,2.2)(6.4,.7)(6.5,0)};
\fill[euorange] (axis cs:4.8,3.55) circle (3pt);
\node[anchor=west,font=\small] at (axis cs:4.9,3.55) {ponto de potência máxima};
\end{axis}
\end{tikzpicture}
\caption{Curva de potência ilustrativa de um pequeno painel fotovoltaico para uma irradiância fixa.}
\end{figure}


\begin{exemplo}[title=Exemplo resolvido -- energia solar recebida e aquecimento]
Uma lata expõe uma área de $\SI{140}{cm^2}=\SI{0.0140}{m^2}$ a uma irradiância média de \SI{600}{W/m^2} durante \SI{90}{min}. A energia solar incidente é
\[
E_{\rm inc}=GA\Delta t=600\times0{,}0140\times(90\times60)
\approx\SI{4.54e4}{J}.
\]
Se \SI{0.340}{kg} de bebida, com $c=\SI{4.2e3}{J/(kg.K)}$, aumentar \SI{16.5}{K}, recebe aproximadamente
\[
Q=mc\Delta T=0{,}340\times4{,}2\times10^3\times16{,}5
\approx\SI{2.36e4}{J}.
\]
A fração da energia incidente associada ao aquecimento da bebida é, neste modelo,
\[
\frac{Q}{E_{\rm inc}}\times100\approx\SI{52}{\percent}.
\]
As restantes transferências incluem aquecimento da lata e perdas para a vizinhança.
\end{exemplo}

\begin{atividade}[title=AL 3.1 -- Radiação e potência elétrica de um painel fotovoltaico]
Mede-se $U$ e $I$ para diferentes cargas e calcula-se
\[
P_{\rm el}=UI.
\]
Repete-se o estudo para diferentes irradiâncias, ou diferentes distâncias/ângulos relativamente à fonte, mantendo controladas as restantes variáveis. O objetivo é identificar as condições que maximizam a potência elétrica fornecida.

A irradiância é potência radiante incidente por unidade de área, em \si{W/m^2}. Não se deve confundir irradiância com energia total recebida.
\end{atividade}

\begin{atividade}[title=AL 3.2 -- Capacidade térmica mássica de um material]
Um bloco metálico é aquecido eletricamente. Mede-se a energia elétrica fornecida, por exemplo através de $E=UI\Delta t$, e a variação de temperatura. Num modelo simples,
\[
E_{\rm fornecida}\approx mc\Delta T+E_{\rm perdas}.
\]
Quanto melhor for o isolamento e a caracterização das perdas, mais próximo o resultado experimental estará do valor de referência.
\end{atividade}

\begin{atividade}[title=AL 3.3 -- Balanço energético num sistema água-gelo]
Mistura-se uma massa conhecida de gelo próximo de \SI{0}{\celsius} com água a temperatura superior num calorímetro. Num sistema aproximadamente isolado,
\[
E_{\text{cedida pela água}}\approx E_{\text{recebida pelo gelo}}.
\]
A energia recebida pelo gelo pode incluir fusão e aquecimento da água resultante. O balanço permite estimar a entalpia mássica de fusão ou testar um valor conhecido.
\end{atividade}

\begin{exemplo}[title=Exemplo resolvido -- balanço água-gelo]
Suponha que \SI{0.100}{kg} de gelo a \SI{0}{\celsius} funde completamente. Só para a fusão, usando $L_f=\SI{334}{kJ/kg}$,
\[
Q_f=mL_f=0{,}100\times334=\SI{33.4}{kJ}.
\]
Se a água resultante aquecer depois acima de \SI{0}{\celsius}, é necessário acrescentar ainda $mc\Delta T$.
\end{exemplo}

\begin{atividade}[title=Extensão -- painel fotovoltaico]
\textbf{Painel fotovoltaico.} Variar a iluminação, o ângulo de incidência ou a carga elétrica e medir tensão e corrente. Calcular $P=UI$. Representar a potência em função de uma variável escolhida. Identificar variáveis de controlo e discutir limitações experimentais.
\end{atividade}

\section{Segunda Lei, energia útil e sustentabilidade}
A energia total conserva-se, mas os processos espontâneos têm um sentido preferencial. Energia térmica dispersa a baixa diferença de temperatura é, em geral, menos capaz de originar trabalho útil do que energia elétrica ou mecânica organizada. Esta ideia está ligada à \textbf{Segunda Lei da Termodinâmica}.

Não existe máquina térmica ou processo real com 100\% de rendimento útil para todas as finalidades. Melhorar eficiência reduz recursos consumidos e energia degradada, mas a análise sustentável deve considerar também materiais, manutenção, duração e impactos ambientais.

\begin{exemplo}[title=Exemplo resolvido -- eficiência energética]
Um dispositivo recebe \SI{2.0}{MJ} de energia e fornece \SI{1.5}{MJ} na forma útil pretendida:
\[
\eta=\frac{1{,}5}{2{,}0}\times100=\SI{75}{\percent}.
\]
Os restantes \SI{0.5}{MJ} não desapareceram; foram transferidos para formas menos úteis para o objetivo do dispositivo, frequentemente como energia interna do meio.
\end{exemplo}

\begin{stem}
\textbf{Casa energeticamente eficiente.} Compara três soluções para reduzir necessidades de aquecimento: isolamento, vidros duplos e controlo solar. Pesquisa propriedades, custo aproximado, durabilidade e impacto ambiental. Constrói uma matriz de decisão com critérios quantitativos e justifica a solução proposta.
\end{stem}


% ---------- Enriquecimento a partir das fichas Fis_FT13--Fis_FT17 ----------
\section{Modelar sistemas termodinâmicos}
Um sistema termodinâmico é um sistema em que a energia interna é relevante para o estudo. A definição clara da fronteira ajuda a distinguir transferências de energia por calor ou trabalho das transformações que ocorrem no interior do sistema.

\begin{figure}[H]
\centering
\begin{tikzpicture}[>=Stealth]
\node[rounded corners=3mm,draw=eublue,fill=eulightblue,minimum width=4.2cm,minimum height=2.3cm] (s) at (0,0) {\textbf{sistema}};
\draw[dashed,eudark,thick] (-2.5,-1.45) rectangle (2.5,1.45);
\node[font=\small] at (0,1.75) {fronteira};
\node[font=\small] at (5.2,0) {vizinhança};
\draw[<->,very thick,euorange] (2.5,.55)--(4.4,.55) node[midway,above,font=\small] {energia};
\draw[<->,very thick,eugreen] (2.5,-.55)--(4.4,-.55) node[midway,below,font=\small] {matéria};
\end{tikzpicture}
\caption{Sistema, fronteira e vizinhança: a escolha do sistema faz parte do modelo.}
\end{figure}

\section{Temperatura, agitação térmica e escalas}
A temperatura está relacionada com o estado de agitação microscópica das partículas. Para partículas de uma mesma substância e estado físico, uma temperatura maior corresponde, em média, a maior energia cinética microscópica.

As escalas Celsius e kelvin relacionam-se por
\[
T(\mathrm K)=\theta(^{\circ}\mathrm C)+273{,}15.
\]
Uma variação de \SI{1}{K} tem a mesma amplitude que uma variação de \SI{1}{\celsius}; apenas a origem das escalas é diferente.

\begin{figure}[H]
\centering
\begin{tikzpicture}
\draw[very thick,eublue] (0,0)--(0,6.5);
\draw[very thick,euorange] (5,0)--(5,6.5);
\foreach \y/\c/\k in {0/-273/0,2.4/0/273,4.15/100/373,6/200/473}{
  \draw[eublue] (-.15,\y)--(.15,\y) node[left=4pt,font=\small] {\c};
  \draw[euorange] (4.85,\y)--(5.15,\y) node[right=4pt,font=\small] {\k};}
\node[eublue,below,font=\bfseries] at (0,-.3) {$\theta/^{\circ}$C};
\node[euorange,below,font=\bfseries] at (5,-.3) {$T$/K};
\draw[<->,gray,dashed] (.3,2.4)--(4.7,2.4) node[midway,above,font=\small] {fusão da água};
\draw[<->,gray,dashed] (.3,4.15)--(4.7,4.15) node[midway,above,font=\small] {ebulição da água};
\end{tikzpicture}
\caption{Correspondência entre as escalas Celsius e kelvin (valores arredondados).}
\end{figure}

\section{Calor, capacidade térmica e leitura de gráficos}
Na ausência de mudança de estado, a energia transferida como calor pode ser modelada por
\[
Q=mc\Delta T.
\]
Se $m$ e $c$ forem constantes, um gráfico $Q=f(\Delta T)$ é linear, com declive $mc$.

\begin{figure}[H]
\centering
\begin{tikzpicture}
\begin{axis}[width=11cm,height=6cm,xlabel={$\Delta T$ / K},ylabel={$Q$ / kJ},xmin=0,xmax=50,ymin=0,ymax=55,grid=major,legend pos=north west]
\addplot[domain=0:50,very thick,eublue]{1.05*x};
\addlegendentry{amostra A}
\addplot[domain=0:50,very thick,euorange]{0.46*x};
\addlegendentry{amostra B}
\end{axis}
\end{tikzpicture}
\caption{Em $Q(\Delta T)$, um maior declive significa maior capacidade térmica $C=mc$.}
\end{figure}

\begin{exemplo}[title=Exemplo resolvido -- capacidade térmica a partir do gráfico]
Uma amostra recebe \SI{18}{kJ} e a temperatura aumenta \SI{24}{K}. O declive de $Q(\Delta T)$ é
\[
C=\frac{Q}{\Delta T}=\frac{18{,}0\times10^3}{24}=\SI{750}{J/K}.
\]
Se a massa for \SI{2.50}{kg}, então
\[
c=\frac{C}{m}=\frac{750}{2{,}50}=\SI{300}{J/(kg.K)}.
\]
\end{exemplo}

\section{Condução, convecção e radiação: três mecanismos}
A energia pode ser transferida entre sistemas a diferentes temperaturas por diferentes mecanismos.
\begin{itemize}
  \item \textbf{Condução}: transferência microscópica de energia através das interações entre partículas, sem transporte macroscópico de matéria; é particularmente importante em sólidos.
  \item \textbf{Convecção}: transferência associada ao movimento macroscópico de um fluido. Pode ser natural, quando diferenças de massa volúmica originam circulação, ou forçada, quando existe uma ventoinha ou bomba.
  \item \textbf{Radiação}: transferência através de radiação eletromagnética, podendo ocorrer no vazio.
\end{itemize}

\begin{figure}[H]
\centering
\begin{tikzpicture}[>=Stealth]
% condução
\draw[fill=gray!20,draw=eudark] (0,0) rectangle (3.2,.55);
\foreach \x in {.35,.75,1.15,1.55,1.95,2.35,2.75}{\fill[eublue] (\x,.275) circle (2pt);}
\draw[very thick,eured] (0,.75)--(0,-.2);
\foreach \x in {.55,1.1,1.65,2.2,2.75}{\draw[->,euorange] (\x,.9)--++(.38,0);}
\node[font=\small\bfseries] at (1.6,-.45) {condução};
% convecção
\draw[fill=eulightblue,draw=eublue] (4.2,0) rectangle (7.4,2.1);
\draw[very thick,eured] (4.7,-.15)--(6.9,-.15);
\draw[->,very thick,euorange] (5, .35).. controls (4.8,1.3) and (5.1,1.7)..(5.8,1.8);
\draw[->,very thick,eucyan] (6.8,1.7).. controls (7.1,1.0) and (6.9,.55)..(6.4,.35);
\node[font=\small\bfseries] at (5.8,-.45) {convecção};
% radiação
\fill[euorange] (8.6,1.1) circle (.25);
\draw[decorate,decoration={snake,amplitude=1.2mm,segment length=4mm},very thick,eured,->] (8.9,1.1)--(11.7,1.1);
\draw[fill=eublue!15,draw=eublue] (11.9,.35) rectangle (12.25,1.85);
\node[font=\small\bfseries] at (10.5,-.45) {radiação};
\end{tikzpicture}
\caption{Representação esquemática dos três mecanismos de transferência de energia.}
\end{figure}

\begin{nota}
A convecção natural é frequentemente explicada pela variação da massa volúmica do fluido com a temperatura: regiões aquecidas podem tornar-se menos densas e subir, enquanto regiões mais frias descem. Não é correto dizer que ``o calor sobe''; o que pode subir é o \textbf{fluido aquecido}.
\end{nota}

\section{Radiação, irradiância e geometria}
A irradiância é
\[
G=\frac{P}{A}.
\]
Para uma fonte que distribua a potência aproximadamente de forma isotrópica e sem absorção significativa no meio, a mesma potência atravessa superfícies esféricas de área $4\pi r^2$. Nesse modelo,
\[
G\propto\frac{1}{r^2}.
\]
Esta relação é uma extensão geométrica útil; fontes reais e painéis próximos podem não obedecer a este modelo simples.

\begin{figure}[H]
\centering
\begin{tikzpicture}
\begin{axis}[width=10.8cm,height=5.8cm,xlabel={distância relativa $r$},ylabel={irradiância relativa $G/G_0$},xmin=1,xmax=5,ymin=0,ymax=1.05,grid=major]
\addplot[domain=1:5,samples=120,very thick,euorange]{1/x^2};
\addplot[only marks,mark=*,eublue] coordinates{(1,1)(2,.25)(3,.111)(4,.0625)};
\end{axis}
\end{tikzpicture}
\caption{Modelo geométrico de irradiância de uma fonte pontual isotrópica: $G\propto1/r^2$.}
\end{figure}

\begin{mundo}
\textbf{Câmaras térmicas.} Um corpo à temperatura ambiente emite radiação predominantemente no infravermelho. Uma câmara térmica deteta diferenças nessa radiação e constrói uma imagem relacionada com a temperatura aparente das superfícies. A interpretação rigorosa depende também da emissividade e de reflexões do ambiente.
\end{mundo}

\begin{exercicios}
\begin{enumerate}
  \item Calcula a energia necessária para aquecer \SI{2.0}{kg} de alumínio em \SI{30}{K}, usando $c=\SI{900}{J/(kg.K)}$.
  \item Um sistema recebe \SI{500}{J} por calor e é realizado sobre ele um trabalho de \SI{120}{J}. Pela convenção deste livro, calcula $\Delta U$.
  \item Que massa de gelo pode fundir com \SI{66.8}{kJ}, usando $L_f=\SI{334}{kJ/kg}$?
  \item Uma máquina recebe \SI{800}{J} e fornece \SI{220}{J} de energia útil. Determina o rendimento.
  \item Um painel de área \SI{0.80}{m^2} recebe irradiância de \SI{750}{W/m^2} e fornece \SI{108}{W}. Calcula o rendimento.
  \item Identifica o mecanismo dominante de transferência de energia: (a) através de uma colher metálica mergulhada em sopa quente; (b) do Sol para a Terra; (c) numa corrente ascendente de ar quente.
  \item \textbf{Integração.} Um aquecedor elétrico de \SI{1.2}{kW} aquece \SI{8.0}{kg} de água em \SI{20}{K} em \SI{12}{min}. Estima a energia elétrica fornecida, a energia ganha pela água e o rendimento do aquecimento. Usa $c=\SI{4180}{J/(kg.K)}$.
  \item Um bloco de \SI{0.50}{kg} recebe \SI{9.0}{kJ} e a sua temperatura aumenta \SI{20}{K}. Estima $c$.
  \item Explica por que razão uma câmara térmica consegue ``ver'' um corpo num quarto escuro.
  \item Um painel de \SI{0.50}{m^2} recebe irradiância de \SI{800}{W/m^2} e produz \SI{72}{W}. Determina o rendimento.
  \item Converte \SI{15.0}{\celsius} para kelvin. Uma amplitude térmica de \SI{10.0}{\celsius} corresponde a quantos kelvin?
  \item Explica a Lei Zero da Termodinâmica e a sua relação com o funcionamento de um termómetro.
  \item Determina a energia necessária para aquecer \SI{0.250}{kg} de zinco de \SI{15.0}{\celsius} para \SI{70.0}{\celsius}, usando $c=\SI{388}{J/(kg.K)}$.
  \item Um corpo cúbico de aresta \SI{5.0}{cm} emite \SI{2.682e4}{J} em \SI{60}{s}. Admitindo emissão uniformemente distribuída pelas seis faces, determina a potência emitida e a potência por unidade de área da superfície.
  \item Uma lata com \SI{0.340}{kg} de bebida, $c=\SI{4.2e3}{J/(kg.K)}$, expõe \SI{140}{cm^2} a uma irradiância de \SI{600}{W/m^2} durante \SI{90}{min}. A temperatura sobe \SI{16.5}{K}. Estima a percentagem da energia incidente que aumentou a energia interna da bebida.
  \item Um painel de área \SI{50}{dm^2} tem rendimento de 12\% e recebe irradiância média de \SI{580}{W/m^2} durante \SI{8.0}{h}. Calcula a energia elétrica produzida nesse intervalo, em Wh e J.
  \item Explica por que a convecção natural num fluido aquecido por baixo pode estabelecer uma circulação macroscópica de matéria.
\end{enumerate}
\end{exercicios}

% =========================================================
\part{Desafios integradores STEM}

\chapter{Oficina de gráficos: Física e Química}
\begin{objetivos}
\begin{itemize}
  \item extrair informação quantitativa de eixos, declives, interceções, áreas e patamares;
  \item relacionar a forma de um gráfico com um modelo físico ou químico;
  \item justificar conclusões a partir de dados representados graficamente;
  \item reconhecer quando uma linearização facilita a análise experimental.
\end{itemize}
\end{objetivos}

\begin{metodo}[title=Antes de fazer contas]
Em qualquer exercício com gráficos: (1) lê as grandezas e unidades dos eixos; (2) identifica se a relação é linear, quadrática, inversa ou por patamares; (3) procura o significado físico/químico do declive e das interceções; (4) só depois substitui valores nas expressões.
\end{metodo}

\section{Química -- gráficos que contam uma história microscópica}

\subsection*{QG1 -- abundância isotópica}
\begin{figure}[H]
\centering
\begin{tikzpicture}
\begin{axis}[ybar,width=10cm,height=5.5cm,bar width=18pt,xlabel={número de massa},ylabel={abundância / \%},symbolic x coords={35,37},xtick=data,ymin=0,ymax=85,grid=major]
\addplot[eublue,fill=eublue!45] coordinates{(35,75)(37,25)};
\end{axis}
\end{tikzpicture}
\caption{Abundância isotópica simplificada de uma amostra de cloro.}
\end{figure}
\begin{exercicios}[title=Exercício QG1]
Usando o gráfico, estima a massa atómica relativa da amostra admitindo massas isotópicas aproximadas de 35 u e 37 u. Explica por que o resultado não é um número inteiro.
\end{exercicios}

\subsection*{QG2 -- energia de um fotão}
\begin{figure}[H]
\centering
\begin{tikzpicture}
\begin{axis}[width=10.5cm,height=5.5cm,xlabel={$f$ / $10^{14}$ Hz},ylabel={$E$ / $10^{-19}$ J},xmin=0,xmax=9,ymin=0,ymax=6.3,grid=major]
\addplot[only marks,mark=*,eublue] coordinates{(2,1.33)(4,2.65)(6,3.98)(8,5.30)};
\addplot[domain=0:9,very thick,eublue]{0.6625*x};
\end{axis}
\end{tikzpicture}
\caption{Energia de fotões em função da frequência (dados arredondados).}
\end{figure}
\begin{exercicios}[title=Exercício QG2]
Determina o declive da reta nas unidades indicadas e converte-o para J s. Compara o valor obtido com a constante de Planck, $h\approx\SI{6.63e-34}{J.s}$.
\end{exercicios}

\subsection*{QG3 -- energia potencial e ligação química}
\begin{figure}[H]
\centering
\begin{tikzpicture}
\begin{axis}[width=10.5cm,height=5.7cm,xlabel={distância internuclear / pm},ylabel={energia potencial / kJ mol$^{-1}$},xmin=50,xmax=350,ymin=-450,ymax=180,grid=major]
\addplot[domain=65:350,samples=180,very thick,eublue]{350*((100/x)^12-2*(100/x)^6)};
\draw[dashed,euorange] (axis cs:100,-350) -- (axis cs:100,0);
\end{axis}
\end{tikzpicture}
\caption{Curva qualitativa de energia potencial para uma ligação diatómica.}
\end{figure}
\begin{exercicios}[title=Exercício QG3]
Identifica no gráfico a distância de equilíbrio aproximada. Explica o significado físico do mínimo de energia e descreve o que acontece quando os átomos ficam muito mais próximos do que essa distância.
\end{exercicios}

\subsection*{QG4 -- volume molar de um gás}
\begin{figure}[H]
\centering
\begin{tikzpicture}
\begin{axis}[width=10cm,height=5.5cm,xlabel={$n$ / mol},ylabel={$V$ / L},xmin=0,xmax=2.1,ymin=0,ymax=52,grid=major]
\addplot[only marks,mark=*,eublue] coordinates{(.5,12)(1,24)(1.5,36)(2,48)};
\addplot[domain=0:2.1,very thick,eugreen]{24*x};
\end{axis}
\end{tikzpicture}
\caption{Volume de uma amostra gasosa em função da quantidade de matéria, a $T$ e $p$ constantes.}
\end{figure}
\begin{exercicios}[title=Exercício QG4]
Determina o declive e interpreta-o. Que volume preverias para \SI{0.75}{mol} nas mesmas condições?
\end{exercicios}

\subsection*{QG5 -- perfil energético de uma reação}
\begin{figure}[H]
\centering
\begin{tikzpicture}
\begin{axis}[width=10.8cm,height=5.8cm,xlabel={coordenada da reação},ylabel={energia / kJ mol$^{-1}$},xmin=0,xmax=10,ymin=0,ymax=130,xtick=\empty,grid=major]
\addplot[very thick,eured,smooth] coordinates{(0,80)(2,82)(4,118)(5,120)(6,100)(8,45)(10,45)};
\draw[dashed,eublue] (axis cs:0,80)--(axis cs:10,80);
\draw[dashed,eugreen] (axis cs:0,45)--(axis cs:10,45);
\node[font=\small] at (axis cs:1.2,88) {reagentes};
\node[font=\small] at (axis cs:8.8,53) {produtos};
\end{axis}
\end{tikzpicture}
\caption{Perfil energético simplificado de uma transformação química.}
\end{figure}
\begin{exercicios}[title=Exercício QG5]
Estima $\Delta H$ e a energia de ativação no sentido direto. Classifica a transformação como endotérmica ou exotérmica. Num mesmo diagrama, que grandeza seria reduzida pela presença de um catalisador?
\end{exercicios}

\section{Física -- gráficos para determinar grandezas}

\subsection*{FG1 -- massa a partir de $E_c(v^2)$}
\begin{figure}[H]
\centering
\begin{tikzpicture}
\begin{axis}[width=10cm,height=5.5cm,xlabel={$v^2$ / \si{m^2.s^{-2}}},ylabel={$E_c$ / J},xmin=0,xmax=40,ymin=0,ymax=85,grid=major]
\addplot[only marks,mark=*,eublue] coordinates{(4,8)(9,18)(16,32)(25,50)(36,72)};
\addplot[domain=0:40,very thick,eublue]{2*x};
\end{axis}
\end{tikzpicture}
\caption{Dados experimentais de energia cinética.}
\end{figure}
\begin{exercicios}[title=Exercício FG1]
Determina a massa do corpo a partir do declive. Prevê $E_c$ para $v=\SI{5.5}{m/s}$.
\end{exercicios}

\subsection*{FG2 -- trabalho em função do deslocamento}
\begin{figure}[H]
\centering
\begin{tikzpicture}
\begin{axis}[width=10cm,height=5.5cm,xlabel={$d$ / m},ylabel={$W$ / J},xmin=0,xmax=10,ymin=0,ymax=130,grid=major]
\addplot[domain=0:10,very thick,euorange]{12*x};
\end{axis}
\end{tikzpicture}
\caption{Trabalho realizado por uma força de componente constante ao longo do movimento.}
\end{figure}
\begin{exercicios}[title=Exercício FG2]
Determina a componente da força na direção do deslocamento. Se a força tiver intensidade \SI{24}{N}, calcula o ângulo que faz com o deslocamento.
\end{exercicios}

\subsection*{FG3 -- curva característica de um gerador}
\begin{figure}[H]
\centering
\begin{tikzpicture}
\begin{axis}[width=10.5cm,height=5.5cm,xlabel={$I$ / A},ylabel={$U$ / V},xmin=0,xmax=4.5,ymin=0,ymax=10,grid=major]
\addplot[only marks,mark=*,eublue] coordinates{(0,9)(1,8.4)(2,7.8)(3,7.2)(4,6.6)};
\addplot[domain=0:4.5,very thick,eublue]{9-.6*x};
\end{axis}
\end{tikzpicture}
\caption{Curva $U(I)$ de uma pilha real.}
\end{figure}
\begin{exercicios}[title=Exercício FG3]
Determina a força eletromotriz e a resistência interna da pilha. Calcula a tensão nos terminais quando $I=\SI{2.5}{A}$ e o rendimento elétrico $U/\varepsilon$ nessa situação.
\end{exercicios}

\subsection*{FG4 -- resistividade pelo declive}
\begin{figure}[H]
\centering
\begin{tikzpicture}
\begin{axis}[width=10cm,height=5.5cm,xlabel={$L$ / m},ylabel={$R$ / $\Omega$},xmin=0,xmax=5.5,ymin=0,ymax=4.8,grid=major]
\addplot[only marks,mark=*,eublue] coordinates{(1,.8)(2,1.6)(3,2.4)(4,3.2)(5,4)};
\addplot[domain=0:5.5,very thick,eugreen]{.8*x};
\end{axis}
\end{tikzpicture}
\caption{Resistência de fios do mesmo material e secção, com comprimentos diferentes.}
\end{figure}
\begin{exercicios}[title=Exercício FG4]
A área da secção dos fios é $A=\SI{2.0e-7}{m^2}$. Determina a resistividade do material a partir do declive e compara a ordem de grandeza com a de um metal ($10^{-8}$ a $10^{-6}\,\Omega\,\mathrm m$).
\end{exercicios}

\subsection*{FG5 -- aquecimento e capacidade térmica}
\begin{figure}[H]
\centering
\begin{tikzpicture}
\begin{axis}[width=10.5cm,height=5.5cm,xlabel={$Q$ / kJ},ylabel={$T$ / $^{\circ}$C},xmin=0,xmax=20,ymin=15,ymax=65,grid=major]
\addplot[only marks,mark=*,eublue] coordinates{(0,20)(4,28)(8,36)(12,44)(16,52)(20,60)};
\addplot[domain=0:20,very thick,eublue]{20+2*x};
\end{axis}
\end{tikzpicture}
\caption{Aquecimento de uma amostra sem mudança de estado.}
\end{figure}
\begin{exercicios}[title=Exercício FG5]
Determina a capacidade térmica $C=Q/\Delta T$ da amostra. Se a massa for \SI{0.25}{kg}, calcula a capacidade térmica mássica.
\end{exercicios}

\subsection*{FG6 -- irradiância e distância}
\begin{figure}[H]
\centering
\begin{tikzpicture}
\begin{axis}[width=10.5cm,height=5.5cm,xlabel={$r$ / m},ylabel={$G$ / W m$^{-2}$},xmin=.8,xmax=4.2,ymin=0,ymax=1050,grid=major]
\addplot[only marks,mark=*,eublue] coordinates{(1,1000)(2,250)(3,111)(4,62.5)};
\addplot[domain=1:4.2,samples=100,very thick,euorange]{1000/x^2};
\end{axis}
\end{tikzpicture}
\caption{Irradiância medida a diferentes distâncias de uma fonte aproximadamente pontual.}
\end{figure}
\begin{exercicios}[title=Exercício FG6]
Verifica, usando dois pares de pontos, se os dados são compatíveis com $G\propto1/r^2$. Estima a irradiância a \SI{2.5}{m}.
\end{exercicios}

% =========================================================

\chapter{Problemas que misturam Física e Química}

\section{Como resolver problemas integradores}
Em situações reais, as fronteiras entre disciplinas desaparecem. Uma bateria envolve reações químicas, transferência de carga, potência elétrica e energia térmica. Uma atmosfera planetária envolve composição química, radiação, temperatura e energia. Uma célula fotovoltaica liga estrutura eletrónica, luz e circuitos elétricos.

Use a estratégia:
\begin{enumerate}
  \item identificar o sistema e os dados;
  \item separar grandezas conhecidas e desconhecidas;
  \item desenhar um esquema;
  \item escolher modelos físicos e químicos adequados;
  \item calcular com unidades coerentes;
  \item testar a plausibilidade do resultado;
  \item explicar o significado físico/químico da resposta.
\end{enumerate}

\begin{exemplo}[title=Exemplo integrador resolvido -- aquecer uma solução com energia solar]
Um painel fornece \SI{60}{W} elétricos durante \SI{10}{min}. Essa energia alimenta um aquecedor com rendimento de 70\% que aquece \SI{0.50}{kg} de uma solução aquosa, com $c\approx\SI{4.0e3}{J/(kg.K)}$.

A energia elétrica é
\[
E=Pt=60\times600=\SI{36}{kJ}.
\]
A energia útil para a solução é
\[
E_u=0{,}70\times36=\SI{25.2}{kJ}.
\]
Logo,
\[
\Delta T=\frac{E_u}{mc}
=\frac{25\,200}{0{,}50\times4000}
\approx\SI{12.6}{K}.
\]
O problema combina potência elétrica, rendimento e energia térmica.
\end{exemplo}

\begin{exemplo}[title={Exemplo integrador resolvido -- combustível, moléculas e energia}]
Uma amostra de \SI{0.50}{mol} de um combustível liberta \SI{400}{kJ/mol}. A energia total libertada é
\[
E=0{,}50\times400=\SI{200}{kJ}.
\]
Se apenas 30\% for convertida em trabalho mecânico útil,
\[
W_u=0{,}30\times200=\SI{60}{kJ}.
\]
O número de moléculas de combustível consumidas é
\[
N=0{,}50N_A\approx3{,}01\times10^{23}.
\]
Assim, uma descrição macroscópica de energia pode ser ligada ao número microscópico de entidades.
\end{exemplo}

\begin{exercicios}[title=Desafios integradores]
\begin{enumerate}
  \item \textbf{Bateria e energia química.} Uma bateria fornece tensão média de \SI{3.7}{V} e corrente de \SI{0.80}{A} durante \SI{2.5}{h}. (a) Calcula a energia elétrica fornecida. (b) Se apenas 85\% da energia química disponível se transforma em energia elétrica útil, estima a energia química consumida. (c) Indica duas formas de dissipação energética possíveis.

  \item \textbf{Combustível e movimento.} Um veículo utiliza uma quantidade de combustível que liberta \SI{12}{MJ} de energia química num percurso curto. Apenas 25\% chega como trabalho mecânico útil às rodas. (a) Determina a energia mecânica útil. (b) Se essa energia fosse usada apenas para aumentar a energia cinética de um veículo de \SI{1000}{kg} inicialmente em repouso, qual seria a rapidez teórica? (c) Explica por que esse resultado não representa uma situação real de condução.

  \item \textbf{Espectro e temperatura.} Um sensor ótico usa radiação de \SI{850}{nm}. (a) Calcula a energia de um fotão. (b) Compara qualitativamente com um fotão de \SI{425}{nm}. (c) Explica por que duplicar a frequência duplica a energia de cada fotão.

  \item \textbf{Solução aquecida.} Preparam-se \SI{250}{mL} de uma solução aquosa contendo \SI{0.100}{mol} de soluto. A solução, de massa aproximada \SI{260}{g}, é aquecida de \SI{20}{\celsius} a \SI{35}{\celsius}. Assume $c\approx\SI{4.0}{J/(g.K)}$. (a) Calcula a concentração molar. (b) Estima a energia recebida pela solução. (c) Se o aquecedor tem potência de \SI{100}{W} e rendimento de 65\%, estima o tempo mínimo.

  \item \textbf{Painel fotovoltaico e eletrólise — análise energética.} Um pequeno painel fornece \SI{20}{W} durante \SI{30}{min}. Um processo eletroquímico experimental aproveita 60\% dessa energia. (a) Calcula a energia elétrica produzida. (b) Determina a energia efetivamente disponível para o processo químico. (c) Discute por que energia e quantidade de matéria estão relacionadas em processos eletroquímicos, mesmo sem usar leis eletroquímicas avançadas.

  \item \textbf{Atmosfera e moléculas.} Numa amostra de ar de \SI{24.0}{L}, em condições com $V_m=\SI{24.0}{L/mol}$, considera 0,040\% em quantidade de \ce{CO2}. (a) Estima a quantidade de matéria total de gás. (b) Estima os mol de \ce{CO2}. (c) Determina o número de moléculas de \ce{CO2}. (d) Explica por que uma pequena fração pode ainda ser relevante para processos radiativos atmosféricos.

  \item \textbf{Travagem regenerativa.} Um automóvel elétrico de \SI{1500}{kg} reduz a rapidez de \SI{20}{m/s} para \SI{10}{m/s}. O sistema recupera 55\% da diminuição de energia cinética para a bateria. (a) Calcula a diminuição de $E_c$. (b) Calcula a energia recuperada. (c) Se a bateria recebe essa energia a uma tensão média de \SI{350}{V}, estima a carga elétrica armazenada usando $E=UQ$.

  \item \textbf{Design STEM — miniestufa.} Pretende-se construir uma miniestufa para manter plantas alguns graus acima da temperatura exterior durante a noite. Propõe um projeto que inclua: materiais, mecanismos de transferência de energia, forma de medir temperatura, hipótese testável, variáveis a controlar, tabela de recolha de dados e critério para comparar dois protótipos.
  \item \textbf{Sensor e bateria.} Um sensor consome \SI{0.30}{W} continuamente e é alimentado por uma bateria que pode fornecer \SI{18}{kJ} de energia elétrica útil. Estima a autonomia ideal. Se a eficiência global do sistema for 80\%, que energia química mínima deve estar disponível?
  \item \textbf{Fotão e ligação.} Uma ligação química requer aproximadamente \SI{250}{kJ/mol} para ser quebrada. Determina a energia por ligação e estima o comprimento de onda de um fotão com essa energia. Discute por que, numa molécula real, absorver um fotão não garante automaticamente a quebra dessa ligação.
\end{enumerate}
\end{exercicios}

\begin{exemplo}[title=Exemplo resolvido -- painel solar a aquecer água]
Um painel fotovoltaico fornece em média \SI{120}{W} durante \SI{25}{min}. A energia elétrica alimenta uma resistência que aquece \SI{1.5}{kg} de água, inicialmente a \SI{18}{\celsius}. Se 80\% da energia elétrica chegar efetivamente à água, estima a temperatura final. Usa $c=\SI{4180}{J/(kg.K)}$.

A energia elétrica produzida é
\[
E=P\Delta t=120\times(25\times60)=\SI{1.80e5}{J}.
\]
A energia útil transferida para a água é
\[
Q=0{,}80\times1{,}80\times10^5=\SI{1.44e5}{J}.
\]
Logo,
\[
\Delta T=\frac{Q}{mc}=\frac{1{,}44\times10^5}{1{,}5\times4180}\approx\SI{23}{K}.
\]
A temperatura final é aproximadamente $18+23=\SI{41}{\celsius}$. O problema liga potência elétrica, rendimento e aquecimento.
\end{exemplo}

\begin{exemplo}[title=Exemplo resolvido -- energia luminosa e quantidade de matéria]
Um processo fotoquímico necessita de \SI{3.0}{kJ} de radiação de comprimento de onda \SI{365}{nm}. Quantos fotões correspondem a essa energia e qual a quantidade de matéria de fotões?
\[
E_{fot}=\frac{hc}{\lambda}\approx5{,}44\times10^{-19}\,\mathrm{J}.
\]
Assim,
\[
N=\frac{3{,}0\times10^3}{5{,}44\times10^{-19}}\approx5{,}5\times10^{21}\ \text{fotões}
\]
e
\[
n=\frac{N}{N_A}\approx9{,}1\times10^{-3}\,\mathrm{mol}.
\]
Este exercício mostra como uma ideia de Física -- energia da radiação -- se liga diretamente à linguagem quantitativa da Química.
\end{exemplo}

\chapter{Ciência, cidadania e decisões informadas}

As Aprendizagens Essenciais de 2026 reforçam a articulação da Física e da Química com problemas sociais, ambientais e tecnológicos. Esta secção não acrescenta uma nova lista de fórmulas; mostra como usar o conhecimento científico para analisar decisões do mundo real de forma fundamentada.

\section{Recursos naturais, materiais e sustentabilidade}
Metais, semicondutores, água, combustíveis e matérias-primas para baterias ou painéis solares são recursos limitados. Uma tecnologia não deve ser analisada apenas durante a utilização: importa considerar extração de matérias-primas, produção, transporte, durabilidade, reparação, reutilização e reciclagem.

\begin{stem}
\textbf{Análise de ciclo de vida simplificada.} Escolhe dois materiais candidatos para um objeto escolar. Compara massa, durabilidade, possibilidade de reciclagem, energia de produção e distância de transporte. Define critérios antes de procurar os dados e apresenta a decisão como uma conclusão condicionada aos critérios adotados.
\end{stem}

\section{Risco laboratorial e comunicação de segurança}
O risco depende do perigo intrínseco e da exposição. No laboratório, reduzir quantidades, usar proteção adequada, conhecer pictogramas, impedir misturas não planeadas e encaminhar resíduos corretamente são medidas de prevenção. Uma boa comunicação de segurança deve indicar \textbf{o perigo, a prevenção e o procedimento de resposta}, sem recorrer a linguagem alarmista.

\begin{metodo}[title=Mini-protocolo de emergência]
Antes de uma atividade prática identifica: (1) substâncias e equipamentos envolvidos; (2) perigos previsíveis; (3) medidas de prevenção; (4) procedimento perante derrame, contacto ou quebra de material; e (5) destino dos resíduos. As instruções específicas do laboratório e do professor prevalecem sempre sobre este esquema geral.
\end{metodo}

\section{Radiação: aplicações, benefícios e risco}
A palavra ``radiação'' inclui fenómenos muito diferentes. Ondas de rádio, micro-ondas, infravermelho e luz visível são radiações não ionizantes nas condições usuais. Raios X e radiação gama possuem energia por fotão suficiente para ionizar matéria e exigem controlo de exposição. A radiação ultravioleta ocupa uma região intermédia em que os efeitos fotobiológicos são relevantes e aumentam, em geral, com a energia dos fotões.

Na medicina, radiação pode ser usada para diagnóstico e tratamento; na indústria, para inspeção de materiais; em sensores, comunicações e energia, outras regiões do espetro são indispensáveis. A avaliação de risco deve considerar tipo de radiação, intensidade, tempo de exposição e proteção disponível.

\section{Energia, mobilidade e segurança rodoviária}
A energia cinética cresce com o quadrado da rapidez. Esta relação ajuda a compreender por que pequenas alterações de velocidade podem ter grande efeito na energia que tem de ser dissipada numa travagem ou colisão.

\begin{exemplo}[title=Exemplo resolvido -- rapidez e distância de travagem]
Considere o mesmo automóvel a \SI{50}{km/h} e a \SI{100}{km/h}. A segunda rapidez é o dobro da primeira, logo
\[
\frac{E_{c,100}}{E_{c,50}}=\left(\frac{100}{50}\right)^2=4.
\]
Se a força média de travagem for aproximadamente a mesma e se ignorarmos o tempo de reação, o trabalho dissipativo $F d$ necessário para parar também quadruplica. Assim, neste modelo simples, a distância de travagem a \SI{100}{km/h} é aproximadamente quatro vezes a de \SI{50}{km/h}. O modelo real deve ainda considerar aderência, pneus, inclinação, sistema de travagem e condições meteorológicas.
\end{exemplo}

\section{Literacia energética e acesso à tecnologia}
Comparar soluções energéticas exige distinguir \textbf{potência} de \textbf{energia}, calcular consumos, considerar rendimento e reconhecer que custo, acesso à tecnologia e qualidade das habitações condicionam as possibilidades de redução de consumo. Uma proposta tecnicamente eficiente pode não ser social ou economicamente acessível a todas as famílias.

\begin{mundo}
Uma auditoria energética escolar pode combinar Física, Matemática e cidadania: mede-se ou estima-se consumo, identificam-se perdas, calculam-se poupanças e discutem-se medidas de baixo custo e intervenções estruturais. A recomendação deve indicar dados, pressupostos, incertezas e quem seria afetado pela solução.
\end{mundo}

\section{Avaliar alegações científicas}
Nem todas as afirmações apresentadas como ``científicas'' são sustentadas por evidência. Antes de aceitar uma alegação, pergunta: a fonte é identificável? os dados podem ser verificados? a amostra é adequada? existe grupo de comparação? foram consideradas explicações alternativas? há concordância com resultados independentes? A autoridade de quem fala não substitui a qualidade das evidências.

\begin{metodo}[title=Checklist contra desinformação científica]
Se uma publicação usa um gráfico sem eixos, percentagens sem valores absolutos, uma correlação como se fosse causalidade ou um único caso como prova geral, há razões para investigar melhor. Procura a fonte primária, verifica a data, compara com instituições científicas credíveis e identifica o que permanece incerto.
\end{metodo}

\chapter{Guia de projetos STEM}

\begin{figure}[H]
\centering
\begin{tikzpicture}[>=Stealth,node distance=1.4cm and 1.2cm]
\tikzset{d/.style={rounded corners=3mm,draw=eugreen,fill=eulightgreen,minimum width=2.7cm,minimum height=.95cm,align=center,font=\small\bfseries}}
\node[d] (p) {Definir\\o problema};
\node[d,right=of p] (i) {Investigar\\e medir};
\node[d,right=of i] (c) {Criar\\soluções};
\node[d,below=of c] (t) {Testar\\protótipos};
\node[d,left=of t] (a) {Analisar\\dados};
\node[d,left=of a] (m) {Melhorar e\\comunicar};
\draw[->,thick,eugreen] (p)--(i); \draw[->,thick,eugreen] (i)--(c); \draw[->,thick,eugreen] (c)--(t);
\draw[->,thick,eugreen] (t)--(a); \draw[->,thick,eugreen] (a)--(m); \draw[->,thick,eugreen] (m)--(p);
\end{tikzpicture}
\caption{Ciclo de projeto de engenharia usado nos desafios STEM.}
\end{figure}

Um projeto STEM não termina quando o protótipo ``funciona''. Deve existir uma pergunta clara, critérios mensuráveis, dados que permitam avaliar a solução e uma fase de melhoria. A documentação do processo é tão importante como o produto final: regista decisões, falhas, alterações e limitações.

\section{Projeto 1 — Auditoria energética de uma sala}
\textbf{Questão orientadora:} onde se pode reduzir consumo energético sem perder conforto?

\textbf{Etapas:}
\begin{enumerate}
  \item inventariar equipamentos e respetivas potências;
  \item estimar tempos de utilização;
  \item calcular energia diária em kWh;
  \item identificar transferências térmicas relevantes;
  \item propor duas medidas de melhoria;
  \item estimar poupança energética e justificar pressupostos.
\end{enumerate}

\section{Projeto 2 — Espectroscópio simples}
Construir um espectroscópio educativo de baixa resolução com rede de difração adequada ou material próprio para demonstração, observar fontes luminosas seguras e comparar padrões. Não olhar diretamente para o Sol nem para fontes laser.

\section{Projeto 3 — Propriedades de materiais}
Comparar materiais quanto a uma propriedade física mensurável: isolamento térmico, resistência elétrica, massa volúmica ou absorção de água. Formular hipótese, construir protótipo, recolher dados, calcular incerteza e comunicar limitações.

\section{Projeto 4 — Caixa térmica inteligente}
Construir dois protótipos de uma pequena caixa térmica com materiais de isolamento diferentes. Colocar no interior a mesma massa de água inicialmente à mesma temperatura, registar $T(t)$ e comparar a taxa de arrefecimento. Uma extensão tecnológica pode usar um sensor digital de temperatura e registo automático. A solução final deve equilibrar desempenho, massa, custo e possibilidade de reutilização dos materiais.

\section{Projeto 5 — Qualidade do ar e mobilidade}
Usar dados públicos ou sensores escolares para comparar, ao longo de vários dias, indicadores de qualidade do ar com intensidade de tráfego ou condições meteorológicas. A análise deve separar observação de correlação e explicação causal, indicar limitações e propor uma ação baseada nos dados.

\section{Rubrica simples de avaliação STEM}
\begin{table}[H]
\centering
\caption{Rubrica indicativa para projetos.}
\renewcommand{\arraystretch}{1.18}\begin{tabularx}{\textwidth}{lXXXX}
\toprule
Critério & 1 & 2 & 3 & 4 \\
\midrule
Problema & pouco definido & definido & claro & claro e mensurável \\
Modelo científico & frágil & parcial & adequado & adequado e bem justificado \\
Dados & insuficientes & alguns dados & dados organizados & dados + incerteza/análise \\
Solução técnica & incipiente & funcional & bem construída & iterada e otimizada \\
Comunicação & pouco clara & compreensível & clara & clara, crítica e visual \\
\bottomrule
\end{tabularx}\renewcommand{\arraystretch}{1.0}
\end{table}

% =========================================================
\clearpage
\begin{center}
\begin{tikzpicture}
\draw[rounded corners=5mm,fill=eulightblue,draw=eublue,line width=1pt] (0,0) rectangle (13.2,4.3);
\node[font=\Large\bfseries,text=eublue] at (6.6,3.35) {Antes de consultar as soluções...};
\node[align=center,text width=11.5cm,font=\small] at (6.6,2.0) {Tenta resolver cada problema com um esquema, identifica as grandezas conhecidas e desconhecidas, escreve primeiro a relação física ou química e só depois substitui os valores. Usa a solução para comparar o \textbf{processo}, não apenas o resultado numérico.};
\draw[euorange,line width=2pt] (2.0,.75)--(11.2,.75);
\end{tikzpicture}
\end{center}
\vspace{1cm}
\part{Soluções}
% =========================================================

\chapter{Soluções dos exercícios}

\section*{Capítulo 1 — Medição, incerteza, dados e trabalho científico}
\begin{enumerate}
  \item (a) \SI{20}{m/s}; (b) $3{,}5\times10^{-6}\,\mathrm{kg}$; (c) \SI{2.4}{\micro m}.
  \item $\bar m=\SI{52.425}{g}\approx\SI{52.43}{g}$. Desvios: 0,005; 0,015; 0,015; 0,005 g. Desvio absoluto médio $\approx\SI{0.010}{g}$.
  \item $\varepsilon_r=|2{,}64-2{,}70|/2{,}70\times100\approx2{,}2\%$.
  \item Um erro sistemático pode deslocar todas as medições de forma semelhante: ficam próximas entre si, mas afastadas do valor de referência.
  \item $\mathrm{J/m}=\mathrm{N}$; o declive sugere uma força resultante constante de aproximadamente \SI{3.2}{N}.
\end{enumerate}

\section*{Capítulo 2 — Átomos, isótopos, quantidade de matéria e composição}
\begin{enumerate}
  \item Átomo: 11 p, 12 n, 11 e. \ce{Na+}: 11 p, 12 n, 10 e.
  \item $N=0{,}250\NA\approx1{,}51\times10^{23}$ moléculas; átomos O $\approx3{,}01\times10^{23}$.
  \item $m=nM=0{,}150\times100{,}1\approx\SI{15.0}{g}$.
  \item $18/250\times100=7{,}2\%$.
  \item $A_r=10{,}0(0{,}199)+11{,}0(0{,}801)=10{,}801\approx10{,}80$.
  \item Massa da solução $=500\times1{,}05=\SI{525}{g}$; massa de soluto $=0{,}080\times525=\SI{42}{g}$.
  \item \SI{0.1}{nm} $<$ \SI{2}{\micro m} $<$ \SI{0.08}{mm} $<$ \SI{3}{cm}.
  \item $n=N/N_A\approx1{,}20\times10^{22}/6{,}022\times10^{23}\approx\SI{1.99e-2}{mol}$.
  \item $V_g=4{,}6/100=\SI{0.046}{mL}$. Repetir reduz a influência de variações aleatórias e permite avaliar a dispersão/precisão.
  \item $n=1{,}0\times10^{-3}/12{,}01\approx8{,}33\times10^{-5}\,\mathrm{mol}$; $N\approx5{,}02\times10^{19}$ átomos.
  \item $m_{Ag}=0{,}72\times0{,}030=\SI{0.0216}{g}$; $n\approx2{,}00\times10^{-4}\,\mathrm{mol}$; $N\approx1{,}21\times10^{20}$ átomos.
  \item $n(\ce{O2})=N_{\text{átomos}}/(2N_A)\approx\SI{15.0}{mol}$.
  \item $m=0{,}050\times0{,}997=\SI{0.0499}{g}$; $n\approx2{,}77\times10^{-3}\,\mathrm{mol}$; $N_{\text{moléculas}}\approx1{,}67\times10^{21}$; total de átomos $\approx5{,}00\times10^{21}$.
  \item $n(\ce{CO2})\approx0{,}500\,\mathrm{mol}$; massas: \SI{22.0}{g}, \SI{14.0}{g} e \SI{38.4}{g}; total $\approx\SI{74.4}{g}$.
\end{enumerate}

\section*{Capítulo 3 — Energia dos eletrões, luz e Tabela Periódica}
\begin{enumerate}
  \item $f=c/\lambda=3{,}00\times10^8/(600\times10^{-9})=5{,}0\times10^{14}\,\mathrm{Hz}$.
  \item $E=hf\approx3{,}31\times10^{-19}\,\mathrm{J}$.
  \item Porque os níveis eletrónicos são quantizados; apenas certas diferenças de energia são possíveis.
  \item O: $1s^2 2s^2 2p^4$; Na: $1s^2 2s^2 2p^6 3s^1$; Cl: $1s^2 2s^2 2p^6 3s^2 3p^5$; Ca: $[\mathrm{Ar}]4s^2$; V: $[\mathrm{Ar}]4s^2 3d^3$.
  \item Na tem maior raio que Cl no mesmo período; a carga nuclear efetiva aumenta globalmente da esquerda para a direita.
  \item Mg tem maior primeira energia de ionização do que Ca, por estar acima no grupo, com eletrão de valência menos afastado do núcleo.
  \item $E\approx3{,}37\times10^{-19}\,\mathrm{J}$ por fotão; por mole, $E_m=E\NA\approx\SI{203}{kJ/mol}$.
  \item Como $E=hc/\lambda$, a ordem crescente é \SI{900}{nm} $<$ \SI{600}{nm} $<$ \SI{300}{nm}.
  \item Os eletrões internos estão, em geral, mais próximos do núcleo e sentem maior atração efetiva, exigindo mais energia para serem removidos.
  \item $d=5{,}00/(80{,}25-78{,}40)\approx2{,}70$.
  \item $f\approx5{,}41\times10^{14}\,\mathrm{Hz}$; $T\approx1{,}85\times10^{-15}\,\mathrm{s}$; $E\approx3{,}58\times10^{-19}\,\mathrm{J}$.
  \item $E\propto1/\lambda$, logo $E_{\rm violeta}/E_{\rm vermelha}=700/400=1{,}75$.
  \item Porque absorção e emissão envolvem diferenças entre os mesmos níveis de energia permitidos; os fotões envolvidos têm as mesmas energias e, portanto, os mesmos comprimentos de onda.
  \item Frequência crescente: vermelho $<$ verde $<$ violeta; energia crescente: vermelho $<$ verde $<$ violeta.

  \item $E_2=-5{,}45\times10^{-19}\,\mathrm J$; $E_4=-1{,}36\times10^{-19}\,\mathrm J$. Para $4\to2$, $E_{\text{fotão}}=|E_2-E_4|\approx4{,}09\times10^{-19}\,\mathrm J$.
  \item Série de Lyman; região ultravioleta.
  \item Para $n=2$, $E_2=-5{,}45\times10^{-19}\,\mathrm J$; logo $E_{\text{ion}}=5{,}45\times10^{-19}\,\mathrm J$ por átomo.
  \item $1s^2 2s^2 2p^6 3s^2 3p^5$; 17 eletrões, cloro.
  \item Maior $n=3$: período 3; termina em $p$: bloco $p$; $3s^23p^3$ dá 5 eletrões de valência: grupo 15.
  \item Todas têm 10 eletrões. Quanto maior $Z$, menor o raio: $\ce{Mg^{2+}}<\ce{Na+}<\ce{F^-}<\ce{O^{2-}}$.
\end{enumerate}

\section*{Capítulo 4 — Ligação química, geometria e forças intermoleculares}
\begin{enumerate}
  \item \ce{NaCl}: iónica; \ce{O2}: covalente; Cu: metálica.
  \item Estruturas de Lewis usuais com pares não ligantes apropriados em O e N.
  \item \ce{CO2}: linear; \ce{CH4}: tetraédrica; \ce{NH3}: piramidal trigonal; \ce{H2O}: angular.
  \item Os dipolos das duas ligações C--O têm igual intensidade e sentidos opostos numa geometria linear, anulando-se.
  \item Etano saturado; eteno e etino insaturados.
  \item Etanol: álcool; ácido etanoico: ácido carboxílico; propanona: cetona.
  \item Etanol, devido a interações intermoleculares mais fortes, incluindo pontes de hidrogénio.
  \item Comprimento: C--C $>$ C=C $>$ C$\equiv$C; energia de ligação: C--C $<$ C=C $<$ C$\equiv$C.
  \item Apolar: os dipolos das três ligações anulam-se na geometria trigonal planar simétrica.
  \item \ce{NH3} e \ce{H2O}; \ce{CH4} não possui H ligado a N, O ou F.

  \item (a) átomos praticamente sem interação; (b) distância de equilíbrio e estado mais estável; (c) predominam fortemente interações repulsivas.
  \item H: 1 eletrão de valência; O: 6; Na: 1; Cl: 7. A notação de Lewis deve mostrar esse número de pontos em torno do símbolo.
  \item No sólido iónico existe uma rede tridimensional extensa de catiões e aniões; \ce{NaCl} indica a proporção 1:1, não moléculas discretas.
  \item A energia potencial aumentaria acentuadamente devido às repulsões, tornando a configuração muito menos estável; os núcleos tenderiam a afastar-se para a distância de equilíbrio.
\end{enumerate}

\section*{Capítulo 5 — Gases, atmosfera, dispersões e soluções}
\begin{enumerate}
  \item $n=V/V_m=4{,}80/24{,}0=\SI{0.200}{mol}$.
  \item $V=nV_m=0{,}300\times24{,}5=\SI{7.35}{L}$.
  \item $n=5{,}85/58{,}5=\SI{0.100}{mol}$; $c=0{,}100/0{,}500=\SI{0.200}{mol/L}$.
  \item $V_i=c_fV_f/c_i=0{,}100\times250/1{,}00=\SI{25.0}{mL}$.
  \item Estratosférico: absorve UV; troposférico: em excesso atua como poluente oxidante.
  \item $n=12{,}0/24{,}0=\SI{0.500}{mol}$; $M=m/n=22{,}0/0{,}500=\SI{44.0}{g/mol}$.
  \item $n_{\rm tot}=2{,}50$ mol; $x_{\ce{N2}}=0{,}80$ e $x_{\ce{O2}}=0{,}20$.
  \item $12/300\times100=4{,}0\%$ (m/m).
  \item $V_i=c_fV_f/c_i=0{,}0200\times500/0{,}500=\SI{20.0}{mL}$.
\end{enumerate}

\section*{Capítulo 6 — Transformações químicas e energia}
\begin{enumerate}
  \item Exotérmico.
  \item $\Delta H=+\SI{40}{kJ/mol}$.
  \item $\Delta H\approx620-760=-\SI{140}{kJ/mol}$.
  \item É necessário fornecer energia para superar a atração que mantém os átomos ligados e levar o sistema a maior energia potencial.
  \item $E=hc/\lambda\approx6{,}63\times10^{-19}\,\mathrm{J}$.
  \item Para \SI{400}{nm}, $E_f\approx4{,}97\times10^{-19}\,\mathrm{J}$; $N=2000/E_f\approx4{,}0\times10^{21}$ fotões.
  \item $\Delta H\approx1520-1380=+\SI{140}{kJ/mol}$; endotérmica.
  \item Funciona como controlo: permite atribuir a mudança observada à exposição à luz e não apenas ao tempo ou a outra variável.
  \item O ozono estratosférico absorve UV e tem função protetora; o ozono troposférico, em concentrações elevadas, é um poluente oxidante.
\end{enumerate}

\section*{Capítulo 7 — Energia e movimentos}
\begin{enumerate}
  \item $E_c=\frac12(1200)(20)^2=\SI{2.4e5}{J}$.
  \item $\Delta E_{p,g}=mgh=0{,}50\times9{,}8\times3{,}0=\SI{14.7}{J}$.
  \item $W=10\times5\times\cos60^\circ=\SI{25}{J}$.
  \item $W_R=\frac12(1{,}5)(6^2-2^2)=\SI{24}{J}$.
  \item $v=\sqrt{2gh}=\sqrt{98}\approx\SI{9.9}{m/s}$.
  \item A energia mecânica transformou-se sobretudo em energia interna do bloco/rampa e eventualmente som.
  \item $\Delta E_{p,g}=900\times9{,}8\times12\approx\SI{1.06e5}{J}$; $P\approx\SI{5.29}{kW}$.
  \item Fração recuperada $\approx0{,}75/1{,}20=0{,}625$; dissipação $\approx37{,}5\%$.
  \item O declive tem unidade N e representa a força resultante aproximadamente constante: \SI{2.8}{N}.
  \item $P=36\,000/12=\SI{3.0}{kW}$.
  \item $W=Fd\cos30^\circ\approx50\times4{,}0\times0{,}866\approx\SI{173}{J}$.
  \item O declive é $m/2$; $m=2\times40=\SI{80}{kg}$.
  \item $E_{m,i}=\SI{58.8}{J}$; $W_{atr}=\SI{-35}{J}$; $E_{c,f}=\SI{23.8}{J}$; $v\approx\SI{4.9}{m/s}$.
  \item $P=500/4{,}0=\SI{125}{W}$.
  \item $E_c\propto v^2$. Duplicar $v$ quadruplica a energia que a força de travagem deve retirar. Se $W=F d$ e $F$ se mantiver aproximadamente constante, $d$ também quadruplica.
\end{enumerate}

\section*{Capítulo 8 — Energia e fenómenos elétricos}
\begin{enumerate}
  \item $Q=I\Delta t=0{,}40\times120=\SI{48}{C}$.
  \item $I=U/R=\SI{0.50}{A}$; $P=UI=\SI{3.0}{W}$.
  \item Série: \SI{30}{\ohm}; paralelo: $R_{eq}=\SI{6.67}{\ohm}$.
  \item $E=9\times5=\SI{45}{Wh}=\SI{1.62e5}{J}$.
  \item $\varepsilon=\SI{1.55}{V}$; $r=\SI{0.80}{\ohm}$.
  \item O efeito Joule aumenta com a corrente; em condutores óhmicos $P=RI^2$.
  \item $P=18\times2{,}5=\SI{45}{W}$; $E=45\times3=\SI{135}{Wh}=\SI{4.86e5}{J}$.
  \item $U=\varepsilon-rI=9{,}0-1{,}0\times0{,}50=\SI{8.5}{V}$.
  \item Cada ramo: $I=6/12=\SI{0.50}{A}$; corrente total \SI{1.0}{A}.
  \item $R\propto L$ para o mesmo material e área; o fio com o dobro do comprimento tem o dobro da resistência.
  \item (a) eletrões de condução; (b) catiões e aniões; (c) eletrões e iões.
  \item $R\approx\SI{0.223}{\ohm}$; $U=RI\approx\SI{3.23}{V}$.
  \item $P=216/60=\SI{3.60}{W}$; $U=P/I=\SI{12.0}{V}$; $R=U/I=\SI{40.0}{\ohm}$.
  \item $I=Q/t=(5{,}0\times10^{18}\times1{,}602\times10^{-19})/60\approx\SI{1.34e-2}{A}=\SI{13.4}{mA}$.
  \item Para mesma área, o declive de $R(L)$ é $\rho/A$. Maior declive implica maior resistividade; portanto $\rho_Y>\rho_X$.
\end{enumerate}

\section*{Capítulo 9 — Energia, fenómenos térmicos e radiação}
\begin{enumerate}
  \item $Q=2{,}0\times900\times30=\SI{5.4e4}{J}$.
  \item $\Delta U=Q+W=500+120=\SI{620}{J}$.
  \item $m=Q/L_f=66{,}8/334=\SI{0.200}{kg}$.
  \item $\eta=220/800\times100=27{,}5\%$.
  \item $P_{inc}=750\times0{,}80=\SI{600}{W}$; $\eta=108/600=18\%$.
  \item (a) condução; (b) radiação; (c) convecção.
  \item $E_{el}=1200\times720=\SI{8.64e5}{J}$; $Q_{\text{água}}=8\times4180\times20=\SI{6.69e5}{J}$; $\eta\approx77{,}4\%$.
  \item $c=Q/(m\Delta T)=9000/(0{,}50\times20)=\SI{900}{J/(kg.K)}$.
  \item Porque corpos à temperatura ambiente emitem radiação infravermelha; a câmara deteta essa radiação, independentemente da luz visível.
  \item $P_{\rm inc}=800\times0{,}50=\SI{400}{W}$; $\eta=72/400=18\%$.
  \item $T=15{,}0+273{,}15=\SI{288.15}{K}$; uma amplitude de \SI{10.0}{\celsius} corresponde a \SI{10.0}{K}.
  \item Se A e B estão separadamente em equilíbrio térmico com C, então A e B estão à mesma temperatura e estariam em equilíbrio térmico entre si. Um termómetro funciona estabelecendo equilíbrio térmico com o sistema e associando uma propriedade mensurável a essa temperatura.
  \item $Q=mc\Delta T=0{,}250\times388\times55\approx\SI{5.34e3}{J}$.
  \item Área total $A=6(0{,}050)^2=\SI{0.0150}{m^2}$; $P=26820/60=\SI{447}{W}$; $P/A\approx\SI{2.98e4}{W/m^2}$.
  \item $E_{inc}=600\times0{,}0140\times5400\approx\SI{4.54e4}{J}$; $Q\approx\SI{2.36e4}{J}$; percentagem $\approx52\%$.
  \item $A=50\,\mathrm{dm^2}=\SI{0.50}{m^2}$; $P_{\text{útil}}=0{,}12\times580\times0{,}50=\SI{34.8}{W}$; $E=34{,}8\times8=\SI{278}{Wh}\approx\SI{1.00e6}{J}$.
  \item Ao aquecer, certas regiões do fluido expandem-se e tornam-se menos densas; sob ação gravítica sobem enquanto regiões mais frias e densas descem, estabelecendo correntes de convecção. Há transporte macroscópico de matéria e energia.
\end{enumerate}


\section*{Oficina de gráficos -- soluções}
\addcontentsline{toc}{section}{Oficina de gráficos -- soluções}
\begin{enumerate}[label=\textbf{QG\arabic*.},leftmargin=*]
\item $A_r\approx0{,}75(35)+0{,}25(37)=35{,}5$. Não é inteiro porque é uma média ponderada das massas isotópicas da amostra.
\item O declive no gráfico é cerca de $0{,}6625\times10^{-19}/10^{14}=6{,}625\times10^{-34}\,\mathrm{J\,s}$, em excelente acordo com $h$.
\item O mínimo ocorre perto de \SI{100}{pm} no gráfico esquemático. Nesse ponto a configuração é energeticamente mais estável; a distâncias muito menores, a energia cresce rapidamente devido às interações repulsivas.
\item Declive $V/n=\SI{24}{L/mol}$, que corresponde ao volume molar nas condições do gráfico. Para \SI{0.75}{mol}, $V\approx\SI{18}{L}$.
\item $\Delta H\approx45-80=\SI{-35}{kJ/mol}$: reação exotérmica. $E_a\approx120-80=\SI{40}{kJ/mol}$. Um catalisador reduz a energia de ativação, não altera $\Delta H$.
\end{enumerate}

\begin{enumerate}[label=\textbf{FG\arabic*.},leftmargin=*]
\item Declive $a=\SI{2.0}{kg}$; como $a=m/2$, $m=\SI{4.0}{kg}$. Para $v=\SI{5.5}{m/s}$, $E_c=\frac12mv^2\approx\SI{60.5}{J}$.
\item Declive $=\SI{12}{N}=F\cos\theta$. Com $F=\SI{24}{N}$, $\cos\theta=0{,}50$, logo $\theta=\SI{60}{\degree}$.
\item $\varepsilon=\SI{9.0}{V}$ e $r=\SI{0.60}{\ohm}$. Para \SI{2.5}{A}, $U=9{,}0-0{,}60(2{,}5)=\SI{7.5}{V}$ e $\eta=7{,}5/9{,}0\approx83\%$.
\item Declive $R/L=\SI{0.80}{\ohm/m}=\rho/A$. Logo $\rho=0{,}80\times2{,}0\times10^{-7}=\SI{1.6e-7}{\ohm.m}$, compatível com a ordem de grandeza de um metal.
\item O gráfico mostra $\Delta T=\SI{40}{K}$ quando $Q=\SI{20}{kJ}$; assim $C=Q/\Delta T=\SI{500}{J/K}$. Para \SI{0.25}{kg}, $c=\SI{2000}{J/(kg.K)}$.
\item De $r=1$ para $2$ m a irradiância passa de 1000 para \SI{250}{W/m^2}, isto é, fica a $1/4$; de 2 para 4 m volta a dividir por 4. Logo é compatível com $1/r^2$. A \SI{2.5}{m}, $G\approx1000/2{,}5^2=\SI{160}{W/m^2}$.
\end{enumerate}

\section*{Desafios integradores}
\begin{enumerate}
  \item $P=3{,}7\times0{,}80=\SI{2.96}{W}$; $E=2{,}96\times2{,}5=\SI{7.40}{Wh}=\SI{2.66e4}{J}$. Energia química $\approx2{,}66\times10^4/0{,}85=\SI{3.13e4}{J}$. Dissipações: aquecimento interno, eletrónica, perdas resistivas.
  \item Energia útil $=0{,}25\times12=\SI{3.0}{MJ}$. $v=\sqrt{2E/m}=\sqrt{6000}\approx\SI{77}{m/s}$. Na realidade há resistência do ar, rolamento, perdas no motor/transmissão e energia cinética inicial não é a única componente relevante.
  \item Para \SI{850}{nm}, $E\approx2{,}34\times10^{-19}\,\mathrm{J}$. A \SI{425}{nm} a frequência é aproximadamente dupla e a energia por fotão também.
  \item $c=0{,}100/0{,}250=\SI{0.400}{mol/L}$. $Q=260\times4{,}0\times15=\SI{1.56e4}{J}$. Potência útil $=65\,\mathrm{W}$; $t\approx15600/65\approx\SI{240}{s}\approx\SI{4.0}{min}$.
  \item $E=20\times1800=\SI{3.6e4}{J}$; útil $=0{,}60E=\SI{2.16e4}{J}$. Processos eletroquímicos ligam transferência de carga a transformações químicas e quantidade de matéria.
  \item $n_{tot}=1{,}00\,\mathrm{mol}$; fração $=0{,}00040$; $n_{CO_2}=4{,}0\times10^{-4}\,\mathrm{mol}$; $N\approx2{,}4\times10^{20}$ moléculas. Pequenas frações podem ter absorção radiativa relevante em bandas específicas.
  \item $\Delta E_c=\frac12(1500)(20^2-10^2)=\SI{2.25e5}{J}$. Recuperada $=0{,}55\times2{,}25\times10^5\approx\SI{1.24e5}{J}$. $Q=E/U\approx1{,}24\times10^5/350\approx\SI{354}{C}$.
  \item Resposta aberta. Deve apresentar um projeto mensurável, variáveis definidas, método de comparação justo, dados repetidos e critérios de avaliação.
  \item Autonomia ideal $t=E/P=18\,000/0{,}30=\SI{6.0e4}{s}\approx\SI{16.7}{h}$. Energia química mínima $=18/0{,}80=\SI{22.5}{kJ}$.
  \item Energia por ligação $E=(250\times10^3)/N_A\approx4{,}15\times10^{-19}\,\mathrm{J}$. $\lambda=hc/E\approx4{,}79\times10^{-7}\,\mathrm{m}\approx\SI{479}{nm}$. Numa molécula real, a absorção depende de estados permitidos, redistribuição de energia e competição com outros processos.
\end{enumerate}

% =========================================================
\backmatter
\chapter*{Formulário essencial}
\addcontentsline{toc}{chapter}{Formulário essencial}

\begin{longtable}{p{.38\textwidth}p{.52\textwidth}}
\toprule
Tema & Relações úteis \\
\midrule
Quantidade de matéria & $n=m/M$, $N=nN_A$ \\
Fotões & $E=hf=hc/\lambda$ \\
Ondas eletromagnéticas & $c=\lambda f$, $f=1/T$ \\
Concentração & $c=n/V$, $c_iV_i=c_fV_f$ \\
Gases & $V=nV_m$, $\rho=M/V_m$, $x_i=n_i/n_{\rm total}$ \\
Energia cinética & $E_c=\frac12mv^2$ \\
Energia potencial gravítica & $E_{p,g}=mgh$, $W_P=-\Delta E_{p,g}$ \\
Trabalho & $W=Fd\cos\theta$ \\
Teorema da energia cinética & $W_R=\Delta E_c$ \\
Energia mecânica & $E_m=E_c+E_{p,g}$, $W_{\rm nc}=\Delta E_m$ \\
Potência mecânica & $P=W/\Delta t$ \\
Corrente & $I=\Delta Q/\Delta t$ \\
Lei de Ohm & $U=RI$ \\
Resistividade & $R=\rho L/A$ \\
Potência elétrica & $P=UI=RI^2=U^2/R$ \\
Energia elétrica & $E=P\Delta t$ \\
Aquecimento & $Q=mc\Delta T$, $C=mc$ \\
Mudança de estado & $Q=mL$ \\
Primeira Lei & $\Delta U=Q+W$ (convenção adotada no livro) \\
Rendimento & $\eta=E_{\text{útil}}/E_{\text{fornecida}}$ \\
Irradiância e painel fotovoltaico & $G=P_{inc}/A=E_{inc}/(A\Delta t)$, $P_{inc}=GA$, $\eta=P_{el}/(GA)$ \\
Entalpia de reação & $\Delta H=H_p-H_r$ \\
\bottomrule
\end{longtable}

\chapter*{Constantes e dados úteis}
\addcontentsline{toc}{chapter}{Constantes e dados úteis}
\begin{itemize}
  \item $N_A=6{,}022\times10^{23}\,\mathrm{mol^{-1}}$
  \item $h=6{,}626\times10^{-34}\,\mathrm{J\,s}$
  \item $c=3{,}00\times10^8\,\mathrm{m\,s^{-1}}$
  \item $g\approx\SI{9.8}{m/s^2}$ junto à superfície terrestre
  \item $1\,\mathrm{kWh}=3{,}6\times10^6\,\mathrm{J}$
\end{itemize}

\chapter*{Referências}
\addcontentsline{toc}{chapter}{Referências}
\begin{hangparas}{1.27cm}{1}
Direção-Geral da Educação. (2026). \textit{Aprendizagens essenciais: Física e Química A — 10.º ano} (março de 2026). \url{https://www.dge.mec.pt/sites/default/files/es_10_fisico-quimica_a.pdf}

Direção-Geral da Educação. (s.d.). \textit{Aprendizagens essenciais — Ensino secundário}. Ministério da Educação. \url{https://www.dge.mec.pt/aprendizagens-essenciais-ensino-secundario}

Direção-Geral da Educação. (s.d.). \textit{Física e Química A — Currículo Nacional}. \url{https://curriculonacional.dge.mec.pt/organizacao-curricular/fisica-e-quimica-0}

International Union of Pure and Applied Chemistry. (2019). \textit{Compendium of chemical terminology (the Gold Book)}. IUPAC. \url{https://goldbook.iupac.org/}

Bureau International des Poids et Mesures. (2019). \textit{The International System of Units (SI)} (9th ed.). BIPM.

National Institute of Standards and Technology. (2019). \textit{The International System of Units (SI)} (NIST Special Publication 330). U.S. Department of Commerce.

OpenStax. (2019). \textit{Chemistry 2e}. Rice University. \url{https://openstax.org/details/books/chemistry-2e}

OpenStax. (2022). \textit{College physics 2e}. Rice University. \url{https://openstax.org/details/books/college-physics-2e}

Serway, R. A., \& Jewett, J. W. (2018). \textit{Physics for scientists and engineers with modern physics} (10th ed.). Cengage Learning.

Zumdahl, S. S., Zumdahl, S. A., \& DeCoste, D. J. (2018). \textit{Chemistry} (10th ed.). Cengage Learning.
\end{hangparas}

\chapter*{Nota editorial}
Esta edição integra e uniformiza material pedagógico proveniente das fichas de trabalho de Química fornecidas pelo autor, incluindo esquemas sobre o modelo de Bohr, transições eletrónicas, modelo quântico, espetroscopia fotoeletrónica, organização e tendências da Tabela Periódica e ligação química. As figuras externas dependentes de ficheiros não incluídos foram redesenhadas em TikZ/PGFPlots para manter o livro autónomo e compilável.


\addcontentsline{toc}{chapter}{Nota editorial}
\addcontentsline{toc}{chapter}{Nota editorial}
\textbf{EUFISICA books}\par
Autor: José Gonçalves\par
Título: \textit{Física e Química A — 10.º ano: Conceitos, gráficos, laboratório e desafios STEM}.\par

Esta edição integra e harmoniza conteúdos selecionados das fichas de trabalho de Física e de Química elaboradas pelo autor, com atualização de notação, unidades, exemplos, gráficos e enquadramento pedagógico. O manuscrito pode continuar a ser adaptado com fotografias próprias, QR codes para simulações, propostas laboratoriais adicionais, grelhas de avaliação e versões específicas para professor e aluno.

\end{document}
